\section{二十次の接触部分空間と階数の未決定部分}
\label{sec:v15-contact-subspace}

\subsection{20次の明示的接触部分空間と通常階数の上下界}
\label{subsec:v15-twentieth-contact-subspace}

本小節で決定するのは、通常接触像に含まれる明示的な328次元部分空間である。
通常接触像そのものの次元には一方向の未決定部分を残すが、
後述の実際の20次項の消去には、この部分空間への所属だけで十分である。
以下では
\[
 P=P_F=F(q(a)),\qquad g=\nabla F(q(a)),\qquad
 J_{20}=dP(\mathcal G_{15})
\]
とおく。前節の独立な整数指標計算から
\[
 \dim\mathcal C_{20}=330,\qquad
 \dim\mathcal G_{15}=1244,\qquad
 \dim\mathcal U_{14}=164,\qquad
 \dim\mathcal D_{16}=688
\]
を用いる。命題\ref{prop:v12-uniform-pfaffian-surjectivity}により
$\delta q:\mathcal G_{15}\twoheadrightarrow\mathcal D_{16}$ は全射である。
従って $dP=dF\circ\delta q$ の二つの写像 $dP,dF$ は同じ階数を持つ。
以下のテンソルの値域は、前小節と同じくtrace線を外した表示で記す。

\begin{proposition}[確定した20次接触部分空間]
\label{prop:v15-twentieth-contact-subspace}
公開された164個の普遍座標場
$\Gamma_1,\ldots,\Gamma_{164}\in\mathcal G_{15}$ を用いて
\[
 S_{20}=P\mathcal U_{14}\oplus
     \operatorname{Span}_{\CC}\{dP(\Gamma_1),\ldots,dP(\Gamma_{164})\}
\]
とおくと、
\[
 \dim S_{20}=328,\qquad S_{20}\subset J_{20}.
\]
さらに、第13・14版と同じ正規化の $U_4,J_6,\Theta_8$ に対して
\[
 \mathcal C_{20}
   =S_{20}\oplus
       \CC\,\Theta_8U_4^3\oplus\CC\,\Theta_8J_6^2.
 \tag{20S}\label{eq:v15-twentieth-certified-splitting}
\]
特に $\mathcal C_{20}/S_{20}$ は二次元であり、右辺の二つの類を基底とする。
この段階では、これらの類が $\mathcal C_{20}/J_{20}$ でも独立であるとは主張しない。
\end{proposition}

\begin{proof}
既存のテンソル規約を用い、
$F_{ijk}=F_{(ijk)}$、
$\tau_{ijk}=\epsilon_{ijklm}a_{lm}$ とする。
ここで $l<m$ は $i,j,k$ の補集合である。
左・右に各15頂点を持つ三価二部グラフについて、
最後の左頂点を除き、左の14頂点に $F$、
右の15頂点に $\tau$ を置く。
三本の自由脚を持つ縮約を $R_{ijk}$ として
\[
 \Gamma^{lm}=\frac16\sum_{i,j,k}\epsilon_{lmijk}R_{ijk}
\]
と定める。$\Gamma$ は $F^{14}a^{15}$ の次数を持ち、
determinant twist を含めて $\mathcal G_{15}$ に属する。
具体的には値域は $\bigwedge^2V\otimes(\det V)^{14}$ である。
また
\[
 T_{ijk}=\epsilon_{ijklm}\frac{\partial P}{\partial a_{lm}}
\]
を最後の左頂点に置いた閉じたグラフの縮約を6で割ると、
$dP(\Gamma)$ が得られる。

整数入力を180個取り、各入力で
\[
              q(a)=(1,1,1,1,1),\qquad P=0
\]
が有理数上で成立するようにする。
各対称テンソル係数を整数として選び、
$F_{444}$ によって $P=\sum_{i,j,k}F_{ijk}=0$ を課している。
この制限で $P\mathcal U_{14}$ は消える。
選んだ164グラフの評価行列は、1009を法として
\[
                  \det M_{\rm G}\equiv342\pmod{1009}
\]
となる $164\times164$ 小行列を持つ。
したがって $dP(\Gamma_i)$ は $P\mathcal U_{14}$ を法として特性零で独立である。
入力、グラフ、選択行、全評価値の剰余は
\texttt{contact20\_search.json} に記録した。

Eulerの関係
\[
                         dP(au/6)=Pu
\]
から $P\mathcal U_{14}\subset J_{20}$ である。
普遍係数と $a$ の多項式環は整域なので、乗法 $u\mapsto Pu$ は単射であり、
この部分空間の次元は164である。
従って $S_{20}$ は上記の直和を持ち、次元は328となる。

一方、
\[
             W_1=\Theta_8U_4^3,\qquad W_2=\Theta_8J_6^2
\]
はいずれも $F^{15}a^{20}$ の次数と $(\det V)^{17}$ の変換則を持つので
$\mathcal C_{20}$ に属する。
同じ180入力で $W_1,W_2$ を有理数により評価し、上の164列に追加すると、
\[
                  \det M_{\rm aug}\equiv667\pmod{1009}
\]
となる $166\times166$ 小行列を得る。
値と正規化は \texttt{complement20\_certificate.json} に記録した。
したがって、この166方向は $P\mathcal U_{14}$ を法として独立である。
さらに $164+166=330=\dim\mathcal C_{20}$ なので、
式\eqref{eq:v15-twentieth-certified-splitting}が従う。
\end{proof}

\begin{proposition}[通常20次階数の確定した範囲]
\label{prop:v15-twentieth-contact-rank-bounds}
通常接触写像について
\[
 328\le
 \operatorname{rank}(dP:\mathcal G_{15}\to\mathcal C_{20})
 =
 \operatorname{rank}(dF:\mathcal D_{16}\to\mathcal C_{20})
 \le329.
\]
従って
\[
 \begin{split}
 \dim\ker dP&\in\{915,916\},\\
 \dim\ker dF&\in\{359,360\},\\
 \dim(\mathcal C_{20}/J_{20})&\in\{1,2\}.
 \end{split}
\]
$\mathcal C_{20}/J_{20}$ は $[W_1],[W_2]$ によって生成される。
これらの間に通常接触像を法とする一つの関係があるかどうかは未決定である。
\end{proposition}

\begin{proof}
下界328は命題\ref{prop:v15-twentieth-contact-subspace}から従う。
上界には359個の明示的な核ベクトルを用いる。
左・右に各13頂点を持つ三価二部グラフで最後の右頂点を除き、
左に13個の $F$、右に12個の $\tau$ を置く。
三本の自由脚を持つ縮約を $R^{ijk}$ として
\[
       \beta_{lm}=\sum_{i,j,k}\epsilon_{lmijk}R^{ijk},
       \qquad h=\beta g
\]
と定める。$\beta$ は $F^{13}a^{12}$ の次数と
$\bigwedge^2V^*\otimes(\det V)^{13}$ の変換則を持つ。
従って $h$ は $F^{14}a^{16}$ の次数と
$V^*\otimes(\det V)^{15}$ の変換則を持つ $\mathcal D_{16}$ の元である。
交代性から
\[
                      dF(h)=g^{\rm t}\beta g=0
\]
が多項式恒等式として成立する。

この方法で得た359個の $h$ を90個の整数入力で評価し、
五成分を並べる。この評価では $P=0$ を課さない。
評価行列には
\[
                   \det M_{\rm K}\equiv269\pmod{1009}
\]
となる $359\times359$ 小行列がある。
全データは \texttt{koszul20\_search.json} に記録した。
従って核には特性零で独立な359方向があり、
\[
                        \operatorname{rank}dF\le688-359=329
\]
となる。$\delta q$ の全射性と各空間の次元を用いれば、残りの主張も従う。
\end{proof}

\begin{remark}[非零minorの意味と再現可能性]
\label{rem:v15-twentieth-rank-certificate-scope}
採用したグラフと整数入力を固定した後は、
上記の小行列式は全て厳密な有限計算である。
テンソル公式は整係数であり、接触場の正規化に現れる分母6は1009で可逆である。
従って非零の法1009小行列式は、有理数上の対応する小行列式の非零性を証明する。
これは階数の確率的推測ではない。

独立監査では、発見時とは異なる右頂点による辺の番号付けと
整数 \texttt{einsum} を用いて、全ての採用グラフの全入力を再縮約した。
各縮約では加算項数を数え、64 bit整数のoverflowが起こらないことを確認してから、
演算ごとに1009で還元した。
この監査は $342,269$ を再現する。
二つの補元の180点での値は有理数上で独立に再計算した。
\nolinkurl{v15_rank_certificate.py} は、これらの非零minor、
独立監査の入力hash、指標次元、および上下界をまとめて検証する。

本計算は、Koszul複体の完全性も、Euler和からの階数決定も用いない。
特に359は、構成したKoszul像の\emph{独立方向数}であり、
Koszul像全体の次元が359だと証明したものではない。
さらに多くのグラフを探索して階数が増えなかったという事実も、
数学的な上界としては用いていない。
\end{remark}

\subsection{通常商を完全決定しなくても20次消去を証明できる理由}
\label{subsec:v15-certified-subspace-suffices}

後述する三点の有理評価汎関数を $\rho_1,\rho_2$ とする。
それらが全ての $dP(\Gamma_i)$ と $P\mathcal U_{14}$ を消し、
\[
                  \rho_i(W_j)=\delta_{ij}
\]
を満たすことを有理数上で確認すれば、
式\eqref{eq:v15-twentieth-certified-splitting}から直ちに
\[
 S_{20}=\ker\rho_1\cap\ker\rho_2,\qquad
 R\in S_{20}\ \Longleftrightarrow\
       \rho_1(R)=\rho_2(R)=0\quad(R\in\mathcal C_{20})
 \tag{20rho}\label{eq:v15-twentieth-subspace-detection}
\]
が従う。
ここでは $\rho_i$ が $J_{20}$ 全体を消すことを仮定しない。
二つの汎関数は、確定した商 $\mathcal C_{20}/S_{20}$ の座標を読み取る。

特に、rawなTate係数 $d_{20}^{\mathrm T}$ に対して
$\rho_1(d_{20}^{\mathrm T})=\rho_2(d_{20}^{\mathrm T})=0$ が厳密に分かれば、
\[
                    d_{20}^{\mathrm T}\in S_{20}\subset J_{20}
\]
となる。低次の座標補正を全て含む正規化された20次項を $R_{20}$ とすると、
本版では別途、接触作用の計算から
\[
                         R_{20}-d_{20}^{\mathrm T}\in J_{20}
\]
を証明する。この二つの所属を合わせて $R_{20}\in J_{20}$ が従い、
20次項を通常接触変換で消去できる。
補正全体が $S_{20}$ に入ることや、補正全体に両汎関数が零となることは、
この結論に必要ないし、ここでは主張しない。
本版で後に示す22次から始まる正規形はこの道筋を用い、
通常20次商が一次元か二次元かという未解決問題とは論理的に独立している。

有限評価の一般的な背景も記しておく。
$R=\CC[a_{ij},F_\alpha]$ において
$P=\sum F_\alpha q(a)^\alpha$ は三次係数 $F_\alpha$ に関する一次式である。
Pfaffianの共通因子は1であり、係数に $q_i^3$ が含まれるため、
この一次式は $\CC[a]$ 上primitiveとなる。
$\CC(a)[F_\alpha]$ 上での既約性とGaussの補題から、$P$ は $R$ でも既約である。
従って $\mathcal C_{20}$ の元が $P=0$ 上で恒等的に消えれば $P$ で割り切れ、
商の次数と変換則からその商は $\mathcal U_{14}$ に属する。
よって
\[
                       \dim(\mathcal C_{20}/P\mathcal U_{14})=166.
\]
この166次元空間の基底を上記164接触像と $W_1,W_2$ が与えている。
ただし追加の式が法1009でその像に入るという観察だけから、
特性零での所属を結論してはならない。

\begin{remark}[残る一方向の問題]
通常商を完全決定するには、例えば次のいずれかが十分である。
特性零で独立な360番目の $\ker dF$ の元を構成すれば、
$\dim J_{20}=328$ と $\mathcal C_{20}/J_{20}\simeq\CC^2$ が従う。
逆に $P\mathcal U_{14}$ を法として165番目の独立接触像を構成すれば、
$\dim J_{20}=329$ と $\dim(\mathcal C_{20}/J_{20})=1$ が従う。
あるいは、上記二汎関数が $\mathcal G_{15}$ 全体の像を消すことを
普遍的に証明しても、前者の結論を得る。
これらは今後の課題であり、本版の実際の20次消去の前提ではない。
\end{remark}

% END rank/rank20_section_ja.tex


% BEGIN contact/contact_twentieth_ja.tex
% 第15版の研究用追記。2026-10-29。
% 構造命題は第13・14版の確定結果から証明する。
% 二十次の新しい有限所属判定は、このファイルでは仮定として区別する。

\section{二十次の接触作用と低次正規化の補正}
\label{sec:v15-contact-twentieth}

$P=P_F=F(q(a))$、$H=\det\operatorname{Hess}_F(q(a))$、$c=1/96$、
$J_{20}=dP(\mathcal G_{15})$ とする。
全ての場と単元は $F$ に多項式的かつ普遍同変とし、
座標変更の線形項と単元の定数項をそれぞれ恒等写像と1に固定する。
本節では、十八次における追加方向の消滅を二十次へ外挿しない。

\subsection{低次条件を保つ対数作用の二本の鎖}

第14版の接触対数補題\ref{lem:v14-contact-logarithm}を用いる。
低次の共変量の消滅により、$L=\log T$ の最小の次数増分は4であり、
\[
 L=\sum_{k\ge4}Q_k,\qquad
 Q_k=Y_{k+1}+b_k\in\mathcal G_{k+1}\oplus\mathcal U_k
\]
と書ける。$P+cH$ の19次までを保つ条件は
\[
 Q_kP+cQ_{k-4}H=0\quad(4\le k\le13),\qquad Q_j=0\ (j<4),
\]
二十次の作用は $Q_{14}P+cQ_{10}H$ である。
偶数増分の条件は
\begin{equation}\label{eq:v15-even-contact-chain}
 \begin{split}
 Q_4P&=0,\qquad Q_6P=0,\\
 Q_8P+cQ_4H&=0,\qquad Q_{10}P+cQ_6H=0,\\
 Q_{12}P+cQ_8H&=0.
 \end{split}
\end{equation}
このうち $Q_4,Q_8,Q_{12}$ の鎖は二十次の作用に現れず、全て零に選べる。
従って有効像を決めるのは $Q_6,Q_{10}$ の鎖である。
奇数増分の条件は独立であり、奇数増分を零にしてもこの像は変わらない。
よって全ての普遍接触変換と、奇座標・偶単元による接触変換の
二十次有効像は一致する。

\begin{lemma}[十一次の接触安定化場のHessian作用]
\label{lem:v15-eleventh-stabilizer-hessian}
$S_{10}=Y_{11}+b_{10}$ が $S_{10}P=0$ を満たせば
\[
 S_{10}H\equiv-\frac23b_{10}H\pmod{J_{20}}.
\]
特に $dP(Y_{11})=0$ なら $dH(Y_{11})\in J_{20}$ である。
\end{lemma}

\begin{proof}
$g=\nabla F(q)$、$h(p)=\det\operatorname{Hess}_F(p)$ と置く。
第13版の命題\ref{prop:v13-contact-rank-sixteenth}により
$\ker(dF:\mathcal D_{12}\to\mathcal C_{16})$ は二十のKoszul像で張られる。
従って $S_{10}P=0$ から
\[
 \delta q(Y_{11})=-\frac13b_{10}q+\beta g,
 \qquad \beta\in\mathcal B_{16}
\]
を得る。$h$ が $q$ に五次であることと $\beta$ の交代性から
\[
 S_{10}H=-\frac23b_{10}H-g^t\beta\nabla h(q).
\]
$\beta\nabla h(q)$ は $F^{14}a^{16}$ の次数を持つ
$\mathcal D_{16}$ の元である。
全次数のPfaffian微分の全射性
$\delta q:\mathcal G_{15}\twoheadrightarrow\mathcal D_{16}$ を用いれば、
最後の項は $J_{20}$ に入る。
\end{proof}

第13版の $H\mathcal U_6\subset J_{16}$ により、線形な持ち上げ
\[
 Z_{11}:\mathcal U_6\longrightarrow\mathcal G_{11},\qquad
 dP(Z_{11}(w))=Hw
\]
を一つ選べる。ここで $\dim\mathcal U_6=2$ である。

\begin{theorem}[二十次の有効像の有限縮約]
\label{thm:v15-effective-twentieth}
$c\ne0$ のとき、19次までの $P+cH$ を保つ普遍接触変換の二十次像は
\begin{equation}\label{eq:v15-effective-image}
 I_{20}=J_{20}+H\mathcal U_{10}
             +dH\bigl(Z_{11}(\mathcal U_6)\bigr)
\end{equation}
である。最後の項は高々二次元であり、各 $w$ に対する
$dH(Z_{11}(w))$ の $\mathcal C_{20}/J_{20}$ における類は純粋な持ち上げに依存しない。
\end{theorem}

\begin{proof}
$Q_6=Y_7+b_6$、$Q_6P=0$ を考える。
第10版の $dF:\mathcal D_8\to\mathcal C_{12}$ の単射性から
\[
 \delta q(Y_7)=-\frac13b_6q,\qquad Q_6H=-\frac23b_6H.
\]
$w=-b_6/6$ と置けば $Q_6H=4Hw$ であり、
この作用は $Q_6=w(\mathcal E-6)$ によって実現する。
ここで $\mathcal E=\sum a_i\partial_{a_i}$ はEuler導分である。
条件 $Q_{10}P+cQ_6H=0$ の一般解は
\[
 Q_{10}=-4cZ_{11}(w)+S_{10},\qquad S_{10}P=0
\]
と書ける。前補題と $P\mathcal U_{14}\subset J_{20}$ から
二十次作用は表示した右辺に含まれる。

逆に、$J_{20}$ は $Q_{14}$ で、$H\mathcal U_{10}$ は
$Q_{10}=u_{10}(\mathcal E-6)$ で実現する。
$w\in\mathcal U_6$ に対し
\[
 Q_6=w(\mathcal E-6),\qquad Q_{10}=-4cZ_{11}(w),
 \qquad Q_4=Q_8=Q_{12}=0
\]
とすれば全低次条件を満たし、二十次作用は $-4c^2dH(Z_{11}(w))$ である。
$c\ne0$ より全方向が実現する。
持ち上げの差については前補題の $b_{10}=0$ の場合を適用する。
\end{proof}

\begin{remark}[単元を残した持ち上げ]
$Hw=dP(Z'_{11}(w))+PB_{10}(w)$ であれば、
$Z_{11}(w)=Z'_{11}(w)+aB_{10}(w)/6$ と純粋な持ち上げに直せる。
そのHessian作用の差は $(5/3)HB_{10}(w)$ なので、
\eqref{eq:v15-effective-image}では $Z_{11}$ を $Z'_{11}$ に置き換えてよい。
現時点では右辺の追加部分が $J_{20}$ に含まれるとは主張していない。
\end{remark}

\begin{proposition}[十次単元の $P$ 倍部分]
\label{prop:v15-pdivisible-tenth-unit}
$\mathcal U_{10}$ の $P$ で割り切れる部分は二次元であり、
その $H$ 倍は $J_{20}$ に入る。
したがって $H\mathcal U_{10}$ の判定には、$P$ 倍部分を法とする
十六次元の商の基底を用いれば十分である。
\end{proposition}

\begin{proof}
商となる共変量の空間は
\[
 \operatorname{Hom}_{\mathrm{GL}(V)}\left(
 \operatorname{Sym}^{9}(\operatorname{Sym}^3V)
 \otimes\operatorname{Sym}^{4}(\Lambda^2V),\det(V)^7\right)
\]
である。Littlewood分解の補分割は
$(7,7,7,3,3)$、$(7,6,6,4,4)$、$(7,5,5,5,5)$ となる。
第14版の九次plethysmの全分解の証明書から、重複度は順に1、1、0である。
よって次元は二であり、整域における $P$ 倍の単射性と商の変換則により
これが $\mathcal U_{10}$ の $P$ 倍部分に一致する。
$u_{10}=Pv$ なら $Hv\in\mathcal U_{14}$ なので
$Hu_{10}=P(Hv)\in J_{20}$ である。
$\dim\mathcal U_{10}=18$ から商の次元は十六となる。
\end{proof}

\paragraph{二つの二次安定化方向の具体的な入力.}
第13版の \texttt{v13\_tensor\_rank\_certificate.json} の
\texttt{HU6\_coefficients} は、二つの $\mathcal U_6$ の基底について、
二十一の接触グラフ像と一つの補元による係数を保存している。
補元の係数はいずれも零である。先頭二十一係数で $Z'_{11}$ を作れば、
上の単元を残した持ち上げを得る。
\nolinkurl{v15_lower_contact_inputs.py} は全二十二行の有理数恒等式を再検算し、
この二つの持ち上げと十次単元の指標入力を集約する。
これは二十次の新しい所属判定とは区別する。

\subsection{二十次までの正確な補正台帳}

未正規化Tate係数を $f_n=d_n^{\mathrm T}$ と書く。
第10版のmerged規約を固定して
\[
 X_7=-\frac{15800}{6137}\Phi_7-\frac{689855}{883728}\Psi_7,
 \qquad X_8=\frac{4600}{2189}\operatorname{Alt}K_8,
\]
\[
 T_{10}(f)=(1+3U_4-V_6)f(a+X_7+X_8)
\]
とする。既存の普遍恒等式は
\[
 f_{10}=cH-3PU_4,\qquad
 f_{12}+dP(X_7)=PV_6,\qquad f_{13}+dP(X_8)=0
\]
である。単元を $(1+3U_4)(1-V_6)$ に置き換えることはしない。

第14次消去は $a-\Xi_9$ と $1-B_8$、第15次消去は10次場と9次単元を用いる。
これらの後の16次係数は
\begin{equation}\label{eq:v15-sixteenth-input}
 A_{16}=f_{16}+d(f_{10})(X_7)
             +3U_4\bigl(f_{12}+dP(X_7)\bigr)-V_6f_{10}
\end{equation}
であり、第13版により $A_{16}\in J_{16}$ である。
純粋な持ち上げ $\Xi_{11}\in\mathcal G_{11}$、
$dP(\Xi_{11})=A_{16}$ を選んで $a-\Xi_{11}$ で16次を消す。

\begin{proposition}[二十次の有限正規化補正式]
\label{prop:v15-exact-twentieth-correction}
上の有限正規化の後の二十次係数は正確に
\begin{equation}\label{eq:v15-exact-twentieth-correction}
 \begin{split}
 R_{20}^{\mathrm{fin}}={}&f_{20}+d(f_{14})(X_7)+d(f_{13})(X_8)
       +\frac12d^2P(X_8,X_8)\\
   &+3U_4\bigl(f_{16}+d(f_{10})(X_7)\bigr)-V_6f_{14}
       -c\,dH(\Xi_{11}).
 \end{split}
\end{equation}
16次を $A_{16}=dP(\Xi'_{11})+PB_{10}$ の規約で消す場合は、
最後の項を $-c(dH(\Xi'_{11})+B_{10}H)$ に置き換える。
その後の18次消去は二十次係数を変えない。
\end{proposition}

\begin{proof}
初回の座標変更で二十次に到達する項は
$f_{20}$、$d(f_{14})(X_7)$、$d(f_{13})(X_8)$、
$\tfrac12d^2P(X_8,X_8)$ である。
二つの $X_7$ は $P$ には18次、$X_7,X_8$ の混合は19次を与える。
それらが20次に達する他の元の係数は7、8次等であり零である。
三つ以上の座標補正による項は24次以上となる。
四次単元が掛かる16次係数は $f_{16}+d(f_{10})(X_7)$、
六次単元が掛かる14次係数は $f_{14}$ なので、最初の六項を得る。

14次消去が二十次へ作用するには12次係数が必要だが、初回変更により零である。
15次消去には11次係数が必要だが同様に零である。
16次消去は10次項 $cH$ へ作用するので、表示した最後の項を与える。
その二次Taylor項は26次以上である。
18次消去の13次座標場と12次単元が二十次へ作用するには8次係数が必要で、零である。
従って18次を消しても表示式は変わらない。
\end{proof}

ここには、13次項を消す偶数次座標場 $X_8$ の二階作用が現れる。
元の13次項が奇数次であることや、後で全次数偶化できることから、
この二十次の補正を省くことはできない。

\subsection{有効商における三つの補正への簡約}

$f_{14}\in J_{14}$ より、$Z_{14}\in\mathcal G_9$ を
$dP(Z_{14})=f_{14}$ と選ぶ。

\begin{theorem}[実際の二十次類の代表元]
\label{thm:v15-twentieth-effective-representative}
第14版の正規形における実際の二十次類は
\begin{equation}\label{eq:v15-twentieth-effective-representative}
 \boxed{\displaystyle
 [R_{20}]=\left[f_{20}+d^2P(X_7,Z_{14})
       -\frac12d^2P(X_8,X_8)-c\,dH(\Xi_{11})\right]
       \quad\text{in }\mathcal C_{20}/I_{20}.
 }
\end{equation}
純粋な持ち上げ $Z_{14}$、$\Xi_{11}$ の選択はこの類に影響しない。
$P$ 倍の単元項を残した持ち上げを用いても有効類は同じである。
\end{theorem}

\begin{proof}
恒等式 $f_{13}=-dP(X_8)$ を $X_8$ で微分すると
\[
 d(f_{13})(X_8)=-d^2P(X_8,X_8)-dP(DX_8\cdot X_8).
\]
最後の場は $\mathcal G_{15}$ に属するので、奇数次項を消した二階補正は
$-\tfrac12d^2P(X_8,X_8)$ に簡約される。
同様に $f_{14}=dP(Z_{14})$ を $X_7$ で微分すれば
\[
 d(f_{14})(X_7)\equiv d^2P(X_7,Z_{14})\pmod{J_{20}}.
\]

既存の所属 $f_{16}\in J_{16}$、$f_{14}\in J_{14}$ から
$U_4f_{16},V_6f_{14}\in J_{20}$ である。
さらに $f_{10}=cH-3PU_4$ と、第13版の
$dH(\mathcal G_7)\subset J_{16}$ を用いると
$U_4d(f_{10})(X_7)\in J_{20}$ である。
これで表示した三補正の式を得る。

二つの純粋な $Z_{14}$ の差 $K_9$ には $dP(K_9)=0$ が成り立ち、
これを $X_7$ で微分すれば $d^2P(X_7,K_9)\in J_{20}$ である。
$\Xi_{11}$ の差については補題\ref{lem:v15-eleventh-stabilizer-hessian}を適用する。
単元項を残した持ち上げと純粋なものとの差はEuler場であり、
$P$ の二階微分への作用は $J_{20}$、$H$ への作用は
$H\mathcal U_{10}\subset I_{20}$ に入る。
\end{proof}

\begin{corollary}[十七次・十九次の明示消去は不要]
\label{cor:v15-twentieth-without-odd-normalizers}
上の有限正規化と18次消去の後に17次・19次の項が残っていても、
式\eqref{eq:v15-twentieth-effective-representative}は全次数偶正規形の
二十次有効類を与える。
\end{corollary}

\begin{proof}
有限正規化後を $f=P+cH+R_{17}+R_{19}+R_{20}+O(a^{21})$ とし、
第14版の偶正規形と比較する接触作用素を $T$ とする。
$Tf-f$ は17次以上なので、対数補題から
$Tf-f\equiv Lf\pmod{(a)^{21}}$ である。
最小次数増分が4であるため、$L(R_{17}+R_{19})$ も21次以上である。
従って二十次までの比較は $L(P+cH)$ だけによる。
偶数増分の条件は\eqref{eq:v15-even-contact-chain}のままであり、
17次・19次消去の非零な右辺は奇数増分の条件にだけ入る。
二十次の差は $I_{20}$ に含まれる。

一般の偶化に含まれ得る六次や八次の座標場をTaylor展開すると、
それぞれの二階作用が二十次に見える場合がある。
ここでは低次を保つ全条件を満たした作用素全体の対数を用いており、
そのような項の相殺も含んだ結論である。
\end{proof}

\subsection{二階補正を小さい有限族へ落とす一般補題}

\begin{lemma}[Jacobian像を法とした二階作用]
\label{lem:v15-hessian-pairing-quotient}
$Y\in\mathcal G_r$、$Z\in\mathcal G_s$、$N=r+s+4$ とする。
$d^2P(Y,Z)$ の $\mathcal C_N/J_N$ における類は、
\[
 [dP(Y)]\in J_{r+5}/P\mathcal U_{r-1},\qquad
 [dP(Z)]\in J_{s+5}/P\mathcal U_{s-1}
\]
だけに依存する。したがって、これらの商からの対称な双線形写像を定める。
\end{lemma}

\begin{proof}
$dP(K)=0$ を $Z$ で微分すると
$d^2P(K,Z)=-dP(DK\cdot Z)\in J_N$ である。
また $dP(Y)=Pu$ のとき、$Y-au/6$ は $dP$ の核に入る。
Eulerの関係を微分して
\[
 d^2P(a u/6,Z)=\frac56u\,dP(Z)\in J_N
\]
を得る。よって表示した各商を通る。
\end{proof}

\begin{proposition}[二十次の補正を全て判定する有限族]
\label{prop:v15-finite-lower-tests}
以下の四群の全てが $J_{20}$ に含まれると仮定する。
\begin{enumerate}
 \item $H\mathcal U_{10}$。$P$ 倍の単元方向は恒等的に $J_{20}$ に入るので省ける。
 \item 第13版の $J_{16}/P\mathcal U_{10}$ の基底を持ち上げる
 二十一の場 $\Gamma_j^{(11)}$ に対する $dH(\Gamma_j^{(11)})$。
 \item 第14版の $J_{12}/P\mathcal U_6$ の基底を持ち上げる二つの場
 $\Gamma_i^{(7)}$ と $J_{14}/P\mathcal U_8$ の基底を持ち上げる七つの場
 $\Gamma_j^{(9)}$ に対する、十四の共変量
 \[
 \operatorname{Hess}_F(q)
    (\delta q\Gamma_i^{(7)},\delta q\Gamma_j^{(9)}).
 \]
 \item $J_{13}=\mathcal C_{13}$ の非零生成元を持ち上げる一つの場
 $\Gamma^{(8)}$ に対する
 \[
 \operatorname{Hess}_F(q)(\delta q\Gamma^{(8)},\delta q\Gamma^{(8)}).
 \]
\end{enumerate}
このとき $I_{20}=J_{20}$ であり、実際の二十次類は
$[R_{20}]=[d_{20}^{\mathrm T}]$ となる。
本命題は二十次の有限所属判定が完了するまでは条件付きである。
\end{proposition}

\begin{proof}
(1),(2)と補題\ref{lem:v15-eleventh-stabilizer-hessian}から
$dH(\mathcal G_{11})\subset J_{20}$ である。
実際、任意の場から二十一の基底像とEuler単元場を除けば $dP$ の核に入る。
従って二つの安定化方向も $dH(\Xi_{11})$ も $J_{20}$ に入り、$I_{20}=J_{20}$ となる。

合成関数の微分により
\[
 d^2P(Y,Z)=\operatorname{Hess}_F(q)(\delta qY,\delta qZ)
                      +dF(d^2q(Y,Z)).
\]
ここでは $d^2q(Y,Z)\in\mathcal D_{16}$ なので最後の項は $J_{20}$ に入る。
前補題と $\dim(J_{12}/P\mathcal U_6)=2$、
$\dim(J_{14}/P\mathcal U_8)=7$ により混合補正は(3)の十四方向で尽くされる。
また $\dim J_{13}=1$、$\mathcal U_7=0$ なので、
二つの八次場による対称二階補正は(4)の一方向だけで尽くされる。
これで式\eqref{eq:v15-twentieth-effective-representative}の三補正が全て消える。
\end{proof}

\begin{remark}[判定が零でない場合]
上の有限族のいずれかが通常商で非零なら、その所属を仮定してはならない。
その場合でも有効像は\eqref{eq:v15-effective-image}で決まり、
実際の類は\eqref{eq:v15-twentieth-effective-representative}で計算できる。
特に八次場の平方補正は元のTateの $K_8$ から直接評価できる。
混合補正は $d(f_{14})(X_7)$ として方向微分してもよく、
その違いは $J_{20}$ に入る。
\end{remark}

% END contact/contact_twentieth_ja.tex


% BEGIN detector/v15_detector_notes_ja.tex
% 第15版用。S20の328次元独立性はrank証明書の節を参照する。
\section{20次の既知接触像を判定する三点汎関数}
\label{sec:v15-three-point-detectors}

20次の普遍境界空間は330次元である。
本節では、既に構成した328次元の接触像の部分空間への所属を、
三つの値だけで判定する二つの汎関数を与える。
この所属判定は、通常接触写像の全階数やKoszul完全性を必要としない。
通常接触像全体との一致がまだ証明されていない段階では、
汎関数の非零値だけから非零接触類を結論しない。

\subsection{既知像と二つの補元}

164個の15次座標場を $\Gamma_{15,\alpha}$ とし、
\[
 \mathcal S_{20}:=
 P\mathcal U_{14}
 +\left\langle dP(\Gamma_{15,\alpha}):1\le\alpha\le164\right\rangle
 \subset\mathcal J_{20}:=dP(\mathcal G_{15})
\]
とおく。ここで $P\mathcal U_{14}\subset\mathcal J_{20}$ は
$dP(a u/6)=Pu$ というEulerの等式による。
一般の $P=0$ の整数入力における164列の非零minorにより、
これらの列は $P\mathcal U_{14}$ を法として独立である。
$\dim\mathcal U_{14}=164$ と $P$ 倍の単射性から
\begin{equation}\label{eq:v15-known-image-dimension}
 \dim\mathcal S_{20}=328
\end{equation}
を得る。
この下界の証明と、通常像 $\mathcal J_{20}$ の上界の証明は区別する。

第11版と同じ正規化の6頂点単元グラフを $J_6$ とし、
\[
 W_1=\Theta_8U_4^3,\qquad W_2=\Theta_8J_6^2
\]
とおく。両者は $F$ 次数15、$a$ 次数20を持つ自然な共変量であり、
$\mathcal C_{20}$ に属する。

\subsection{三つの入力}

\begin{align*}
 F_D&=x_0^3+x_1^3+x_2^3+x_3^3-2x_4^3,\\
 F_H&=x_0^3+x_1^3+x_2^3+x_3^3+4x_4^3-6x_2x_3x_4,\\
 a_b&=b\,e_0\wedge e_1
       +b^{-1}(e_2-e_4)\wedge(e_3-e_4)
\end{align*}
とおく。Pfaffianの規約は前節までと同じである。
任意の $b\ne0$ に対し
\[
 q(a_b)=(0,0,1,1,1),\qquad F_D(q(a_b))=F_H(q(a_b))=0
\]
となる。三つの評価点を
\[
 z_1=(F_D,a_1),\qquad z_2=(F_D,a_2),\qquad z_3=(F_H,a_1)
\]
とする。$F_D$ は非零対角係数を持つ滑らかな三次式である。
$F_H$ の最後の三変数の部分はHesse型であり、
対角係数の積と混合tensor係数の判別条件は $4+8(-1)^3=-4\ne0$ である。
最初の二変数は分離しているので $F_H$ も滑らかである。

この族では、グラフの整数縮約から
\[
\begin{array}{c|ccc}
 &\Theta_8&U_4&J_6\\\hline
(F_D,a_b)&-2b^2&0&24b^6\\
(F_H,a_b)&-4b^2&5b^4&528b^6
\end{array}
\]
を得る。$F_H$ のtensorの規約は
$F_{234}=F_{243}=\cdots=-1$ であり、
多項式の混合項は $-6x_2x_3x_4$ となる。
この係数6を省略しない。

\begin{proposition}[三点の厳密な双対汎関数]
\label{prop:v15-exact-dual-detectors}
$\operatorname{ev}_i(f)=f(z_i)$ とおき、
\begin{align}
 \rho_1&=\frac{-22168\operatorname{ev}_1+82\operatorname{ev}_2
                    -1365\operatorname{ev}_3}{682500},
                    \label{eq:v15-rho1}\\
 \rho_2&=\frac{4\operatorname{ev}_1-\operatorname{ev}_2}{18869760}
                    \label{eq:v15-rho2}
\end{align}
と定める。このとき
\[
 \rho_i(\mathcal S_{20})=0,\qquad
 \rho_i(W_j)=\delta_{ij},\qquad
 \ker(\rho_1,\rho_2)=\mathcal S_{20}
\]
が成り立つ。
\end{proposition}

\begin{proof}
$P(z_i)=0$ なので、全ての $P\mathcal U_{14}$ の方向は三点とも零である。
164個のグラフについて、三つの値を有理数上で再計算した。
各列の値は厳密に $(c,4c,-16c)$ の形を持つ。
従って両汎関数は全ての164列を消す。
この等式は有限体での試験に留めず、非有界整数のtensor縮約と有理数演算で検証した。

選択した一つの接触列と二つの補元の評価行列は
\[
 \begin{pmatrix}
 -96&0&-1152\\
 -384&0&-18874368\\
 1536&-500&-1115136
 \end{pmatrix}
\]
であり、その行列式は
\[
 905748480000\ne0
\]
である。逆行列の対応する二行が
式\eqref{eq:v15-rho1}--\eqref{eq:v15-rho2}を与え、
$\rho_i(W_j)=\delta_{ij}$ を満たす。
両汎関数は $\mathbb Q$ 上で定義され、複素数へ係数拡大して用いる。
従って $(\rho_1,\rho_2):\mathcal C_{20}\to\mathbb C^2$ は階数2で、
その核の次元は $330-2=328$ である。
式\eqref{eq:v15-known-image-dimension}と包含
$\mathcal S_{20}\subset\ker(\rho_1,\rho_2)$ から核の等式を得る。
ここでは $\mathcal S_{20}=\mathcal J_{20}$ は仮定していない。
\end{proof}

従って任意の $f\in\mathcal C_{20}$ に対し
\[
 f-\rho_1(f)W_1-\rho_2(f)W_2\in\mathcal S_{20}\subset\mathcal J_{20}
\]
である。特に二つの値がともに零であることは、接触像への所属の十分条件となる。
通常接触像が既知像より大きい可能性を残しても、この十分条件は有効である。

\subsection{低次正規化から来る補正の消滅}

$h_F(p)=\det\operatorname{Hess}_F(p)$、$H(F,a)=h_F(q(a))$ と書く。
$dH$ は $a$ 変数方向の微分、$d_qh_F$ は $q$ 変数方向の微分とする。
三つの入力において、$\operatorname{Hess}_F(q)$ の最初の二行と二列は零であり、
残りの $3\times3$ ブロックは可逆である。
従って全Hessianの階数は3で、
\[
 H(z_i)=0,\qquad d_qh_F(q(a))(z_i)=0
\]
となる。後者は $5\times5$ 行列の全ての $4\times4$ 余因子が零であることによる。
従って任意の $u\in\mathcal U_{10}$ と $Z\in\mathcal G_{11}$ に対し
\[
 (Hu)(z_i)=0,\qquad
 dH(Z)(z_i)=d_qh_F(q(a))(dq(Z))(z_i)=0
\]
である。これらは $\mathcal C_{20}$ の元なので、
命題\ref{prop:v15-exact-dual-detectors}から
\begin{equation}\label{eq:v15-hessian-secondary-in-known-image}
 H\mathcal U_{10}+dH(\mathcal G_{11})\subset\mathcal S_{20}
\end{equation}
を得る。ここでは $\mathcal U_{10}$ や $\mathcal G_{11}$ の評価基底を列挙する必要はない。

次に $\mathcal J_{12}/P\mathcal U_6$ の2個の代表場と、
$\mathcal J_{14}/P\mathcal U_8$ の7個の代表場を、それぞれ
$X_{7,i}$、$Y_{9,j}$ とする。
これらは第14版で用いた同じグラフ場である。
14個の普遍境界共変量
\[
 B_{ij}:=\operatorname{Hess}_F(q)
          \bigl(dq(X_{7,i}),dq(Y_{9,j})\bigr)
 \quad (1\le i\le2,\ 1\le j\le7)
\]
について三点で厳密評価し、全ての $i,j$ に対して
\[
 \rho_1(B_{ij})=\rho_2(B_{ij})=0
\]
を確認した。従って $B_{ij}\in\mathcal S_{20}$ である。
同様に、既知の1次元像 $\mathcal J_{13}$ を生成する8次場 $X_8$ の代表を一つ選ぶと、
\[
 \operatorname{Hess}_F(q)(dq(X_8),dq(X_8))\in\mathcal S_{20}
\]
となる。この代表場の $dP$ が非零であることも、
別の一般整数入力で有理数上の縮約により確認した。

これらの有限検査から、全ての必要な場へ広げる理由を記す。
$dP(X)=0$ なら、その恒等式を場 $Y$ 方向に微分することで
\[
 P''(X,Y)=-dP(DX[Y])\in\mathcal J_{20}
\]
を得る。Euler単元場が一方に入る場合も、
$\operatorname{Hess}_F(q)(q,v)=2\,dF(v)$ により接触像に入る。
残る方向は上の $2\times7$ 個と8次像の1個で尽くされる。
さらに
\[
 P''(X,Y)=\operatorname{Hess}_F(q)(dq(X),dq(Y))
            +dF\bigl(d^2q(X,Y)\bigr)
\]
の最後の項は $d^2q(X,Y)\in\mathcal D_{16}$ と
全次数Pfaffian全射定理から $\mathcal J_{20}$ に入る。
従って
\begin{equation}\label{eq:v15-quadratic-corrections-contact}
 P''(\mathcal G_7,\mathcal G_9)
  +P''(\mathcal G_8,\mathcal G_8)\subset\mathcal J_{20}
\end{equation}
を得る。
この議論の途中で生じる任意の接触項について、
まだ両 $\rho_i$ が零になると仮定しているのではない。
それらは直接 $\mathcal J_{20}$ に属するため、接触同値の計算で捨ててよい。

\subsection{実際の20次項への適用と再現範囲}

filtered contact作用の計算で得られた代表は、既存の規約では
\[
 R_{20}^{\mathrm{eff}}
 =d_{20}^{T}+P''(X_7,Z_{14})
       -\tfrac12P''(X_8,X_8)-\tfrac1{96}dH(\Xi_{11})
       \pmod{\mathcal J_{20}}
\]
という形に書ける。
式\eqref{eq:v15-hessian-secondary-in-known-image}と
\eqref{eq:v15-quadratic-corrections-contact}により
\[
 [R_{20}^{\mathrm{eff}}]=[d_{20}^{T}]
 \quad\text{in }\mathcal C_{20}/\mathcal J_{20}
\]
である。三点で独立に再計算したfresh Tate値は、
\[
 \bigl(d_{20}^{T}(z_1),d_{20}^{T}(z_2),d_{20}^{T}(z_3)\bigr)
 =(A,4A,-16A),\qquad
 A=\frac{1112989093269228481}{187535250000000}
\]
である。式\eqref{eq:v15-rho1}--\eqref{eq:v15-rho2}に代入すると
$\rho_1(d_{20}^{T})=\rho_2(d_{20}^{T})=0$ となる。
従って $d_{20}^{T}\in\mathcal S_{20}\subset\mathcal J_{20}$ であり、
実際の20次項は接触像に入り消去できる。
この結論には、通常商の次元が2であることの証明を先に必要としない。
fresh Tateブロックの再構成と次数の切断の検査は、Tate計算の節とその証明書で行う。

\texttt{v15\_rho20\_three\_cases.json} は三つの入力を保持する。
\nolinkurl{v15_rho20_dual_certificate.py} は
164接触列、二つの補元、14個の混合Hessianと1個の平方を有理数上で再計算する。
分母を持つ $a$ は一度 $Da$ にして整数縮約し、20次の斉次性により $D^{20}$ で割り戻す。
$U_4,J_6,\Theta_8$ についても各々の次数で同じ処理を行う。
探索時には、先に完成した二つの対角点の厳密値を再利用してHesse点を計算した。
最終証明書では全三点をfreshに再計算した。
既定の再実行も全三点をfreshに再計算する。
値を入力とする監査と、tensor縮約そのものの再計算は区別して記録する。

% END detector/v15_detector_notes_ja.tex


% BEGIN tate/v15_tate_notes_ja.tex
\section{二十次Tate係数の新規再構成とHesse型Clifford計算}
\label{sec:v15-tate-twentieth}

本節では、前節の三点における未正規化係数 $d_{20}^{\mathrm T}$ を、
TateのMaurer--Cartan再帰から新たに計算する。
既存の十八次ブロックを二十次に流用するのではなく、
必要な内部変数の次数を改めて導き、全ての入力点で最初から再構成した。
二つの対角点に加え、混合項 $-6x_2x_3x_4$ を持つHesse点も扱う。

\subsection{二十次に必要な切断の範囲}

補正を $A=\sum_{r\ge1}A_r$、定数ブロックを
$K_r:\Lambda^4V\to V$ と書く。規約は
\[
 R_n=\sum_{i=1}^{n-1}A_iA_{n-i},\qquad A_n=-hR_n,
 \qquad d^{\mathrm T}=\tfrac12\det\!\left(\sum_{r\ge1}K_r\right)
\]
である。各 $K_r$ は $a$ に $r$ 次で、$K_1=-a$、
$K_2=\operatorname{Hess}_F(q)/3$ となる。

\begin{lemma}[二十次の厳密な内部次数切断]
\label{lem:v15-tate-twentieth-truncation}
二十次の行列式係数には $K_1,\ldots,K_{16}$ だけが必要である。
これらは、$A_1,\ldots,A_{12}$ を全て求めた後、
$R_{13},R_{14},R_{15},R_{16}$ をそれぞれ内部 $x$ 次数
$4,3,2,1$ 以下で計算すれば厳密に得られる。
そのとき $A_{13},A_{14},A_{15},A_{16}$ はそれぞれ
$x$ 次数 $3,2,1,0$ 以下まで正確である。
\end{lemma}

\begin{proof}
行列式の各単項式は、次数が少なくとも1の五つのブロックの積である。
その総次数を20以下にするには、各因子の次数は高くとも16でよい。
Clifford積の収縮は $x$ 次数を下げず、Eulerホモトピー $h$ はちょうど1下げる。
また内部重みの関係
\[
 2\deg_x-\#\theta+\#\eta=3
\]
により、$A_r$ の $x$ 次数は常に4以下である。
従って $A_1,\ldots,A_{12}$ の全計算で前半を確保できる。
後半で既に求めた $A_i$ は $x$ 次数 $16-i$ 以下まで正確である。
$R_n$ の積で $i\le n-1$ なら $16-i\ge17-n$ なので、
$x$ 次数 $17-n$ 以下の曲率に必要な係数は全て保存されている。
これを帰納的に用いれば、表示した切断で各定数ブロックを失わない。
\end{proof}

第14版の $A_{12}$ は $x$ 次数2以下までの計算であった。
今回の全 $A_{12}$ は新たに計算しており、前版の切断を全項と見なしてはいない。

\subsection{三変数の混合項を持つ自然なClifford積}

$F=x_0^3+x_1^3+f(x_2,x_3,x_4)$ とし、
\[
 B_{ij}(x)=\tfrac16\frac{\partial^2F}{\partial x_i\partial x_j}
\]
とおく。自然なClifford生成元は
\[
 \eta_i=\iota_i-\sum_jB_{ij}(x)\theta_j,
 \qquad \{\eta_i,\eta_j\}=-2B_{ij},\qquad
 \{\theta_i,\eta_j\}=\delta_{ij}
\]
を満たす。ここで $\theta_i$ は外積作用、$\iota_i$ は収縮作用である。
記号としての $\eta$ は外積で反対称化し、$\theta$ を通常順に置く。
これが既存の普遍同変なEulerホモトピーを適用する記号基底である。

\begin{lemma}[Hesse型の外積記号とPBW基底]
\label{lem:v15-hesse-symbol-product}
混合項が $x_2,x_3,x_4$ のみに含まれる場合、
$0,1$ の各一変数Clifford代数と、最後の三変数のClifford代数との
次数付きテンソル積によって全代数を計算できる。
最後の三変数で外積記号から順序付きPBW基底への変換を $S$ とすると
\[
 S=1+C,\qquad C^2=0,\qquad S^{-1}=1-C.
\]
従って自然な記号の積は
\[
 \mu_{\mathrm{ext}}=S^{-1}\circ\mu_{\mathrm{PBW}}\circ(S\otimes S)
\]
で計算できる。
\end{lemma}

\begin{proof}
異なるブロック間のClifford係数は零なので、ブロック間には次数付き符号だけが現れる。
三変数ブロックは六つの奇生成元を持ち、PBW基底は $2^6=64$ 個である。
$i<j<k$ について外積記号は
\[
 \eta_i\wedge\eta_j=\eta_i\eta_j+B_{ij},
\]
\[
 \eta_i\wedge\eta_j\wedge\eta_k
  =\eta_i\eta_j\eta_k+B_{ij}\eta_k-B_{ik}\eta_j+B_{jk}\eta_i
\]
となる。$C$ はこの一対の収縮であり、最後の三変数には
$\eta$ が高々三つしかないので二対の収縮はできない。
従って $C^2=0$ であり、表示した逆変換が厳密に成立する。
$C$ は $x$ 次数を上げるため、前補題の切断の議論も保たれる。
$x=0$ では $S=1$ なので、定数ブロックの取り出しは従来の規約と一致する。
\end{proof}

実装では、三変数ブロックの語をClifford関係で正規順序へ還元してから、
上の $S^{-1}$ を適用する。
PBW順序のままEulerホモトピーを適用すると自然な外積記号と異なるため、
この変換を省略しない。
基本作用素は
\[
 D=\sum_i x_i\eta_i+\frac23\sum_i\partial_iF\,\theta_i
\]
であり、$F_H$ の混合項からは各対応する二次単項式に係数 $-4$ が現れる。
これを含む完全な $B_{ij}$ を用いて、独立な32次元外積表現も実装した。
独立表現ではPBW還元も $S$ の式も使わず、実際の $\eta_i$ 作用を反対称化する。

\subsection{三点の二十次係数}

\begin{proposition}[新規に再構成した三点Tate値]
\label{prop:v15-fresh-tate-twentieth}
第\ref{sec:v15-three-point-detectors}節の三点 $z_1,z_2,z_3$ に対し、
\[
 A_*=\frac{1112989093269228481}{187535250000000}
\]
とおくと
\[
 \boxed{\bigl(d_{20}^{\mathrm T}(z_1),d_{20}^{\mathrm T}(z_2),
                       d_{20}^{\mathrm T}(z_3)\bigr)
          =A_*(1,4,-16).}
\]
従って $\rho_1(d_{20}^{\mathrm T})=\rho_2(d_{20}^{\mathrm T})=0$ であり、
$d_{20}^{\mathrm T}\in\mathcal S_{20}\subset\mathcal J_{20}$ である。
\end{proposition}

\begin{proof}
各点で $A_1$ から再帰を開始し、前補題の範囲を有理数上で再構成した。
各段の曲率の閉性とMaurer--Cartan方程式を保存した次数範囲で直接照合した。
二十次係数は
\[
\begin{array}{c|c}
 z_1&1112989093269228481/187535250000000\\
 z_2&1112989093269228481/46883812500000\\
 z_3&-1112989093269228481/11720953125000
\end{array}
\]
である。五行の行列式を120置換による展開と列部分集合による動的計画法で
それぞれ計算し、0次から20次まで全係数の一致を確認した。
最後の二つの分母は最初の分母のそれぞれ $1/4,1/16$ である。
両汎関数の整数分子は、この比に対して
\[
 -22168+82\cdot4-1365(-16)=0,\qquad4-4=0
\]
となる。命題\ref{prop:v15-exact-dual-detectors}を用いて所属を得る。
\end{proof}

\paragraph{独立照合の範囲.}
三点以外にも、前版で用いた $F=x_0^3+x_1^3+x_2^3+x_3^3-4x_4^3$ の
六点を二十次まで新たに再構成した。
その標準点では、各係数を個別の有理数として計算する方法と、
共通分母を分離して整数で累積する方法の両方を実行した。
全16ブロック、0次から20次までの全行列式係数、および低次場の値が完全に一致した。
同じ点の18次以下は第14版の値にも一致した。
従って新規計算は異なる九点、標準点の独立再構成を含めて十回である。

Maurer--Cartan再帰と行列式の照合は有理数上で行った。
独立な32次元外積作用素による0次から12次の照合、および定数ブロックの
独立作用素との照合は素数 $1000003$ 上で行っている。
この二つの検証範囲を混同しない。
Hesse型Clifford生成元の反交換関係は有理数上で全て確認した。

\paragraph{再現用ソース.}
\texttt{source\_v15/tate/} 内に、対角入力の
\nolinkurl{v15_tate_twentieth.py}、混合項の積を実装する
\nolinkurl{v15_hesse_clifford.py} とその再帰、
三点の最終所属を確認する
\nolinkurl{v15_actual_twentieth_certificate.py} を置いた。
公開入力、全定数ブロック、有理数の行列式係数、各再帰段の記録を添付する。
最終所属の監査は保存済みブロックから行列式を再計算する。
それと、最初からTate再帰を再構成する実行とは区別した。

% END tate/v15_tate_notes_ja.tex


% BEGIN contact/twenty_second_normal_form_ja.tex
% 第15版の確定稿。三点のfresh Tate係数と既知像への所属に基づく。
% 通常接触像の全階数が確定したことは仮定せず、主張もしない。

\section{二十次の消去と二十二次開始の正規形}
\label{sec:v15-twenty-second}

本節では、通常接触像の一つの既知部分空間への所属から、
実際の二十次項を消去する。
以下、$P=P_F$、$H=H_{10}$、$c=1/96$、
$J_{20}=dP(\mathcal G_{15})$ とし、
$\mathcal S_{20}\subset J_{20}$ は
命題\ref{prop:v15-exact-dual-detectors}の328次元の既知像とする。

\begin{theorem}[二十二次開始の普遍代数的収束偶正規形]
\label{thm:v15-twenty-second-normal-form}
二十次の有効接触像は通常接触像と一致し、$I_{20}=J_{20}$ である。
実際の境界の二十次類は零であり、任意の滑らかな三次式 $F$ に対して
\begin{equation}\label{eq:v15-twenty-second-normal-form}
 \boxed{\displaystyle
 d^{\mathrm T}\sim
 P_F+\frac1{96}\det\operatorname{Hess}_F(q(a))
       +\sum_{\substack{n\ge22\\ n\ \mathrm{even}}}R_n,
 \qquad R_n\in\mathcal C_n
 }
\end{equation}
という正規形が存在する。
接触変換は普遍同変な代数的冪級数で、座標変更の線形項は恒等写像、
単元の定数項は1であり、各形式係数は $F$ の係数に多項式的である。
各固定 $F$ では、変換・逆変換・正規形の級数は原点近傍で収束する。
この結論は $\mathcal S_{20}=J_{20}$ や通常商の次元の確定を必要とせず、
二十二次以降の消去を主張しない。
\end{theorem}

\begin{proof}
命題\ref{prop:v15-exact-dual-detectors}により
\[
 \ker(\rho_1,\rho_2)=\mathcal S_{20}\subset J_{20}.
\]
三点で $H$ と $d_qh$ が点ごと零であることから、
式\eqref{eq:v15-hessian-secondary-in-known-image}は
\[
 H\mathcal U_{10}+dH(\mathcal G_{11})
       \subset\mathcal S_{20}\subset J_{20}
\]
を与える。従って定理\ref{thm:v15-effective-twentieth}の
二つの追加安定化方向も通常像に入り、$I_{20}=J_{20}$ である。
さらに式\eqref{eq:v15-quadratic-corrections-contact}により
\[
 d^2P(\mathcal G_7,\mathcal G_9)
       +d^2P(\mathcal G_8,\mathcal G_8)\subset J_{20}.
\]
定理\ref{thm:v15-twentieth-effective-representative}の
混合二階補正、八次場の平方補正、十一次場のHessian作用は、
以上により全て $J_{20}$ に属する。
よって、有限正規化後の二十次係数について
\begin{equation}\label{eq:v15-raw-and-normalized-twentieth-class}
 [R_{20}^{\mathrm{fin}}]=[d_{20}^{\mathrm T}]
           \quad\text{in }\mathcal C_{20}/J_{20}
\end{equation}
である。

次に、三点 $z_1,z_2,z_3$ の実際のTate再帰をそれぞれ再構成する。
対角三次式の二点とHesse型三次式の一点から得られた未正規化値は
\begin{equation}\label{eq:v15-three-fresh-twentieth-values}
 \bigl(d_{20}^{\mathrm T}(z_1),d_{20}^{\mathrm T}(z_2),
                d_{20}^{\mathrm T}(z_3)\bigr)
       =(\lambda_{20},4\lambda_{20},-16\lambda_{20}),
 \qquad
 \lambda_{20}=\frac{1112989093269228481}{187535250000000}.
\end{equation}
これらは $A_1,\ldots,A_{12}$ と必要な $A_{13},\ldots,A_{16}$ の切断を
新たに作り、十六個の定数ブロックから求めた値である。
Hesse型入力では自然なClifford積と独立な外積作用素を用いており、
三値の比を仮定して係数を決めてはいない。
命題\ref{prop:v15-exact-dual-detectors}の整数重みを用いると
\[
 -22168\lambda_{20}+82(4\lambda_{20})
       -1365(-16\lambda_{20})=0,\qquad
 4\lambda_{20}-4\lambda_{20}=0.
\]
従って $\rho_1(d_{20}^{\mathrm T})=\rho_2(d_{20}^{\mathrm T})=0$ であり、
\[
 d_{20}^{\mathrm T}\in\mathcal S_{20}\subset J_{20}.
\]
式\eqref{eq:v15-raw-and-normalized-twentieth-class}から
実際の二十次類は零である。

第14版の定理\ref{thm:v14-twentieth-normal-form}の全次数偶正規形を
\[
 f_{\mathrm{ev}}=P+cH+R_{20}^{\mathrm{ev}}
                       +R_{22}^{\mathrm{ev}}+\cdots
\]
と書く。
系\ref{cor:v15-twentieth-without-odd-normalizers}により
$R_{20}^{\mathrm{ev}}-R_{20}^{\mathrm{fin}}\in I_{20}=J_{20}$ である。
従って $R_{20}^{\mathrm{ev}}\in J_{20}$ であり、
普遍多項式な場 $\Xi_{15}\in\mathcal G_{15}$ を
\[
 dP(\Xi_{15})=R_{20}^{\mathrm{ev}}
\]
となるように選べる。
これは有限次元の普遍共変量空間での線形な持ち上げであり、
$F$ の判別式や評価点の値による除算を伴わない。

奇座標変更 $a\mapsto a-\Xi_{15}(F,a)$ で引き戻せば二十次が消える。
$\Xi_{15}$ は奇数次なので、全次数の偶性は保たれる。
$H$ への最初の寄与は24次、$P$ の二階Taylor項は34次からであり、
既定の低次正規形は変わらない。
従って次に残り得る次数は22である。

今回加えた座標変更は有限な多項式で、原点でのJacobianが恒等写像である。
その形式逆は代数的冪級数であり、固定した $F$ では解析的にも収束する。
第14版で構成した代数的・収束的な接触変換との有限合成は、
これらの性質および普遍同変性と係数の $F$ 多項式性を保つ。
これで所要の正規形を得る。
\end{proof}

\begin{remark}[通常商の同定との区別]
二つの共変量 $W_1=\Theta_8U_4^3$ と $W_2=\Theta_8J_6^2$ は、
$\mathcal C_{20}/\mathcal S_{20}$ の基底を与える。
通常接触像が $\mathcal S_{20}$ より大きい可能性を残す段階では、
これらを $\mathcal C_{20}/J_{20}$ の基底とは呼ばない。
同様に、$\rho_i$ の非零値から非零接触類を結論しない。

本定理で必要なのは、二つの汎関数値 $\rho_1,\rho_2$ がともに零になると
既知像 $\mathcal S_{20}$ に入るという十分条件である。
これにより実際の二十次項の消去は確定する。
通常接触商の完全な同定と、二十二次以降の普遍商・実際の類の決定は
別の問題として残る。
\end{remark}

% END contact/twenty_second_normal_form_ja.tex

\clearpage
% BEGIN theory/v16_plucker_hessian_structure_ja.tex
% 第16版：第15版で未採用だった Pluecker/Hessian メモの証明。
% source_v13, source_v14, source_v15 は変更しない。
% 外部入力：Fisher, The invariants of a genus one curve, Theorem 4.4(i), n=3。
% 追加 bibliography key: Fisher2008（別ファイルに収録）。

\section{Pl\"ucker錐とHessian作用の全次数簡約}
\label{sec:v16-plucker-hessian-structure}

本節では、第15版では未採用だった構造案を証明する。
通常接触像の階数や、Pfaffian代入後のKoszul複体の完全性は仮定しない。
全ての係数は $F$ に多項式的であり、接触作用の間に $F$ を動かさない。
以下の全次数の包含は、実際の高次Tate係数の消滅を意味するものではない。

\subsection{線束の規約とPl\"uckerイデアル}

$\dim V=5$、$\Delta=\det V$ とし、
\[
 A=\Lambda^2V,\qquad W=A\otimes\Delta^{-1},\qquad
 E=V^*\otimes\Delta^{-1},\qquad L=\Delta^{-3}
\]
と置く。変数は $F\in\operatorname{Sym}^3V$、$a\in W$ であり、
\[
 q(a)=\frac12a\wedge a\in E,\qquad P=F(q(a))\in L.
\]
自然な同一視
$\operatorname{Sym}^3V=\operatorname{Sym}^3E^*\otimes L$
を用いれば、$F$ は $E$ 上の $L$ 値三次式である。
既存の $\mathcal G_r,\mathcal D_{r+1},\mathcal U_s,\mathcal C_n$ は、
それぞれ $W,E,\mathbb C,L$ 値の普遍同変多項式であり、その二重次数は
\[
\begin{array}{c|cc|c}
 &\deg_F&\deg_a&\text{値域}\\\hline
 \mathcal G_r&r-1&r&W\\
 \mathcal D_{r+1}&r-1&r+1&E\\
 \mathcal U_s&s&s&\mathbb C\\
 \mathcal C_n&n-5&n&L
\end{array}
\]
である。$J_n=dP(\mathcal G_{n-5})$ と書く。
Euler場によって $P\mathcal U_{n-6}\subset J_n$ である。

体積形式を選んで $a_{ij}$ を座標とし、
\[
 \Sigma=\{a\in W:\operatorname{rank}a\le2\},\qquad
 \mathfrak I=(q_0,q_1,q_2,q_3,q_4)
\]
と置く。以下での $\mathfrak I$ は、必要に応じて
$\mathbb C[F,a]$ へ拡大したイデアルも表す。

\begin{lemma}[Pl\"ucker錐の被約なイデアル]
\label{lem:v16-plucker-prime-ideal}
$\mathfrak I$ は素イデアルであり、$\Sigma$ のイデアルに等しい。
従って、普遍多項式が全ての $(F,a)$、$a\in\Sigma$ で消えることと、
それが $\mathbb C[F,a]$ の $\mathfrak I$ に入ることは同値である。
\end{lemma}

\begin{proof}
次の五式を $g_1,\ldots,g_5$ とする。
\[
\begin{split}
 g_1&=a_{12}a_{34}-a_{13}a_{24}+a_{14}a_{23},\\
 g_2&=a_{02}a_{34}-a_{03}a_{24}+a_{04}a_{23},\\
 g_3&=a_{03}a_{14}-a_{04}a_{13}-a_{01}a_{34},\\
 g_4&=a_{02}a_{14}-a_{04}a_{12}-a_{01}a_{24},\\
 g_5&=a_{02}a_{13}-a_{03}a_{12}-a_{01}a_{23}.
\end{split}
\]
これらは順に $q_0,-q_1,-q_2,q_3,-q_4$ である。
変数の順序を
\[
 a_{02}>a_{03}>a_{04}>a_{12}>a_{13}>a_{14}>
 a_{23}>a_{24}>a_{34}>a_{01}
\]
とする辞書式順序では、この五式はGr\"obner基底となる。
実際、十組の $S$ 多項式を表示した順に除法すると全て余りが零になる。
その十組の除法の商全体を
\nolinkurl{v16_structure_certificate.json} に記録し、
\nolinkurl{v16_structure_certificate.py} は各等式を整数・有理数で再構成する。
初期単項式はいずれも $a_{01}$ を含まないので、
$\mathfrak I:a_{01}^{\infty}=\mathfrak I$ である。

$a_{01}$ を可逆にすると、最後の三式から
\[
 a_{ij}=\frac{a_{0i}a_{1j}-a_{0j}a_{1i}}{a_{01}}
 \quad(2\le i<j\le4)
\]
を得る。これを最初の二式へ代入すると恒等的に零となるため、局所化した商は
\[
 \mathbb C[a_{01}^{\pm1},a_{02},a_{03},a_{04},a_{12},a_{13},a_{14}]
\]
である。飽和性により元の商はこの整域へ単射するので、$\mathfrak I$ は素である。
この開集合の点は二つのベクトルの外積であり、その閉包は階数二以下の錐である。
また五つの4行4列Pfaffianが同時に零となる交代行列の階数は二以下である。
従って零点集合は正確に $\Sigma$ である。
$F$ の係数を加える操作は多項式拡大なので、最後の主張も従う。
\end{proof}

\begin{lemma}[普遍Pfaffian方向は錐上で消える]
\label{lem:v16-plucker-vanishing-D}
全ての $r\ge1$ と $w\in\mathcal D_{r+1}$ について
\[
 w(F,a)|_{a\in\Sigma}=0.
\]
さらに、$F,a$ にそれぞれ $r-1$ 次の、$\operatorname{End}(E)$ 値の
普遍同変多項式 $M(F,a)$ を選んで
\begin{equation}\label{eq:v16-equivariant-Mq}
 w=Mq
\end{equation}
と書ける。そのトレースは $\operatorname{tr}M\in\mathcal U_{r-1}$ である。
\end{lemma}

\begin{proof}
$V=L_2\oplus Q_3$、$L_2=\langle e_0,e_1\rangle$、
$Q_3=\langle e_2,e_3,e_4\rangle$ と分け、
$a_0=e_0\wedge e_1\otimes\Delta^{-1}$ を取る。
ここで右辺の $\Delta^{-1}$ の元は選んだ体積による基底を意味する。
一パラメータ部分群
\[
 \lambda(t)=\operatorname{diag}(t,t,1,1,1)
\]
は、外積の重み2と $\Delta^{-1}$ の重み $-2$ が相殺するため $a_0$ を固定する。
入力 $F\in\operatorname{Sym}^3V$ の $\lambda$ 重みは0、1、2、3である。
従って $F$ に多項式的な式には負の重みは現れない。
一方、出力 $E=V^*\otimes\Delta^{-1}$ の重みは、
$L_2^*$ 成分で $-3$、$Q_3^*$ 成分で $-2$ である。
同変性から $w(F,a_0)=0$ を得る。
非零の階数二の点は一つの $\mathrm{GL}(V)$ 軌道をなすので、$w$ は $\Sigma$ 全体で消える。
これは determinant twist を含めた証明であり、固定した滑らかな $F$ のみに関する議論ではない。

前補題から各成分は $\mathfrak I$ に入り、二重次数をそろえて $w=Mq$ と書ける。
この表示を与える有限次元の係数空間の全射は $\mathrm{GL}(V)$ 同変である。
標数零の完全可約性により、同変な $w$ の持ち上げ $M$ を同変に選べる。
つまり $M$ の値域は $E\otimes E^*=\operatorname{End}(E)$ であり、
変換則は共役作用である。従ってトレースはスカラー値で、その二重次数から
$\mathcal U_{r-1}$ に属する。$M$ の一意性は主張しない。

表現論的にも、錐の $a$ 次数 $r+1$ の成分は $S_{(r+1,r+1)}V$ であるのに対し、
trace線を外した出力の最高重みは $(r,r,r,r,r-1)$ である。
前者の第一行は後者より大きいので、多項式 $F$ 表現とのテンソル積から出力は現れない。
これは上の一パラメータ部分群による消失と整合する。
\end{proof}

\subsection{adjugate恒等式と同変な持ち上げ}

$p\in E$ を独立変数として
\[
 g(p)=dF_p,\qquad T(p)=\operatorname{Hess}_F(p),\qquad
 h(p)=\det T(p),\qquad B(p)=\operatorname{adj}T(p)
\]
と置き、$H=h(q(a))$ と書く。
ここで各テンソルの値域を省略してはならない。
\[
\begin{split}
 g&\in E^*\otimes L,\qquad
 T\in\operatorname{Sym}^2E^*\otimes L,\\
 h&\in(\det E^*)^2\otimes L^5=L,\\
 B&\in\operatorname{Sym}^2E\otimes(\det E^*)^2\otimes L^4
      =\operatorname{Sym}^2E.
\end{split}
\]
最後の二つの同一視には $\det E^*=\Delta^6$ を用いた。
特に $B$ は自然な写像 $E^*\to E$ であり、ここで選んだ体積に依存する
余分な determinant factor は残らない。

三次式のEuler恒等式 $Tp=2g$ と $BT=h\,\operatorname{id}$ から
\begin{equation}\label{eq:v16-adjugate-euler}
 B g=\frac h2p
\end{equation}
を得る。これを $p_k$ で微分すると、座標表示では
\begin{equation}\label{eq:v16-differentiated-adjugate}
 p_j\partial_k h
   =h\delta_{jk}+2\sum_i(\partial_k B_{ji})g_i
\end{equation}
となる。これは $\operatorname{End}(E)\otimes L$ 値のテンソル恒等式である。

\begin{proposition}[Hessianの一次作用の全次数簡約]
\label{prop:v16-hessian-derivation-reduction}
全ての $r\ge1$ について
\begin{equation}\label{eq:v16-all-hessian-action-mod-units}
 dH(\mathcal G_r)\subset J_{r+9}+H\mathcal U_{r-1}.
\end{equation}
より正確には、$w\in\mathcal D_{r+1}$ に対し
\[
 dh_{q(a)}(w)=H\operatorname{tr}M+dF_{q(a)}(v),\qquad
 v\in\mathcal D_{r+5},
\]
と書ける。最後の項は $J_{r+9}$ に入る。
\end{proposition}

\begin{proof}
前補題の $w=Mq$ を用いる。
式\eqref{eq:v16-differentiated-adjugate}に $M_{kj}$ を掛けて $k,j$ を縮約すれば
\[
 dh_q(Mq)=H\operatorname{tr}M+\sum_i g_i(q)v_i,
 \qquad
 v_i=2\sum_{k,j}M_{kj}(F,a)(\partial_kB_{ji})(F,q(a)).
\]
ここで微分 $\partial_k$ は $B$ にだけ作用し、$M$ を微分しない。
$dB\in E^*\otimes\operatorname{Sym}^2E$ と $M\in E\otimes E^*$ の縮約なので、
$v$ は $E$ 値の普遍同変多項式である。
$B$ は $F$ と $p$ にそれぞれ4次だから、$v$ の二重次数は
\[
 (\deg_Fv,\deg_av)=(r+3,r+5).
\]
従って $v\in\mathcal D_{r+5}$ である。
命題\ref{prop:v12-uniform-pfaffian-surjectivity}の全次数全射
\[
 dq:\mathcal G_{r+4}\twoheadrightarrow\mathcal D_{r+5}
\]
を用いて $v=dq(Z)$ と持ち上げれば $dF_q(v)=dP(Z)\in J_{r+9}$ となる。
$w=dq(Y)$、$Y\in\mathcal G_r$ とすれば所望の包含を得る。
\end{proof}

\begin{proposition}[錐上で消える単元のHessian倍]
\label{prop:v16-plucker-unit-annihilation}
$u\in\mathcal U_s$ が $\Sigma$ 上で消えるなら
\begin{equation}\label{eq:v16-ideal-unit-annihilation}
 Hu\in J_{s+10}.
\end{equation}
従って $H\mathcal U_s$ の通常接触商における像は、
$\mathcal U_s$ のPl\"ucker錐への制限だけに依存する。
\end{proposition}

\begin{proof}
$s=0,1$ の場合はイデアルの最小次数が二であるため $u=0$ である。
$s\ge2$ では、補題\ref{lem:v16-plucker-prime-ideal}と完全可約性から
\[
 u=\langle b,q\rangle,\qquad b(F,a)\in E^*,\qquad
 (\deg_Fb,\deg_ab)=(s,s-2)
\]
という普遍同変な表示を選べる。
式\eqref{eq:v16-adjugate-euler}と $B$ の対称性により
\[
 Hu=2\langle b,Bg\rangle=\langle g,2Bb\rangle.
\]
$2Bb$ は $E$ 値で、その二重次数は $(s+4,s+6)$ なので
$\mathcal D_{s+6}$ に属する。
全射 $dq:\mathcal G_{s+5}\to\mathcal D_{s+6}$ で持ち上げれば結論を得る。
\end{proof}

\begin{remark}[持ち上げの代数性と有理性]
ここでの二つの持ち上げは、有限次数の普遍係数空間の線形写像に対するものである。
$F$ の判別式、$P$ の値、$H$ の値、行列 $T$ の行列式による除算はない。
adjugate は最初から多項式である。
入力が有理係数なら、同変性も含めた線形条件は有理数体上の条件であり、
複素数体上での可解性から有理数体上の持ち上げを選べる。
従ってこれらは固定した $F$ ごとの存在主張ではなく、$F$ に多項式的な普遍恒等式を与える。
\end{remark}

\subsection{単元制限環と三元三次不変式}

$\mathcal U=\bigoplus_{s\ge0}\mathcal U_s$ とし、その斉次イデアルを
\[
 \mathcal N_s=\{u\in\mathcal U_s:u|_\Sigma=0\},\qquad
 \mathcal N=\bigoplus_{s\ge0}\mathcal N_s
\]
と定める。

\begin{theorem}[単元のPl\"ucker制限環]
\label{thm:v16-unit-restriction-ring}
階数二の点 $a_0=e_0\wedge e_1\otimes\Delta^{-1}$ での制限により、次数付き環の同型
\begin{equation}\label{eq:v16-unit-restriction-ring}
 \mathcal U/\mathcal N
 \simeq\mathbb C[\operatorname{Sym}^3Q_3]^{\mathrm{SL}(Q_3)}
 =\mathbb C[S_4,T_6]
 =\mathbb C[\overline U_4,\overline J_6]
\end{equation}
を得る。従って $\mathcal U_s/\mathcal N_s$ の基底は
\[
 \overline U_4^{\,a}\overline J_6^{\,b},\qquad
 a,b\ge0,\quad 4a+6b=s
\]
である。
\end{theorem}

\begin{proof}
補題\ref{lem:v16-plucker-vanishing-D}と同じ一パラメータ部分群は $a_0$ を固定する。
スカラー値の $u(F,a_0)$ の重みは零であり、$F$ の各成分の重みは非負なので、
この多項式は $F$ の $\operatorname{Sym}^3Q_3$ 成分だけに依存する。
$f\in\operatorname{Sym}^3Q_3$ に対し $\overline u(f)=u(f,a_0)$ と置く。
$h\in\mathrm{GL}(Q_3)$ を $V$ 上で恒等写像と直和して作用させると
\[
 h a_0=(\det h)^{-1}a_0,
 \qquad \overline u(hf)=(\det h)^s\overline u(f).
\]
特に $\overline u$ は $\mathrm{SL}(Q_3)$ 不変式であり、その通常の $f$ 次数は $s$ である。
逆に、三次式の係数に $s$ 次の $\mathrm{SL}_3$ 不変式の $\mathrm{GL}_3$ 重みは
$\det^s$ なので、次数とdeterminant twistは一致する。
非零の階数二の点が一軌道であることから、制限写像の核は正確に $\mathcal N$ である。

三元三次式の全係数空間の $\mathrm{SL}_3$ 不変式環は、次数4、6の二つの生成元の
多項式環である。ここではFisherの一次資料\cite[\S3.3,\S4,
Lemma~4.3, Theorem~4.4(i)]{Fisher2008}を、$n=3$、標数零で適用する。
この定理は判別式を可逆化した環ではなく、特異三次式も含む全係数空間の環についてのものである。

命題\ref{prop:v14-rank-two-unit-restriction}によれば、
$f_t=x_2^3+x_3^3+x_4^3+6t x_2x_3x_4$ 上で
\[
 \overline U_4(f_t)=t^4-t,\qquad
 \overline J_6(f_t)=6-120t^3-48t^6.
\]
両者はそれぞれ次数4、6の非零不変式である。
その次数の不変式空間はそれぞれ一次元なので、$S_4,T_6$ の非零定数倍を与える。
従って制限写像の像は不変式環の生成元を両方含み、全射である。
これで環の同型が従う。
既存の制限計算では、$(s,t)\mapsto(s^4\overline U_4(f_t),s^6\overline J_6(f_t))$
の $(1,0)$ におけるJacobian行列式が36であり、代数的独立性も直接検証されている。
\end{proof}

\begin{theorem}[全ての正次数単元倍の消去]
\label{thm:v16-positive-unit-hessian-annihilation}
全ての整数 $s>0$ について
\begin{equation}\label{eq:v16-all-positive-unit-hessian}
 H\mathcal U_s\subset J_{s+10}.
\end{equation}
\end{theorem}

\begin{proof}
低次の入力は、既存の次の二つの所属だけである。
\[
 HU_4\in J_{14},\qquad HJ_6\in J_{16}.
\]
第一の所属は、式\eqref{eq:v11-rho-fourteen-values}と
命題\ref{prop:v11-contact-rank-fourteenth}から
小節\ref{subsec:v11-fourteenth-complement}で証明されている。
第二の所属は、$J_6\in\mathcal U_6$ と
命題\ref{prop:v13-hessian-unit-membership}から従う。

任意の $u\in\mathcal U_s$ に対し、前定理によって
\[
 u-\sum_{4a+6b=s}c_{ab}U_4^aJ_6^b\in\mathcal N_s
\]
と書ける。左辺の $H$ 倍は命題\ref{prop:v16-plucker-unit-annihilation}により
$J_{s+10}$ に入る。
$s>0$ なので各単項式では $a>0$ または $b>0$ である。
前者なら $HU_4$、後者なら $HJ_6$ の既知の持ち上げに残る単元を掛ければよい。
実際、$v\in\mathcal U_t$、$Y\in\mathcal G_r$ について
\[
 vY\in\mathcal G_{r+t},\qquad v\,dP(Y)=dP(vY)
\]
なので、通常接触像は単元共変量環の加群である。
従って全ての項が $J_{s+10}$ に入り、結論を得る。
\end{proof}

\begin{corollary}[全次数のHessian一次作用と22次への帰結]
\label{cor:v16-all-hessian-first-actions}
全ての $r\ge2$ について
\begin{equation}\label{eq:v16-all-hessian-first-actions}
 dH(\mathcal G_r)\subset J_{r+9}.
\end{equation}
特に
\[
 H\mathcal U_{12}\subset J_{22},\qquad
 dH(\mathcal G_{13})\subset J_{22},\qquad
 dH(\mathcal G_8)\subset J_{17}.
\]
$r=1$ はこの消去の主張に含めない。
実際、$\mathcal G_1=\mathbb C a$ であり $dH(a)=10H$ である。
\end{corollary}

\begin{proof}
命題\ref{prop:v16-hessian-derivation-reduction}と
定理\ref{thm:v16-positive-unit-hessian-annihilation}を、$s=r-1>0$ として適用する。
22次については制限環の12次部分が $\overline U_4^3,\overline J_6^2$ で張られるため、
全次数の記述を使わずこの二方向と錐イデアル部分に分けても同じ結論を得る。
また $r=8$ の場合は、既存の $\mathcal U_7=0$ と
式\eqref{eq:v16-all-hessian-action-mod-units}から直接にも従う。
\end{proof}

\subsection{下位jetが二項正規形である場合の有効像}

\begin{theorem}[二項正規形の後の有効像は全次数で通常像に等しい]
\label{thm:v16-all-effective-images}
$N>10$、$c\in\mathbb C$ とする。
$N-1$ 次までのjetが $P+cH$ である方程式について、
その下位jetを保ち、線形部分が恒等写像、単元の定数項が1である普遍接触変換の
$N$ 次有効像は
\begin{equation}\label{eq:v16-all-effective-images}
 I_N=J_N
\end{equation}
である。偶数 $N$ に対しては、奇座標変更と偶単元だけに制限しても同じ像となる。
\end{theorem}

\begin{proof}
接触対数補題\ref{lem:v14-contact-logarithm}により、下位jetを保つ変換の
最初の差は、対数作用素 $\sum_{k>0}Q_k$ の一次作用で計算できる。
ここで $Q_k=Y_{k+1}+b_k\in\mathcal G_{k+1}\oplus\mathcal U_k$ である。
$N$ 次の差は
\[
 Q_{N-6}P+cQ_{N-10}H
\]
となる。第一項はEuler場も含めて $J_N$ に入る。
第二項は $N-10>0$ なので、前定理と前系により
\[
 dH(\mathcal G_{N-9})+H\mathcal U_{N-10}\subset J_N
\]
に属する。従って $I_N\subset J_N$ である。
逆に $Y\in\mathcal G_{N-5}$ による座標変更 $a\mapsto a+Y$ は
$N-1$ 次までを保ち、$N$ 次に $dP(Y)$ を与えるため、$J_N\subset I_N$ である。
$N$ が偶数なら $N-5$ は奇数なので、この逆包含は奇座標変更だけで実現する。
\end{proof}

\begin{remark}[適用範囲]
本定理の仮定は、対象の下位jetが既に $P+cH$ であることである。
消去されていない別の低次項がある場合には、その項への接触作用を追加する必要がある。
また $I_N=J_N$ は $\mathcal C_N/J_N$ の次元も、実際の $R_N$ の類も決定しない。
とりわけ、この構造結果だけから実際の22次係数の零性を結論しない。
\end{remark}

\subsection{有限検算と証明への依存}

\nolinkurl{v16_structure_certificate.py} はPython標準ライブラリだけを使う。
五つのPl\"ucker関係の十組の $S$ 多項式の正確な除法、
二重次数と中心重み、既存の階数二単元制限のJacobian36を検算する。
さらに三つの明示的な整数係数 $F$ について、$p$ を記号変数のまま保持し、
式\eqref{eq:v16-adjugate-euler}の五成分と
式\eqref{eq:v16-differentiated-adjugate}の25成分を全て正確に検算する。
最後の三例は添字・係数の監査であり、有限個の代入から普遍恒等式を推論するものではない。
普遍恒等式の証明は $Tp=2g$ とadjugateの恒等式を微分する上の議論である。
三元三次不変式環については一次資料の定理を用い、
既存の接触階数やTate計算を今回独立に再認証したとは主張しない。

% END theory/v16_plucker_hessian_structure_ja.tex


% BEGIN theory/v16_normal_slice_injection_ja.tex
% 第16版：rank-two normal slice による通常接触商への全次数単射。
% この節の単射は Fisher の不変式環生成定理には依存しない。
% v14 の rank-two U4,J6 独立性、v15 の既知328次元部分空間を継承する。

\section{階数二の横断三次項と通常接触商の独立な方向}
\label{sec:v16-normal-slice-injection}

前節のPl\"ucker錐上の消失を使うと、通常接触像を消す新しい検出方法が得られる。
それは階数二の点の横断三次項を、三元三次式のJacobian空間を法として取る方法である。
この方法は、有限体の零やグラフ探索の停滞による上界を用いない。
特に第15版で残っていた20次通常階数の二択を解消する。

\subsection{横断切片と三次Jacobian空間}

体積形式を一つ固定し、$p=(p_2,p_3,p_4)$、$\varepsilon$ を独立変数として
\begin{equation}\label{eq:v16-normal-slice}
 \begin{split}
 a_\varepsilon(p)&=e_0\wedge e_1+
 \varepsilon\bigl(p_4e_2\wedge e_3-p_3e_2\wedge e_4+p_2e_3\wedge e_4\bigr),\\
 F_f&=x_0^3+x_1^3+f(x_2,x_3,x_4)
 \end{split}
\end{equation}
と置く。$a$ の $\Delta^{-1}$ 因子はこの体積形式で表示している。
Pfaffianは正確に
\begin{equation}\label{eq:v16-normal-pfaffian}
 q(a_\varepsilon(p))=(0,0,\varepsilon p_2,\varepsilon p_3,\varepsilon p_4)
\end{equation}
となる。$f$ は三元三次式の全係数空間を動く。
\[
 \operatorname{Jac}(f)_3
 =\operatorname{Span}_{\mathbb C}
 \{p_j\partial_i f: i,j\in\{2,3,4\}\}
 \subset\mathbb C[p_2,p_3,p_4]_3
\]
と書く。

\begin{lemma}[通常接触像の横断三次項]
\label{lem:v16-contact-normal-cubic-jet}
$N\ge6$、$Y\in\mathcal G_{N-5}$ とする。
この切片における $dP(Y)$ の $\varepsilon^3$ 係数は
$\operatorname{Jac}(f)_3$ に入る。
\end{lemma}

\begin{proof}
$w=dq(Y)\in\mathcal D_{N-4}$ と置く。
補題\ref{lem:v16-plucker-vanishing-D}により、
$a_0=e_0\wedge e_1$ で $w(F_f,a_0)=0$ である。
従って、$p$ に線形なベクトル $w_1(f,p)$ を用いて
\[
 w(F_f,a_\varepsilon(p))=\varepsilon w_1(f,p)+O(\varepsilon^2)
\]
と展開できる。
一方、$F_f$ の最初の二つの偏微分は $3q_0^2,3q_1^2$ なので、
切片上では正確に零である。残る三成分は
\[
 \partial_iF_f(q(a_\varepsilon))=\varepsilon^2\partial_i f(p)
 \qquad(i=2,3,4)
\]
となる。よって
\[
 [\varepsilon^3]\,dP(Y)(F_f,a_\varepsilon(p))
   =\sum_{i=2}^4\partial_i f(p)\,(w_1)_i(f,p)
   \in\operatorname{Jac}(f)_3.
\]
場 $Y$ の接方向の成分を零と仮定する必要はない。
最初の二成分の寄与が消えるのは、分離した $F_f$ と切片の明示式による。
\end{proof}

\begin{lemma}[$\Theta_8$ の切片上の正確な値]
\label{lem:v16-theta-normal-hessian}
$B_f(p)=\tfrac16\operatorname{Hess}_f(p)$ と置くと
\begin{equation}\label{eq:v16-exact-theta-slice}
 \Theta_8(F_f,a_\varepsilon(p))
   =\varepsilon^3\det B_f(p)
   =\frac{\varepsilon^3}{216}\det\operatorname{Hess}_f(p).
\end{equation}
これは最初の係数だけでなく、切片上での正確な多項式恒等式である。
また $u\in\mathcal U_s$ について
\[
 [\varepsilon^0]\,u(F_f,a_\varepsilon(p))=\overline u(f)
\]
であり、右辺は前節の階数二制限である。
\end{lemma}

\begin{proof}
既存の式\eqref{eq:boundary-theta-eight}を用いる。
切片上の $\operatorname{Hess}_{F_f}(q)$ は、最初の二行・二列が零であり、
残る三行三列のブロックは $\varepsilon\operatorname{Hess}_f(p)$ である。
従って3行3列minorの和に残るのは
$I=J=\{2,3,4\}$ の一項だけである。
その係数は $c_{234}^2=a_{01}^2=1$ なので、表示した恒等式となる。
単元の定数項は $u(F_f,a_0)$ であり、
前節と同じ一パラメータ部分群の重みにより、これは $F_f$ の三元部分 $f$ だけに依存する。
\end{proof}

三元三次式 $f$ について、十個の三次式
\begin{equation}\label{eq:v16-ten-ternary-cubics}
 \{p_j\partial_i f: i,j=2,3,4\},\qquad \det B_f
\end{equation}
の係数行列を $M(f)$ とする。
Fermat式 $f=p_2^3+p_3^3+p_4^3$ では、行を
九個の単項式 $p_i^2p_j$（$i,j$ の順）と $p_2p_3p_4$ の順に並べると
\begin{equation}\label{eq:v16-fermat-normal-minor}
 M(f)=\operatorname{diag}(3,3,3,3,3,3,3,3,3,1),
 \qquad\det M(f)=3^9=19683\ne0.
\end{equation}
従って三元三次式の非空なZariski開集合で
\[
 \mathbb C[p_2,p_3,p_4]_3
   =\operatorname{Jac}(f)_3\oplus\mathbb C\det B_f
\]
となる。この非零性には、Jacobian環の完全交叉性やsocleについての外部定理は不要である。

独立検算 \nolinkurl{v16_normal_slice_certificate.py} は、
上の十行十列行列とPfaffianの切片式を正確な整数・有理数で再構成する。
さらに $f_t=p_2^3+p_3^3+p_4^3+6t p_2p_3p_4$ について
\[
 \det M(f_t)=3^9(1+8t^3)^3
\]
および
\[
 \det B_{f_t}
 =(1+8t^3)p_2p_3p_4-\frac{t^2}{3}\sum_{i=2}^4p_i\partial_i f_t
\]
を検証する。前者は次数上界12と13個の整数値、後者は係数ごとの多項式恒等式による。
これらの追加のHesse計算を用いなくても、Fermatの一つの非零minorと開性で以下の証明は閉じる。

\subsection{通常接触商への全次数単射}

\begin{theorem}[$\Theta_8$ と二つの単元が与える自由部分加群]
\label{thm:v16-hessian-socle-injection}
$\deg U_4=4$、$\deg J_6=6$、$\deg\Theta_8=8$ とすると、
\begin{equation}\label{eq:v16-theta-module-injection}
 \mathbb C[U_4,J_6]\,[\Theta_8]
 \hookrightarrow
 \bigoplus_N\mathcal C_N/J_N
\end{equation}
は次数付き $\mathbb C[U_4,J_6]$ 加群の単射である。
すなわち、各 $s\ge0$ について
\[
 \{[\Theta_8U_4^aJ_6^b]:4a+6b=s\}
 \quad\text{は }\mathcal C_{s+8}/J_{s+8}\text{ で独立である。}
\]
\end{theorem}

\begin{proof}
重み $s$ の多項式 $A(U_4,J_6)$ が
$\Theta_8 A(U_4,J_6)\in J_{s+8}$ を満たすと仮定する。
前の二補題により、全ての三元三次式 $f$ に対して
\[
 A(\overline U_4(f),\overline J_6(f))\det B_f
      \in\operatorname{Jac}(f)_3
\]
となる。式\eqref{eq:v16-fermat-normal-minor}が定める非空な開集合では
$[\det B_f]\ne0$ なので
\[
 A(\overline U_4(f),\overline J_6(f))=0
\]
である。これは三元三次式の係数の多項式だから、全係数空間で恒等的に零となる。
命題\ref{prop:v14-rank-two-unit-restriction}の代数的独立性から $A=0$ が従う。
また通常接触像は単元共変量による積で保たれるため、この単射は加群の単射である。

この証明は、$\overline U_4,\overline J_6$ の既存の独立性だけを使う。
三元三次不変式環の全生成定理\cite{Fisher2008}は本定理の依存には入らない。
\end{proof}

\begin{corollary}[20次通常接触階数の確定]
\label{cor:v16-exact-twentieth-contact-rank}
第15版の記号を保つと
\begin{equation}\label{eq:v16-exact-twentieth-rank}
 J_{20}=S_{20},\qquad\dim J_{20}=328,
 \qquad\dim\mathcal C_{20}/J_{20}=2.
\end{equation}
通常商の基底は
\[
 [\Theta_8U_4^3],\qquad[\Theta_8J_6^2]
\]
であり、
\[
 \dim\ker(dP:\mathcal G_{15}\to\mathcal C_{20})=916,
 \qquad
 \dim\ker(dF:\mathcal D_{16}\to\mathcal C_{20})=360.
\]
\end{corollary}

\begin{proof}
命題\ref{prop:v15-twentieth-contact-subspace}の既存の整数証明書により
$S_{20}\subset J_{20}$、$\dim S_{20}=328$、$\dim\mathcal C_{20}=330$ である。
前定理を $s=12$ に適用すると、二つの候補は $\mathcal C_{20}/J_{20}$ でも独立なので
$\dim J_{20}\le328$ となる。従って $J_{20}=S_{20}$ である。
既存の $\dim\mathcal G_{15}=1244$ から最初の核の次元を得る。
$dq:\mathcal G_{15}\twoheadrightarrow\mathcal D_{16}$ が全射であり
$dP=dF\circ dq$ なので、$dF$ の階数も328である。
$\dim\mathcal D_{16}=688$ から最後の次元を得る。
\end{proof}

\begin{remark}[第15版の検出子とKoszul核の読み替え]
第15版の二つの三点汎関数は $S_{20}$ を正確に消していた。
今回の $S_{20}=J_{20}$ により、それらは通常接触商上の双対基底でもあることが新たに確定した。
これは第15版の時点で仮定していた性質ではない。
第15版の原データ、三点値、実際の20次消去の証明を変更する必要はない。

また既存の359個の独立なKoszul核ベクトルは、確定した360次元核の真部分空間である。
残る一方向が別のKoszul像で生成されるのか、Koszul像の外にあるのかは、
この論証では決まらない。従ってKoszul完全性や、全Koszul像の階数は引き続き主張しない。
\end{remark}

\begin{corollary}[通常接触商の全次数Hilbert級数の下界]
\label{cor:v16-contact-hilbert-lower-bound}
$\mathbb C[U_4,J_6]$ 加群として
\[
 \mathbb C[U_4,J_6]\,[\Theta_8]\oplus\mathbb C\cdot[H]
 \hookrightarrow\bigoplus_N\mathcal C_N/J_N
\]
であり、第二項では正次数イデアル $(U_4,J_6)$ が零作用する。
従ってHilbert級数には係数ごとの下界
\begin{equation}\label{eq:v16-contact-hilbert-lower-bound}
 \sum_N\dim(\mathcal C_N/J_N)t^N
 \ \ge_{\mathrm{coeff}}\
 \frac{t^8}{(1-t^4)(1-t^6)}+t^{10}
\end{equation}
が成り立つ。
\end{corollary}

\begin{proof}
$[H]\ne0$ は定理\ref{thm:boundary-eleven-jet}の十次通常商の計算による。
定理\ref{thm:v16-positive-unit-hessian-annihilation}から、正次数単元は $[H]$ を消す。
$\Theta_8$ 部分加群の次数は $8+4a+6b$ であり、10次は含まないので直和となる。
残りは前定理による。
\end{proof}

\begin{remark}[下界と全商の同定を区別する]
式\eqref{eq:v16-contact-hilbert-lower-bound}を全商のHilbert級数との等式とは主張しない。
横断三次項の検出で見えない追加の接触類が、高次に存在する可能性は残る。
例えば22次にこの方法が与えるのは $[\Theta_8U_4^2J_6]$ の非零性と一次元の下界であり、
22次通常商の全次元も実際のTate係数の類も、これだけでは確定しない。
\end{remark}

% END theory/v16_normal_slice_injection_ja.tex


% BEGIN characters/v16_character_notes_ja.tex
\section{二十二次の普遍共変量とKoszul各項}
\label{sec:v16-characters}

二十二次では、候補となる単元の積だけから通常接触商の次元を決めない。
本節では普遍空間とKoszul源の次元を先に確定し、写像の階数とは区別する。
計算は整数の重み指標、Weyl交代和、Pieri則による。

\subsection{次数とdeterminant twist}

$\dim V=5$、$\Delta=\det V$、$A=\Lambda^2V$、
$W=A\otimes\Delta^{-1}$ とする。
$F\in\operatorname{Sym}^3V$、$a\in W$、
$q(a)=\tfrac12a\wedge a\in V^*\otimes\Delta^{-1}$、
$P=F(q(a))\in\Delta^{-3}$ という既存の規約を保つ。
中心元 $t\operatorname{id}_V$ の重みから、$a$ 次数 $n$ の境界項は
$F$ 次数 $n-5$、$a$ 次数 $r$ の座標補正は $F$ 次数 $r-1$、
$a$ 次数 $s$ の単元補正は $F$ 次数 $s$ を持つ。
従って、trace線を外した普遍空間は
\begin{align*}
 C_n&=\operatorname{Hom}_{\mathrm{GL}(V)}
 \bigl(\operatorname{Sym}^{n-5}(\operatorname{Sym}^3V)
       \otimes\operatorname{Sym}^{n}A,\Delta^{n-3}\bigr),\\
 G_r&=\operatorname{Hom}_{\mathrm{GL}(V)}
 \bigl(\operatorname{Sym}^{r-1}(\operatorname{Sym}^3V)
       \otimes\operatorname{Sym}^{r}A,A\otimes\Delta^{r-1}\bigr),\\
 D_{r+1}&=\operatorname{Hom}_{\mathrm{GL}(V)}
 \bigl(\operatorname{Sym}^{r-1}(\operatorname{Sym}^3V)
       \otimes\operatorname{Sym}^{r+1}A,\Lambda^4V\otimes\Delta^{r-1}\bigr),\\
 U_s&=\operatorname{Hom}_{\mathrm{GL}(V)}
 \bigl(\operatorname{Sym}^{s}(\operatorname{Sym}^3V)
       \otimes\operatorname{Sym}^{s}A,\Delta^s\bigr).
\end{align*}
二十二次のKoszul第 $k$ 項は
\[
 K_{k,22}=\operatorname{Hom}_{\mathrm{GL}(V)}
 \left(\operatorname{Sym}^{17-k}(\operatorname{Sym}^3V)
 \otimes\operatorname{Sym}^{22-4k}A,
 \Lambda^kV^*\otimes\Delta^{19-2k}\right),\qquad 1\leq k\leq5.
\]
$K_{1,22}=D_{18}$、$K_{2,22}=B_{22}$ と書く。
第5項にも $a$ 次数2が残るので、前の次数の終端をそのまま切り捨ててはならない。

\begin{proposition}[二十二次の厳密次元]
\label{prop:v16-character-dimensions}
各普遍空間の次数・標的分割・次元は次の通りである。
\[
\begin{array}{c|r|r|c|r}
\text{空間}&\deg_F&\deg_a&\text{標的分割}&\dim\\\hline
C_{22}&17&22&(19,19,19,19,19)&916\\
G_{17}&16&17&(17,17,16,16,16)&3929\\
U_{16}&16&16&(16,16,16,16,16)&486\\
D_{18}&16&18&(17,17,17,17,16)&2140\\
B_{22}&15&14&(15,15,15,14,14)&1706\\
K_{3,22}&14&10&(13,13,12,12,12)&534\\
K_{4,22}&13&6&(11,10,10,10,10)&57\\
K_{5,22}&12&2&(8,8,8,8,8)&2
\end{array}
\]
これらの値は接触写像の階数やKoszul完全性を仮定しない。
\end{proposition}

\begin{proof}
$\operatorname{Sym}^m(\operatorname{Sym}^3V)$ の指標を、35個の三次単項式重み
$\alpha$ に対するEuler積
\[
 \prod_{|\alpha|=3}(1-tx^\alpha)^{-1}
\]
の $t^m$ 係数として整数演算で計算する。
$\rho=(4,3,2,1,0)$ とすると、Schur多重度は
\[
 [S_\lambda V:\operatorname{Sym}^m(\operatorname{Sym}^3V)]
 =\sum_{\sigma\in S_5}\operatorname{sgn}(\sigma)
 [x^{\lambda+\rho-\sigma\rho}]
 \operatorname{ch}\operatorname{Sym}^m(\operatorname{Sym}^3V)
\]
である。外積二乗側の指標も10個の重みから独立に作り、標的分割に対する
テンソル積の交代和を直接評価した。

また
\[
 \operatorname{Sym}^{d}(\Lambda^2V)
 =\bigoplus_{j=0}^{\lfloor d/2\rfloor}S_{(d-j,d-j,j,j,0)}V
\]
の各多重度1とWeyl次元和を照合した。
$C_{22}$ で必要な $F$ 側の分割は
$(19,16-j,16-j,j,j)$、$0\leq j\leq8$ であり、各多重度は
\[
 (3,4,36,78,226,234,254,61,20).
\]
和は916となる。
$U_{16}$ の分割は $(16,16-j,16-j,j,j)$、$0\leq j\leq8$ で、各多重度は
\[
 (2,1,14,33,112,126,148,37,13).
\]
和は486となる。

$G_{17}$ と $D_{18}$ はそれぞれ外積第3冪・第1冪を用いる垂直帯Pieri計算、
$K_{k,22}$ は外積第 $k$ 冪を用いる垂直帯Pieri計算でも計算し、
テンソル積指標からの値と照合した。
Koszul源で非零となるPieri成分数は $k=2,3,4,5$ の順に25、20、8、1である。
第5項の非零成分は $F$ 側分割 $(8,8,8,6,6)$ で、多重度2である。

全ての $F$ 次数 $1\leq m\leq17$ と $a$ 次数 $1\leq d\leq22$ について、
Euler積とは異なるNewton--Adams再帰
\[
 m h_m=\sum_{k=1}^m p_kh_{m-k}
\]
により全重み係数を再計算し、辞書全体が一致することを確認した。
必要な $F$ 次数12から17までの全Schur分解についても、
多重度とWeyl次元の積の和が $\binom{m+34}{m}$ に一致する。
これらは全て有限の整数恒等式であり、浮動小数点階数を用いない。
\end{proof}

\subsection{Euler交代和と制限計算の範囲}

二十二次のKoszul各項の交代和は
\[
 916-2140+1706-534+57-2=3.
\]
従って $H_k$ を当該有限次数成分のKoszulホモロジーとすると
\[
 \dim H_0-\dim H_1+\dim H_2-\dim H_3+\dim H_4-\dim H_5=3.
\]
この等式だけから $\dim(C_{22}/J_{22})=3$、あるいは1と結論しない。
全次数Pfaffian微分全射性を用いる場合には
\[
 \dim\ker(\delta q:G_{17}\longrightarrow D_{18})=3929-2140=1789
\]
を得るが、$dF:D_{18}\to C_{22}$ の核の次元は別に求める必要がある。

\begin{proposition}[Pl\"ucker錐への制限の標的次元]
\label{prop:v16-plucker-dimensions}
分解可能な $a$ の錐への制限では、$\operatorname{Sym}^dA$ をCartan成分
$S_{(d,d)}V$ に置き換える。
この置換後のHom空間の次元は次の通りである。
\[
\begin{array}{c|rrrrrrrr}
 &C_{22}&G_{17}&U_{16}&D_{18}&B_{22}&K_{3,22}&K_{4,22}&K_{5,22}\\\hline
 \dim&0&2&2&0&0&12&17&2
\end{array}
\]
また $U_{14}$ の制限標的は1次元である。
\end{proposition}

\begin{proof}
前のLittlewood分解の $j=0$ のみを残して同じPieri計算を行う。
$U_{16}$ では $F$ 側分割 $(16,16,16,0,0)$ の多重度2、
$U_{14}$ では $(14,14,14,0,0)$ の多重度1を得る。
$D_{18}$ では第一行18が比較する矩形の幅17を超えるため全成分が消える。
各成分と交代和は保存証明書の \texttt{Plucker} 欄に収録した。
\end{proof}

ここで計算したのは制限先の普遍Hom空間である。
特定の候補がその空間を尽くすことや、通常接触商全体を尽くすことを
この表から自動的に結論してはならない。
特に $U_{14}$ の錐上成分が1次元であっても、それだけで
$C_{22}/J_{22}$ が $[\Theta_8U_4^2J_6]$ の1次元空間になるとはいえない。

Fermat固定群による従来の成分別制限計算からも、次の上界を得た。
\[
\begin{array}{c|r|r|r}
 &\dim&\text{対角値の階数上界}&\text{対角値と一階法jetの階数上界}\\\hline
C_{22}&916&284&803\\
G_{17}&3929&972&3244\\
U_{16}&486&106&371\\
D_{18}&2140&576&1834
\end{array}
\]
これは固定群指標から得る上界であり、実際の評価写像の階数の等号を主張しない。
対角値や一階jetだけを調べて普遍空間の全てを見たとみなすことはできない。

\subsection{再現用証明書}

\texttt{characters/v16\_character\_certificate.py} が全重みのEuler積・Newton再帰、
全Schur次元和、テンソル積の交代和とPieri照合を新規に実行する。
\texttt{v16\_koszul\_pieri\_check.py} は独自の垂直帯列挙を使い、
保存されたSchur分解から全Koszul源を再計算する。
\texttt{v16\_character\_audit.py} は保存証明書の次数、twist、交代和、
Weyl次元、Pieri、制限上界を短時間で監査する。
保存証明書の監査と全重みの新規再計算は区別する。

% END characters/v16_character_notes_ja.tex


% BEGIN rank/rank22_section_ja.tex
\section{二十二次の通常接触商とKoszul非完全性}
\label{sec:v16-contact-rank}

本節では、指標による源の次元とは独立に通常接触像の階数を決める。
下界には普遍場の明示的な評価小行列を、上界には
定理\ref{thm:v16-hessian-socle-injection}の横断三次jetによる非所属を使う。
KoszulのEuler交代和から写像の階数を推測する必要はない。

\begin{proposition}[二十二次の通常接触商]
\label{prop:v16-contact-rank-twenty-second}
通常接触像を $J_{22}=dP(G_{17})=dF(D_{18})$ とすると
\[
 \dim J_{22}=915,\qquad
 C_{22}/J_{22}=\mathbf C\cdot[\Theta_8U_4^2J_6].
\]
従って
\[
 \dim\ker(dP:G_{17}\to C_{22})=3014,\qquad
 \dim\ker(dF:D_{18}\to C_{22})=1225.
\]
さらに429個の明示的な普遍場 $\Gamma_{17,\alpha}$ により
\[
 J_{22}=P U_{16}\oplus
 \operatorname{span}\{dP(\Gamma_{17,\alpha}):1\leq\alpha\leq429\}
\]
と書ける。
\end{proposition}

\begin{proof}
17対17の三価二部グラフを考え、左側の最後の頂点を開口とする。
残り16個の左頂点へ対称三階テンソル $F$、17個の右頂点へ
$\tau_{ijk}=\epsilon_{ijklm}a_{lm}$ を置いて縮約し、自由三脚 $R_{ijk}$ を得る。
\[
 \Gamma^{lm}=\frac16\sum_{i,j,k}\epsilon_{lmijk}R_{ijk}
\]
は $F$ 次数16、$a$ 次数17の普遍場で、$G_{17}$ に属する。
開口へ $T_{ijk}=\epsilon_{ijklm}(\partial P/\partial a_{lm})$ を戻した閉縮約を
6で割ると $dP(\Gamma)$ になる。

採用した429個のグラフのうち286個は連結グラフの探索から、143個は
既存の低次場と単元の積から得た。
後者に用いた候補は、非単元 $J_{18}$ の60方向に $U_4$ を掛けた60列、
$J_{16}$ の21方向と $U_6$ の2方向から作る42列、
$J_{14}$ の7方向と $U_8$ の6方向から作る42列、
$J_{12}$ の2方向と $U_{10}$ の18方向から作る36列の計180列である。
そのうち143列を採用した。
単元 $u$ と場 $\Gamma$ に対して $dP(u\Gamma)=u\,dP(\Gamma)$ なので、
全て二十二次の通常接触像に属する。
積は単元グラフと場のグラフの非連結和でも表せるため、
最終証明書には各列を単独で再構成できる17対17のグラフを記録した。
四頂点単元グラフは未除算の値 $24U_4$ を使うが、非零定数倍は張る空間を変えない。

$q=(1,1,1,1,1)$、$P=0$ を有理数上で満たす450個の整数入力で評価する。
$F_{444}$ を調整して $P=0$ を課しており、有限体上だけの条件ではない。
429列の評価行列には、法1009で行列式527となる429行の小行列がある。
各グラフは整係数縮約で、分母6は1009で可逆なので、これは特性0での独立性を証明する。
選択した小行列の全 $429^2=184041$ 成分を、別の辺番号付けと整数縮約で再計算し、
二種類のGaussian消去でも非零行列式を照合した。

$P$ は非零多項式であり、普遍係数環は整域なので
$\dim(PU_{16})=\dim U_{16}=486$ である。
Eulerの等式 $dP(au/6)=Pu$ により $PU_{16}\subset J_{22}$。
前の429方向は $P=0$ 制限で独立なので $PU_{16}$ と交わらず、
\[
 \dim J_{22}\geq486+429=915.
\]

一方、定理\ref{thm:v16-hessian-socle-injection}により
$\Theta_8U_4^2J_6\notin J_{22}$。
$\dim C_{22}=916$ なので $\dim J_{22}\leq915$ となり、両界が一致する。
全次数Pfaffian微分全射性により $dP(G_{17})=dF(D_{18})$ であるから、
核の次元は $3929-915=3014$ および $2140-915=1225$ となる。
\end{proof}

\begin{proposition}[補元の独立な有限証明書]
\label{prop:v16-twenty-second-complement}
前の429列へ $W_{22}=\Theta_8U_4^2J_6$ の評価列を加えた行列には、
法1009で行列式591となる430行の小行列がある。
従って、明示的な直和分解
\[
 C_{22}=P U_{16}\oplus
 \operatorname{span}\{dP(\Gamma_{17,\alpha}):1\leq\alpha\leq429\}
 \oplus\mathbf C W_{22}
\]
が得られる。
\end{proposition}

\begin{proof}
$U_4$ は四頂点グラフを24で割った既存の正規化、$J_6$ は六頂点グラフの
既存の正規化を使う。450入力全てで $\Theta_8,U_4,J_6$ と $W_{22}$ を
有理数上で新規に評価し、選択小行列の行列式591を得る。
$486+430=916$ という独立な指標次元との一致から直和が従う。
この有限証明書単独は既知915次元部分空間への補元を示すものだが、
前命題の非所属による上界と合わせることで通常商そのものの基底となる。
\end{proof}

\begin{corollary}[二十二次の正次数Koszulホモロジー]
\label{cor:v16-positive-koszul-homology}
同一の二重次数とdeterminant twistを持つ二十二次Koszul成分のホモロジーを
$H_k$、その次元を $h_k$ とすると
\[
 h_2+h_4=2+h_1+h_3+h_5\geq2.
\]
従ってこの成分は正のホモロジー次数で完全ではなく、$H_2$ または $H_4$ は非零である。
\end{corollary}

\begin{proof}
命題\ref{prop:v16-character-dimensions}の全Koszul源次元からEuler交代和は3。
前命題により $h_0=\dim(C_{22}/J_{22})=1$ なので
\[
 3=1-h_1+h_2-h_3+h_4-h_5
\]
を整理すればよい。
\end{proof}

さらに定理\ref{thm:v16-koszul-high-homology-vanishing}を用いると、
$H_i=0$ ($i>2$) により
\[
 h_2=h_1+2\geq2
\]
へ強化される。これは次節で示す普遍勾配の三成分の正則列性に依存する帰結であり、
Euler交代和のみの帰結と区別する。
対応する第5・第4・第3Koszul微分の階数2、55、479は
系\ref{cor:v16-twentysecond-koszul-defect}にまとめる。
第1・第2ホモロジー群の個別の次元は、なおこの計算だけでは決まらない。
$\ker dF$ が1225次元であっても、それが第2Koszul微分の像全体に一致するとは
主張していない。補助的なKoszulグラフ探索で構成した323独立候補は、
通常階数915の証明にも上の非完全性の証明にも使っていない。
同様に、二十次でEuler交代和と通常商の次元がともに2であることは、
正次数ホモロジーの個別消滅を意味しない。

再現用の主入力は \texttt{rank/contact22\_combined.json}、
\texttt{rank/complement22\_certificate.json} と独立縮約監査である。
\texttt{rank/v16\_rank\_certificate.py} は保存された小行列・入力hash・補元値を監査する。
\texttt{--recompute} を付けると選択小行列の全成分を別の縮約方法で作り直す。
幾何学的な横断jetの非所属定理と全次数Pfaffian微分全射性は、
有限計算とは区別された原稿の証明入力である。

% END rank/rank22_section_ja.tex


% BEGIN theory/v16_koszul_height_ja.tex
% 第16版：普遍勾配イデアルのheightによるKoszul高次homologyの消失。
% 全Koszul完全性は主張しない。H1,H2が残り得ることが本質である。

\section{普遍勾配の余次元とKoszul複体の残る二つのhomology}
\label{sec:v16-universal-koszul-height}

22次では、全ての源次元の交代和と通常商の次元が一致しない。
この差の位置を特定するため、$F$ と $a$ を共に変数とする普遍係数環を用いる。
体積形式と基底を選んで
\[
 R=\mathbb C[F,a],\qquad\dim R=35+10=45,
 \qquad g_i=(\partial_iF)(q(a))\quad(0\le i\le4)
\]
と置く。固定した $F$ のJacobian環に置き換えることはしない。

\begin{proposition}[固定した三つの普遍勾配は正則列]
\label{prop:v16-universal-three-gradient-regular}
イデアル $(g_0,g_1,g_2)$ のheightは3であり、
$g_0,g_1,g_2$ は $R$ 上の正則列である。
また五つ全部の勾配イデアル $(g_0,\ldots,g_4)$ のheightも3である。
\end{proposition}

\begin{proof}
$q\ne0$ を固定すると、線形写像
\[
 \operatorname{Sym}^3V\longrightarrow V,\qquad F\longmapsto dF_q
\]
は全射である。例えば $q=e_0^*$ に基底変換すれば、
対称テンソル係数 $f_{000},f_{001},f_{002},f_{003},f_{004}$ をそれぞれ変えることで
$g_0,\ldots,g_4$ を独立に変えられる。
本稿の $F(p)=\sum f_{ijk}p_ip_jp_k$ の規約では、この部分行列は $3I_5$ である。
座標に依存しない表示として、$q(v)=1$ となる $v\in V$ を選べば、任意の $z\in V$ に対し
\[
 F_z=v^2z-\frac23q(z)v^3
 \qquad\Longrightarrow\qquad d(F_z)_q=z
\]
とも書ける。
従って $q(a)\ne0$ の開集合上では、固定した三成分への射影
$F\mapsto(g_0,g_1,g_2)$ も常に階数3である。

$a$ の開集合 $W\setminus\Sigma$ は10次元なので、
その上の $g_0=g_1=g_2=0$ の零点は階数32のベクトル束の全空間であり、次元は42である。
一方、$q=0$ 上では全ての $g_i$ が恒等的に零である。
補題\ref{lem:v16-plucker-prime-ideal}のチャートから $\dim\Sigma=7$ であるため、
この部分の次元は $35+7=42$ である。
両部分を合わせた零点集合の次元は42であり、$R$ 内のheightは3となる。
五成分を同時に零とする場合、開集合上の次元は $10+(35-5)=40$、
錐上の次元は42なので、やはりheightは3である。

多項式環は正則なCohen--Macaulay環である。
ここで大域的な正則列へ戻すために斉次性を使う。
全ての $F,a$ 変数に次数1を与えると各 $g_i$ は5次斉次であり、
斉次極大イデアル $\mathfrak m=(F,a)$ に属する。
斉次性と上のheight計算により、$R_{\mathfrak m}$ でも
$(g_0,g_1,g_2)$ のheightは3である。
従って局所的な正則列判定から、この固定した順序の三元は
$R_{\mathfrak m}$ 上の正則列となる
\cite[Corollary~3.8]{SwanMKN}。
各乗算写像の核は次数付き加群であり、非零なら非零な斉次元を含む。
そのような元は、非零定数項を持つ多項式に掛けても最低次数の項が消えないため、
$\mathfrak m$ での局所化後も非零である。
局所化後の単射性から各大域的な乗算写像も単射となる。
また商は斉次極大で非零なので、$g_0,g_1,g_2$ は $R$ 上の正則列である。
同じ局所的な可換環論の入力は\cite[Tags 00NQ, 02JN]{StacksCMRegular}でも確認できる。
ここで正則列は三つの固定した成分についての結論であり、
五つ全部の成分を正則列と仮定しているのではない。
\end{proof}

\begin{theorem}[普遍Koszul複体の高いhomologyは消える]
\label{thm:v16-koszul-high-homology-vanishing}
$K_\bullet=K(g_0,g_1,g_2,g_3,g_4;R)$ を普遍勾配のKoszul複体とする。
全ての $i>2$ について
\begin{equation}\label{eq:v16-koszul-high-homology-zero}
 H_i(K_\bullet)=0.
\end{equation}
従って、本稿で使う任意の有限二重次数とdeterminant twistの同変成分でも
$H_3,H_4,H_5$ は零である。
\end{theorem}

\begin{proof}
前命題によって $K(g_0,g_1,g_2;R)$ は
$R/(g_0,g_1,g_2)$ の自由分解である。
\[
 K_\bullet=K(g_0,g_1,g_2;R)\otimes_R K(g_3,g_4;R)
\]
と分ける。後者は自由加群からなる長さ二の複体なので、テンソル後の複体は
\[
 R/(g_0,g_1,g_2)\otimes_R K(g_3,g_4;R)
\]
と擬同型である。これはhomological degree $0,1,2$ にしか項を持たない。
これで全 $R$ 加群としての高次homologyの消失が従う。
この証明のために選んだ三成分の分割が $\mathrm{GL}(V)$ 不変である必要はない。
得られた結論は、もとの同変なKoszul複体のhomologyそのものについての消失である。

twistと二重次数も明示する。勾配は
$V\otimes\Delta^{-2}$ 値で、$F,a$ の次数は $(1,4)$ である。
従って、$k$ 番目の外積の値域を
\[
 \Lambda^k(V^*\otimes\Delta^2)\otimes\Delta^{-3}
   =\Lambda^kV^*\otimes\Delta^{2k-3}
\]
とすれば、勾配による縮約が正しい同変微分となる。
全二重次数を $(n-5,n)$ に固定すると、係数多項式の次数は
$(n-5-k,n-4k)$ であり、$a$ のtrace線を外した同変成分は正確に
\[
 \mathcal K_{k,n}=\operatorname{Hom}_{\mathrm{GL}(V)}
 \left(\operatorname{Sym}^{n-5-k}(\operatorname{Sym}^3V)
 \otimes\operatorname{Sym}^{n-4k}A,
 \Lambda^kV^*\otimes\Delta^{n-3-2k}\right)
\]
となる。負の次数を持つ項は零とする。
各二重次数の空間は有限次元であり、標数零の完全可約性によって
同変成分を取る操作は完全である。
従って全複体の式\eqref{eq:v16-koszul-high-homology-zero}は、
この有限次元複体にもそのまま降りる。
\end{proof}

\begin{corollary}[22次のKoszul defectは第二homologyに現れる]
\label{cor:v16-twentysecond-koszul-defect}
22次の同変Koszul複体の微分を
$\delta_k:\mathcal K_{k,22}\to\mathcal K_{k-1,22}$ とし、
$h_i=\dim H_i(\mathcal K_{\bullet,22})$ と書く。
命題\ref{prop:v16-character-dimensions}と
命題\ref{prop:v16-contact-rank-twenty-second}の結果から
\begin{equation}\label{eq:v16-twentysecond-koszul-defect}
 h_2=h_1+2\ge2,
 \qquad
 \operatorname{rank}\delta_5=2,
 \quad\operatorname{rank}\delta_4=55,
 \quad\operatorname{rank}\delta_3=479
\end{equation}
となる。特に22次のKoszul複体は正次数で完全ではない。
\end{corollary}

\begin{proof}
確定した次元は順に
\[
 \dim(\mathcal K_{0,22},\ldots,\mathcal K_{5,22})
   =(916,2140,1706,534,57,2)
\]
であり、その交代和は3である。
$dq$ の全射性と通常階数915により $h_0=916-915=1$ である。
前定理によって $h_3=h_4=h_5=0$ なので、Euler等式は
$1-h_1+h_2=3$ となる。
また $H_5=0$ から $\delta_5$ は単射であり、その階数は2である。
$H_4=0$ より $\ker\delta_4=\operatorname{im}\delta_5$ なので
$\operatorname{rank}\delta_4=57-2=55$、
$H_3=0$ より同様に $\operatorname{rank}\delta_3=534-55=479$ を得る。
\end{proof}

\begin{remark}[まだ決まらない階数]
この結論は $\delta_2:\mathcal B_{22}\to\mathcal D_{18}$ の全階数を決めない。
$b=\operatorname{rank}\delta_2$ と置けば
\[
 h_1=1225-b,\qquad h_2=1227-b
\]
である。構成した少数のKoszulベクトルや探索の停止を、$b$ の上界や正確な値と取り違えない。
また20次では同じ高次homology消失と今回確定した通常商次元2により
$h_2=h_1$ となるが、その共通値の零性はなお決まらない。
既存の359独立Koszulベクトルと360次元核からは、20次のこの共通値が0または1であることまで従う。
\end{remark}


% END theory/v16_koszul_height_ja.tex


% BEGIN contact/contact_twenty_second_ja.tex
% 第16版：22次の filtered contact 作用と低次正規化の補正。
% 普遍的な補正式・有限縮約と、まだ別証明を要する像所属を区別する。

\section{二十二次の接触作用と正規化補正の有限縮約}
\label{sec:v16-contact-twenty-second}

$P=P_F=F(q(a))$、$H=\det\operatorname{Hess}_F(q(a))$、$c=1/96$ とし、
$J_n=dP(\mathcal G_{n-5})$ と書く。
全ての場、単元、持ち上げは $F$ に多項式的な普遍共変量の空間で選ぶ。
固定した $F$ ごとのJacobian商を用いて持ち上げるのではない。
本節は、前版で未決定であった通常20次階数を仮定しない。

\subsection{二十二次の対数作用}

接触作用素 $T(f)=u\,f\circ\phi$ の対数を
\[
 L=\log T=\sum_{k\ge4}Q_k,\qquad
 Q_k=Y_{k+1}+b_k\in\mathcal G_{k+1}\oplus\mathcal U_k
\]
とする。第14版の補題\ref{lem:v14-contact-logarithm}により、
これは導分と乗算作用素の和であり、最小の次数増分は4である。

\begin{proposition}[二十二次の保存条件と有効作用]
\label{prop:v16-effective-twenty-second-chain}
$P+cH$ の21次までを保存する条件は
\begin{equation}\label{eq:v16-twenty-second-preservation}
 Q_kP+cQ_{k-4}H=0\quad(4\le k\le15),\qquad Q_j=0\ (j<4)
\end{equation}
であり、22次の作用は $Q_{16}P+cQ_{12}H$ である。
特に、線形条件の核
\[
 \mathcal L_{22}=\{(Q_4,Q_8,Q_{12}):
 Q_4P=0,\ Q_8P+cQ_4H=0,\ Q_{12}P+cQ_8H=0\}
\]
を用いれば
\begin{equation}\label{eq:v16-effective-twenty-second-kernel}
 I_{22}=J_{22}
    +c\{Q_{12}H:(Q_4,Q_8,Q_{12})\in\mathcal L_{22}\}.
\end{equation}
全ての普遍接触変換と、奇座標・偶単元に制限した変換は同じ22次像を与える。
さらに系\ref{cor:v16-all-hessian-first-actions}と
定理\ref{thm:v16-positive-unit-hessian-annihilation}を用いると
\begin{equation}\label{eq:v16-effective-twenty-second-ordinary}
                         I_{22}=J_{22}.
\end{equation}
\end{proposition}

\begin{proof}
$T(P+cH)-(P+cH)$ が22次以上であることと、
$L(P+cH)$ が22次以上であることは対数補題により同値である。
この場合、22次では両者が一致する。
従って次数ごとに比較すれば保存条件と作用式を得る。

偶数増分の条件は
\[
\begin{array}{lll}
Q_4P=0,&Q_8P+cQ_4H=0,&Q_{12}P+cQ_8H=0,\\
Q_6P=0,&Q_{10}P+cQ_6H=0,&Q_{14}P+cQ_{10}H=0
\end{array}
\]
という二本の鎖に分かれる。
22次作用に現れるのは第一の鎖だけで、第二の鎖は全て零に選べる。
奇数増分の条件も独立なので、それらも零に選べる。
また $P\mathcal U_{16}\subset J_{22}$ より $Q_{16}P$ の像は $J_{22}$ である。
これで\eqref{eq:v16-effective-twenty-second-kernel}を得る。

一般の $Q_{12}$ に対して
\[
 Q_{12}H=dH(Y_{13})+b_{12}H
       \in dH(\mathcal G_{13})+H\mathcal U_{12}\subset J_{22}
\]
である。逆包含 $J_{22}\subset I_{22}$ は $Q_{16}$ の自由な作用から従う。
ここでは通常20次商の次元も、22次の通常像の階数も用いない。
\end{proof}

\subsection{merged規約からの正確な二十二次台帳}

未正規化Tate係数を $f_n=d_n^{\mathrm T}$ とし、前版と同じ
\[
 X_7=-\frac{15800}{6137}\Phi_7-\frac{689855}{883728}\Psi_7,
 \qquad X_8=\frac{4600}{2189}\operatorname{Alt}K_8
\]
を用いる。初回変更は厳密に
\[
 g=(1+3U_4-V_6)f(a+X_7+X_8)
\]
とし、単元を積 $(1+3U_4)(1-V_6)$ に変更しない。
既存の恒等式は
\[
 f_{10}=cH-3PU_4,\quad
 f_{12}+dP(X_7)=PV_6,\quad f_{13}+dP(X_8)=0
\]
である。従って $g_6=P$、$g_{10}=cH$、$g_{12}=g_{13}=0$ であり、
\[
 A_{14}:=g_{14}=f_{14}+3U_4f_{10},\qquad g_{15}=f_{15},
\]
\[
 A_{16}:=g_{16}=f_{16}+df_{10}(X_7)
       +3U_4(f_{12}+dP(X_7))-V_6f_{10}.
\]
これらの14、15、16次項は、それぞれ $J_{14},J_{15},J_{16}$ に属する。

以下では、Euler場で単元項を吸収した純粋な持ち上げを選び、
\[
 dP(\Xi_9)=A_{14},\qquad dP(\Xi_{10})=f_{15},\qquad
 dP(\Xi_{11})=A_{16}
\]
により、順に $a-\Xi_9$、$a-\Xi_{10}$、$a-\Xi_{11}$ で消去する。
ここで
\[
 g_{18}=f_{18}+df_{12}(X_7)+\tfrac12P''(X_7,X_7)
           +3U_4f_{14}-V_6(f_{12}+dP(X_7))
\]
であり、18次消去の入力は
\begin{equation}\label{eq:v16-eighteenth-normalizer-input}
 A_{18}=g_{18}-c\,dH(\Xi_9).
\end{equation}
第14版の所属から $A_{18}\in J_{18}$ なので、
$dP(\Xi_{13})=A_{18}$ として $a-\Xi_{13}$ で消す。
20次も前版の定理により $J_{20}$ に入り、15次場で消せる。

\begin{proposition}[二十二次の正確な有限補正式]
\label{prop:v16-exact-twenty-second-correction}
初回変更の22次係数は
\begin{equation}\label{eq:v16-initial-twenty-second}
\begin{split}
 g_{22}={}&f_{22}+df_{16}(X_7)+df_{15}(X_8)
                +\tfrac12 f_{10}''(X_7,X_7)\\
 &+3U_4\bigl(f_{18}+df_{12}(X_7)+\tfrac12P''(X_7,X_7)\bigr)
             -V_6\bigl(f_{16}+df_{10}(X_7)\bigr).
\end{split}
\end{equation}
上の純粋な有限正規化の後の22次係数は正確に
\begin{equation}\label{eq:v16-exact-twenty-second}
 R_{22}^{\mathrm{fin}}
 =g_{22}-dA_{14}(\Xi_9)+\tfrac12P''(\Xi_9,\Xi_9)
                         -c\,dH(\Xi_{13}).
\end{equation}
15、16、20次消去は22次係数に寄与しない。
\end{proposition}

\begin{proof}
初回の座標場の次数増分は6と7である。
22次に到達する一次項は $df_{16}(X_7)$ と $df_{15}(X_8)$、
二次項は $\tfrac12 f_{10}''(X_7,X_7)$ である。
混合二次項に必要な $f_9$ と、$X_8$ の平方に必要な $f_8$ は零である。
三階以上のTaylor項は、最小次数の $P$ に作用しても24次以上となる。
四次単元を掛けるには18次係数、六次単元を掛けるには16次係数を用いるので、
\eqref{eq:v16-initial-twenty-second}の残る二項を得る。
特に $3U_4P''(X_7,X_7)/2$ も22次に含まれている。

14次消去の増分は8である。
22次には14次項の一次微分と、$P$ の二階微分が寄与する。
15次、16次消去の増分は9、10であるが、対応する13次、12次係数は既に零である。
18次消去の増分12は10次項 $cH$ に作用するので、
$-c\,dH(\Xi_{13})$ を加える。
20次消去の増分14に対応する8次係数も零である。
各段の二階以上の残る寄与は、表示した14次消去の平方を除いて22次を超える。
\end{proof}

\nolinkurl{v16_degree_ledger.py} は、元の係数の次数、各座標場の次数増分、
merged単元の三項を独立に列挙し、初回の九項、単元を残す14次消去の五項、
後段で18次消去のHessian作用だけが残ることを有理数のTaylor係数とともに検算する。
この次数検算は、接触像への所属や全次数偶化の証明とは別の確認である。

\begin{remark}[単元を残す規約での同じ台帳]
14次を $A_{14}=dP(\Xi'_9)+PB_8$ とし、
$(1-B_8)g(a-\Xi'_9)$ で消す場合、その22次係数は正確に
\[
 g_{22}-dA_{14}(\Xi'_9)+\tfrac12P''(\Xi'_9,\Xi'_9)-PB_8^2
\]
となる。最後の項を落としてはならない。
この規約での18次入力は $g_{18}-c(dH(\Xi'_9)+B_8H)$ である。
その18次入力を $dP(\Xi'_{13})+PB_{12}$ として消すなら、
22次へ加える項は $-c(dH(\Xi'_{13})+B_{12}H)$ となる。
15、16、20次の単元からの22次寄与は、それぞれ13、12、8次係数が零なので消える。
これらの規約差を商で捨てる際は、$P\mathcal U_{16}\subset J_{22}$、
二階作用のEuler方向、および $H\mathcal U_{12}\subset J_{22}$ を用いる。
\end{remark}

\subsection{通常商に残る三つの二階作用}

$f_{16}\in J_{16}$、$f_{15}\in J_{15}$ により
\[
 dP(Z_{16})=f_{16},\quad Z_{16}\in\mathcal G_{11},\qquad
 dP(Z_{15})=f_{15},\quad Z_{15}\in\mathcal G_{10}
\]
と選ぶ。添字 $Z_{16},Z_{15}$ は持ち上げる係数の次数を表す。

\begin{theorem}[二十二次の実際の類の代表]
\label{thm:v16-twenty-second-correction-representative}
通常像 $J_{22}$ を法として
\begin{equation}\label{eq:v16-twenty-second-correction-representative}
\begin{split}
 [R_{22}^{\mathrm{fin}}]=\bigl[ f_{22}
  &+P''(X_7,Z_{16})+P''(X_8,Z_{15})
    +\tfrac c2H''(X_7,X_7)\\
  &-\tfrac12P''(\Xi_9,\Xi_9)-c\,dH(\Xi_{13})\bigr].
\end{split}
\end{equation}
最後の項は $J_{22}$ に属する。
さらに $dH(\mathcal G_7)\subset J_{16}$ の純粋な線形持ち上げ
\[
 W:\mathcal G_7\longrightarrow\mathcal G_{11},\qquad
 dP(W(X))=dH(X)
\]
を選ぶと
\begin{equation}\label{eq:v16-three-quadratic-families}
\boxed{\displaystyle
 [R_{22}^{\mathrm{fin}}]=
 \left[f_{22}+P''\!\left(X_7,Z_{16}+\tfrac c2W(X_7)\right)
       +P''(X_8,Z_{15})-\tfrac12P''(\Xi_9,\Xi_9)\right].}
\end{equation}
従って必要な残余補正は
$P''(\mathcal G_7,\mathcal G_{11})$、
$P''(\mathcal G_8,\mathcal G_{10})$、
$P''(\mathcal G_9,\mathcal G_9)$ の三族だけである。
\end{theorem}

\begin{proof}
恒等式 $dP(Z_{16})=f_{16}$ を $X_7$ で微分し、
$dP(Z_{15})=f_{15}$ を $X_8$ で微分すると
\[
 df_{16}(X_7)\equiv P''(X_7,Z_{16}),\qquad
 df_{15}(X_8)\equiv P''(X_8,Z_{15})\pmod{J_{22}}
\]
を得る。同様に $A_{14}=dP(\Xi_9)$ から
\[
 -dA_{14}(\Xi_9)+\tfrac12P''(\Xi_9,\Xi_9)
       \equiv-\tfrac12P''(\Xi_9,\Xi_9)\pmod{J_{22}}
\]
である。

第14版の命題\ref{prop:v14-all-lower-corrections-in-image}により
$P''(\mathcal G_7,\mathcal G_7)\subset J_{18}$ である。
また $f_{18}\in J_{18}$、$f_{16}\in J_{16}$ であり、
$f_{12}=PV_6-dP(X_7)$ を微分すると
\[
 df_{12}(X_7)+\tfrac12P''(X_7,X_7)\in J_{18}.
\]
さらに $df_{10}(X_7)\in J_{16}$ は
$f_{10}=cH-3PU_4$ と $dH(\mathcal G_7)\subset J_{16}$ から従う。
従って\eqref{eq:v16-initial-twenty-second}の単元を含む最後の二群は全て $J_{22}$ に入る。

$f_{10}''=cH''-3(PU_4)''$ において、
\[
 (PU_4)''(X_7,X_7)
 =U_4P''(X_7,X_7)+2dP(X_7)dU_4(X_7)
                        +P\,U_4''(X_7,X_7)
\]
の三項も $J_{22}$ に入る。
これで\eqref{eq:v16-twenty-second-correction-representative}を得る。

最後に $dP(W(X_7))=dH(X_7)$ を $X_7$ で微分すると
\[
 H''(X_7,X_7)+dH(DX_7\cdot X_7)
 =P''(W(X_7),X_7)+dP(D(W(X_7))\cdot X_7).
\]
$DX_7\cdot X_7\in\mathcal G_{13}$ と
$dH(\mathcal G_{13})\subset J_{22}$ により、
$H''(X_7,X_7)\equiv P''(W(X_7),X_7)$ である。
同じ包含により $dH(\Xi_{13})$ も消える。
全ての純粋な持ち上げの差は $dP$ の核に入り、
第15版の補題\ref{lem:v15-hessian-pairing-quotient}により表示した類に影響しない。
\end{proof}

\subsection{十七次の奇項と全次数偶化の比較}

22次では、前版の「17次消去を明示しなくてよい」という議論を、
次数の確認なしにそのまま使うことはできない。
そこで、17次の奇項が通常像に入ることを先に示す。

\begin{lemma}[残る十七次項の純粋消去]
\label{lem:v16-seventeenth-odd-removal}
上の有限正規化の後に残る17次項 $R_{17}$ は $J_{17}$ に入る。
12次の普遍座標場によってこれを消しても、22次係数は変わらない。
\end{lemma}

\begin{proof}
有限正規化の境界と、第15版の全次数偶正規形を比較する。
両者の差は17次以上なので、対数補題により17次の差は
\[
 -R_{17}=Q_{11}P+cQ_7H
\]
である。$Q_{11}P\in J_{17}$ であり、
系\ref{cor:v16-all-hessian-first-actions}の $r=8$ の場合と
$\mathcal U_7=0$ から $Q_7H=dH(Y_8)\in J_{17}$ となる。
従って $R_{17}\in J_{17}$ である。

$dP(\Xi_{12})=R_{17}$ を選び $a-\Xi_{12}$ で消す。
22次への一次寄与には11次係数が必要だが零である。
$P$ の二階寄与は28次からなので、22次を変えない。
18次、20次の既に消した係数も変わらない。
\end{proof}

\begin{corollary}[二十二次の偶正規形との一致]
\label{cor:v16-evenization-twenty-second-class}
第15版の全次数偶正規形における22次係数 $R_{22}^{\mathrm{ev}}$ と、
上の有限正規化の係数は
\[
 R_{22}^{\mathrm{ev}}-R_{22}^{\mathrm{fin}}\in I_{22}=J_{22}
\]
を満たす。従って式\eqref{eq:v16-three-quadratic-families}は、
実際の全次数偶正規形の22次類を与える。
\end{corollary}

\begin{proof}
前補題により17次を消した後は
\[
 f^{\mathrm{fin}}=P+cH+R_{19}+R_{21}+R_{22}^{\mathrm{fin}}+O(a^{23})
\]
である。偶正規形と比較する $T$ に対し $Tf^{\mathrm{fin}}-f^{\mathrm{fin}}$ は19次以上なので、
対数補題により $Tf^{\mathrm{fin}}-f^{\mathrm{fin}}\equiv Lf^{\mathrm{fin}}\pmod{(a)^{23}}$。
最小増分4から $L(R_{19}+R_{21})$ も23次以上である。
19次、21次の非零な比較式は奇数増分の条件にだけ入り、
偶数増分の鎖は\eqref{eq:v16-twenty-second-preservation}のままである。
従って22次の差は\eqref{eq:v16-effective-twenty-second-kernel}の像に入る。
\end{proof}

\begin{remark}[十七次を残した場合の二階対数項]
17次を残して直接比較する場合、22次には
\[
 (Tf-f)_{22}=Q_{16}P+cQ_{12}H+\tfrac12Q_5R_{17}
\]
という追加項が現れる。
実際、$Lf$ の17次項は $-R_{17}$、18次項は零なので、
$\tfrac12L^2f$ の22次項は $-\tfrac12Q_5R_{17}$ である一方、
$Lf$ 自身には $Q_5R_{17}$ が入る。
ただし $Q_5P=0$ なので $Q_5$ は通常像を保存し、
$R_{17}\in J_{17}$ から $Q_5R_{17}\in J_{22}$ である。
具体的に $Q=Y+b$、$QP=0$、$R=dP(Z)$ なら
\[
 QR=dP(DZ\cdot Y-DY\cdot Z)-P\,db(Z)
\]
となる。
従って直接比較しても同じ類を得るが、追加項の次数を無視してよいわけではない。
任意の低次安定化を含む座標変更の三階以上のTaylor項についても、
ここでは作用素全体の対数と低次保存条件がその相殺を含めて扱っている。
\end{remark}

\subsection{七十六列による十分な所属検査}

既存の低次像の次元は
\[
\begin{array}{c|ccccc}
 &J_{12}/P\mathcal U_6&J_{13}/P\mathcal U_7&J_{14}/P\mathcal U_8
 &J_{15}/P\mathcal U_9&J_{16}/P\mathcal U_{10}\\\hline
\dim&2&1&7&6&21
\end{array}
\]
である。各商の基底を持ち上げた場を、それぞれ
$G_{7,i}$、$G_8$、$G_{9,j}$、$G_{10,k}$、$G_{11,l}$ とする。

\begin{proposition}[二十二次補正の有限族への帰着]
\label{prop:v16-seventy-six-lower-pairings}
以下の76個の共変量が全て $J_{22}$ に属すると仮定する。
\[
\begin{array}{ll}
 \operatorname{Hess}_F(q)(dq(G_{7,i}),dq(G_{11,l})),&1\le i\le2,\ 1\le l\le21,\\
 \operatorname{Hess}_F(q)(dq(G_8),dq(G_{10,k})),&1\le k\le6,\\
 \operatorname{Hess}_F(q)(dq(G_{9,i}),dq(G_{9,j})),&1\le i\le j\le7.
\end{array}
\]
このとき
\begin{equation}\label{eq:v16-raw-twenty-second-suffices}
 [R_{22}^{\mathrm{ev}}]=[d_{22}^{\mathrm T}]
                         \quad\text{in }\mathcal C_{22}/J_{22}.
\end{equation}
ここで仮定は22次の普遍所属であり、少数点での零値だけに置き換えてはならない。
\end{proposition}

\begin{proof}
第15版の二階作用の補題により、$P''(Y,Z)$ の通常類は各 $dP$ 像を
$P$ 倍の単元像で割った商だけに依存する。
また三族のいずれも $d^2q(Y,Z)\in\mathcal D_{18}$ なので、
\[
 P''(Y,Z)=\operatorname{Hess}_F(q)(dqY,dqZ)+dF(d^2q(Y,Z))
\]
の最後の項は、全次数Pfaffian全射によって $J_{22}$ に入る。
従って76個の代表の所属から三族全体の所属が従う。
式\eqref{eq:v16-three-quadratic-families}と
系\ref{cor:v16-evenization-twenty-second-class}を適用すればよい。
\end{proof}

\paragraph{計算入力と現在の検査範囲.}
7、9、11次場は既存のグラフ基底、8次場は第15版で非零像を確認したグラフを用いる。
10次については、旧版のTate部分項の代わりに六つのグラフ場を新たに作った。
一般の $P=0$ 入力での $6\times6$ 行列を有理数上で再構成し、
非零行列式と、その1009を法とする値934を確認した。
既存の $\dim(J_{15}/P\mathcal U_9)=6$ と合わせ、これらも基底を与える。
入力と行列は \nolinkurl{v16_g10_graph_basis.json} に保存した。

\nolinkurl{v16_lower_pairings_exact.py} は第15版の $F_D,F_H,a_b$ を用いる。
四点
\[
 (F_D,a_1),\quad(F_D,a_2),\quad(F_H,a_1),\quad(F_H,a_2)
\]
で76列を有理数上で再計算する。
全76列は、この四点の全てで点ごと厳密に零となった。
従って候補汎関数
\[
 \widetilde\rho_{22}
 =\frac{16\operatorname{ev}_{(F_H,a_1)}-
                \operatorname{ev}_{(F_H,a_2)}}{3459456000}
\]
は全76列を消し、$\Theta_8U_4^2J_6$ に対する値は1である。
この評価結果だけでは、上の76列の $J_{22}$ 所属や、
この汎関数が通常接触像全体を消すことは従わない。
本稿では命題\ref{prop:v16-exact-rho22}で通常像と核の一致を証明し、
系\ref{cor:v16-all-lower-corrections-in-image}で76列の普遍所属を確定する。

% END contact/contact_twenty_second_ja.tex


% BEGIN detector/v16_detector_notes_ja.tex
% 第16版：全429接触方向のQ再縮約による22次通常商の厳密検出子。

\section{二十二次通常商を読む二点汎関数}
\label{sec:v16-exact-twenty-second-detector}

命題\ref{prop:v16-contact-rank-twenty-second}により
\[
 \dim\mathcal C_{22}=916,\qquad
 \dim J_{22}=915,\qquad
 J_{22}=P\mathcal U_{16}
    \oplus\operatorname{Span}\{dP(\Gamma_1),\ldots,dP(\Gamma_{429})\}
\]
である。ここで $\dim\mathcal U_{16}=486$、
429方向は286個の連結グラフと143個の下位場・単元の積からなる。
通常商の生成元を
\[
                      W_{22}=\Theta_8U_4^2J_6
\]
とする。本節では、これら全ての方向の二点での値を有理数上で再構成し、
通常商の係数を読む汎関数を確定する。

\subsection{全429方向の有理数による再縮約}

第15版と同じ
\[
 F_H=x_0^3+x_1^3+x_2^3+x_3^3+4x_4^3-6x_2x_3x_4,
\]
\[
 a_b=b\,e_0\wedge e_1+b^{-1}(e_2-e_4)\wedge(e_3-e_4)
\]
を用い、$z_1=(F_H,a_1)$、$z_2=(F_H,a_2)$ とする。
両点で $q=(0,0,1,1,1)$、$P=0$ であり、
\[
 \Theta_8(F_H,a_b)=-4b^2,\qquad
 U_4(F_H,a_b)=5b^4,\qquad J_6(F_H,a_b)=528b^6.
\]
従って
\begin{equation}\label{eq:v16-complement-two-values}
 W_{22}(z_1)=-52800,\qquad
 W_{22}(z_2)=-3460300800.
\end{equation}

\begin{proposition}[二十二次の厳密な係数汎関数]
\label{prop:v16-exact-rho22}
$\operatorname{ev}_i(R)=R(z_i)$ とし、
\begin{equation}\label{eq:v16-exact-rho22}
 \boxed{\displaystyle
 \rho_{22}=
 \frac{16\operatorname{ev}_1-\operatorname{ev}_2}{3459456000}}
\end{equation}
と置く。このとき
\[
 \rho_{22}(J_{22})=0,\qquad
 \rho_{22}(W_{22})=1,\qquad
                       \ker\rho_{22}=J_{22}.
\]
特に任意の $R\in\mathcal C_{22}$ に対して
\[
 [R]=\rho_{22}(R)[W_{22}]
                        \quad\text{in }\mathcal C_{22}/J_{22}
\]
である。汎関数は $\mathbb Q$ 上で定義し、$\mathbb C$ へ係数拡大して用いる。
\end{proposition}

\begin{proof}
429個の公開グラフを全て有理数上で縮約し、各列について
\[
                      dP(\Gamma_j)(z_2)=16dP(\Gamma_j)(z_1)
\]
を確認した。従って両評価の表示した結合は全429方向を消す。
$P(z_i)=0$ なので $P\mathcal U_{16}$ の486方向も点ごと零である。
式\eqref{eq:v16-complement-two-values}から
\[
 16W_{22}(z_1)-W_{22}(z_2)=3459456000
\]
となり、正規化は1である。
具体的に公開配列の33列目の接触方向と補元を並べた評価行列は
\[
 \begin{pmatrix}
 1920&-52800\\
 30720&-3460300800
 \end{pmatrix},\qquad
 \det=-6642155520000\ne0
\]
である。
この検出子は非零なので $\dim\ker\rho_{22}=916-1=915$。
通常像の次元と包含 $J_{22}\subset\ker\rho_{22}$ から核の一致を得る。
\end{proof}

\paragraph{積方向と正規化の検査.}
143個の積方向では、下位場 $\Gamma_r$ と単元グラフ $u_s$ は $r+s=17$ を満たし、
\[
 dP(u_s\Gamma_r)=u_s\,dP(\Gamma_r)
\]
という普遍恒等式を用いる。
単元のグラフを先に、下位場のグラフを後に置いた非連結17頂点グラフは、
最後の左頂点を接触方向の自由頂点に保つ。
\nolinkurl{v16_rho22_exact_certificate.py} は各積について、
この全17頂点グラフの値と、二つの因子を別々に縮約した値の積を、
両点で正確に照合した。
特に四頂点の単元はグラフそのものの値 $24U_4$ を用いており、
143列の途中で24による除算を挿入していない。
一方、補元 $W_{22}$ には既存の正規化された $U_4$ を用いる。

有理数の $a$ は一度 $Da$ にして整数縮約し、
全接触列は22次の斉次性により $D^{22}$ で割り戻す。
積の各因子もそれぞれの次数で割り戻し、同じ値になることを確認する。
これは全450点の有限体minorによる独立性の監査とは異なる、有理数上の所属判定である。

\subsection{低次補正の全消滅}

\begin{corollary}[二十二次の全低次正規化補正は通常像に入る]
\label{cor:v16-all-lower-corrections-in-image}
\[
 \begin{split}
 P''(\mathcal G_7,\mathcal G_{11})&\subset J_{22},\\
 P''(\mathcal G_8,\mathcal G_{10})&\subset J_{22},\\
 P''(\mathcal G_9,\mathcal G_9)&\subset J_{22}.
 \end{split}
\]
従って実際の22次偶正規形の係数は
\begin{equation}\label{eq:v16-all-corrections-vanish}
 R_{22}^{\mathrm{ev}}-d_{22}^{\mathrm T}\in J_{22}
\end{equation}
を満たす。
\end{corollary}

\begin{proof}
命題\ref{prop:v16-seventy-six-lower-pairings}の76列は、
$z_1,z_2$ のそれぞれで有理数上の再縮約により全て零であった。
従って前命題により全76列は $J_{22}$ に入る。
同命題のkernel方向、Euler方向、および $d^2q$ の持ち上げの議論によって、
三族全体が $J_{22}$ に入る。

Hessian一次作用 $dH(\mathcal G_{13})$ と $H\mathcal U_{12}$ の全体は、
前節までの構造定理により既に $J_{22}$ に入っている。
$H''(\mathcal G_7,\mathcal G_7)$ も
定理\ref{thm:v16-twenty-second-correction-representative}で第一の族へ還元した。
従って有限正規化の全補正が消え、
系\ref{cor:v16-evenization-twenty-second-class}により全次数偶化後の係数にも同じ結論が成り立つ。
\end{proof}

\subsection{実際の二十二次類}

\begin{proposition}[二十二次の実際の境界類の消滅]
\label{prop:v16-actual-twenty-second-class}
未正規化Tate係数と実際の正規化後の係数は、ともに $J_{22}$ に属する。
\[
                d_{22}^{\mathrm T}\in J_{22},\qquad
                [R_{22}^{\mathrm{ev}}]=0.
\]
\end{proposition}

\begin{proof}
命題\ref{prop:v16-fresh-tate22-values}による新規Tate再帰の値を
\[
 A_{22}=\frac{1329012984671053041316939}{630118440000000000}
\]
と書けば、
\[
             d_{22}^{\mathrm T}(z_1)=A_{22},\qquad
             d_{22}^{\mathrm T}(z_2)=16A_{22}.
\]
従って $\rho_{22}(d_{22}^{\mathrm T})=0$ であり、
前命題の核の同定から $d_{22}^{\mathrm T}\in J_{22}$ となる。
式\eqref{eq:v16-all-corrections-vanish}を用いれば正規化後の類も零である。
対角入力 $(F_D,a_1),(F_D,a_2)$ の新規値はいずれも零であり、
四点の全データは同じTate命題とその証明書に収めた。
\end{proof}

\paragraph{証明書の範囲.}
\nolinkurl{v16_rho22_exact_certificate.json} には、429グラフの両点での有理数値、
143積方向の各因子と積、補元値、選択した小行列、76補正の汎関数値を保存した。
Tate入力については全三次テンソルと交代行列 $a$ を一致させ、
保存された18個の定数ブロックから0次〜22次の行列式係数を再計算した。
この保存ブロックの監査を、最初からのMaurer--Cartan再帰生成とは区別する。
通常像の独立性、三元切片による補元の非所属、低次補正の数学的な縮約は、
それぞれの命題を別の入力として用いる。

% END detector/v16_detector_notes_ja.tex


% BEGIN tate/v16_tate_notes_ja.tex
\section{二十二次Tate係数の新規再構成}
\label{sec:v16-tate-twentysecond}

二十二次では、前版の十六個の定数ブロックを延長するだけでは足りない。
十八個の定数ブロックに必要な内部次数を確保するため、四つの入力点で
$A_1$ からMaurer--Cartan再帰を再実行した。
さらに一つの対角点では有理数の実装を変えて再実行し、
全十八ブロックと全行列式係数を比較した。

\subsection{十八次ブロックを得る切断}

\begin{lemma}[二十二次の厳密な内部次数切断]
\label{lem:v16-tate22-truncation}
二十二次以下の行列式係数には $K_1,\ldots,K_{18}$ だけが必要である。
これらは $A_1,\ldots,A_{14}$ を全て計算し、続いて
$R_{15},R_{16},R_{17},R_{18}$ をそれぞれ内部 $x$ 次数
$4,3,2,1$ 以下で求めれば厳密に得られる。
対応する $A_{15},A_{16},A_{17},A_{18}$ は、それぞれ
$x$ 次数 $3,2,1,0$ 以下まで正確である。
\end{lemma}

\begin{proof}
行列式の各項は正次数の五つのブロックの積なので、総次数22以下では
各ブロックの次数は18以下である。
Clifford積は $x$ 次数を下げず、Eulerホモトピーは1下げる。
内部重み $2\deg_x-\#\theta+\#\eta=3$ により、各 $A_r$ の
$x$ 次数は4以下である。
$n=15,\ldots,18$ の曲率に必要な $x$ 次数は $19-n$ 以下であり、
その積の各因子の添字 $i\le n-1$ に対して
$18-i\ge19-n$ である。従って、それ以前の切断が必要な項を失わないことが
帰納的に従う。ホモトピーを適用すると $18-n$ 次以下の補正を得る。
\end{proof}

第15版で保存した $A_{13},A_{14}$ は内部次数を切断していたため、
本節では両者の全項を改めて計算している。
Hesse型の計算では補題\ref{lem:v15-hesse-symbol-product}の
$S=1+C$、$S^{-1}=1-C$ を用い、常に自然な外積記号に戻してから
Eulerホモトピーを適用する。
このClifford積の実装は第15版からバイト単位で変更していない。

\subsection{四点の厳密な二十二次値}

\[
\begin{split}
 F_D&=x_0^3+x_1^3+x_2^3+x_3^3-2x_4^3,\\
 F_H&=x_0^3+x_1^3+x_2^3+x_3^3+4x_4^3-6x_2x_3x_4,\\
 a_b&=b\,e_0\wedge e_1+
       b^{-1}(e_2-e_4)\wedge(e_3-e_4)\qquad(b=1,2)
\end{split}
\]
とする。いずれの点でも $q(a_b)=(0,0,1,1,1)$、$P=0$ である。
この四つの入力点では $H_{10}=0$ であり、$q$ 変数に関する
Hessian行列式の勾配も零となる。

\Needspace{7\baselineskip}
\begin{proposition}[新規再構成による四点の二十二次係数]
\label{prop:v16-fresh-tate22-values}
\[
 A_{22}=\frac{1329012984671053041316939}{630118440000000000}
\]
とおくと、未正規化のTate行列式係数は
\[
\boxed{\bigl(d_{22}^{\mathrm T}(F_D,a_1),d_{22}^{\mathrm T}(F_D,a_2),
             d_{22}^{\mathrm T}(F_H,a_1),d_{22}^{\mathrm T}(F_H,a_2)\bigr)
       =(0,0,A_{22},16A_{22}).}
\]
従って命題\ref{prop:v16-exact-rho22}の検出子について
$\rho_{22}(d_{22}^{\mathrm T})=0$ である。
\end{proposition}

\begin{proof}
各点の再帰を最初から有理数上で実行し、保存した内部次数の全範囲で
曲率の閉性とMaurer--Cartan方程式を照合した。
定数ブロックの規約は $K_1=-a$、$K_2=\operatorname{Hess}_F(q)/3$、
$d^{\mathrm T}=\frac12\det(\sum K_r)$ である。
行列式を120置換による展開と列部分集合による動的計画法で別々に計算し、
0次から22次までの全23係数を有理数上で比較した。
得られた非零の二十二次値は
\[
 d_{22}^{\mathrm T}(F_H,a_1)
   =\frac{1329012984671053041316939}{630118440000000000},\qquad
 d_{22}^{\mathrm T}(F_H,a_2)
   =\frac{1329012984671053041316939}{39382402500000000}.
\]
後者は前者の16倍であるため、検出子の分子
$16d_{22}^{\mathrm T}(F_H,a_1)-d_{22}^{\mathrm T}(F_H,a_2)$ は零である。
\end{proof}

\subsection{独立検証と証明の範囲}

各再帰では、PBW積とその基底変換を用いない32次元外積作用素を
$x=(1,2,3,4,5)$、素数 $1000003$ で別に構成し、
完全な $A_1,\ldots,A_{14}$ から得る0次から14次までの作用素関係を検証した。
全十八の定数ブロックにも独立な外積作用素による照合を行った。
これは有理数上の再帰と切断証明に加える実装検査である。

$(F_D,a_1)$ では、共通分母による整数実装と各係数を
\texttt{Fraction} とする実装を独立に走らせた。
両者の全十八ブロック、全23行列式係数、二十二次値が一致した。
さらに、前版に同じ入力がある $(F_D,a_1),(F_D,a_2),(F_H,a_1)$ では、
前版の十六ブロックと0次から20次までの全係数に一致した。
これらは保存データに対する標準ライブラリだけの独立監査でも再確認した。

配布する \texttt{v16\_tate\_saved\_audit.py} は保存ブロックから
行列式を再計算する監査であり、新規Maurer--Cartan再帰そのものを再実行しない。
新規再帰用の二つのドライバ、四点五回の出力、ログを別に同梱する。
局所的なチェックポイントは配布物に含めず、ドライバも既存の
チェックポイントを入力として読まない。

本節の四点での値だけから普遍接触像への所属は従わない。
命題\ref{prop:v16-exact-rho22}の完全な商の検出、
系\ref{cor:v16-all-lower-corrections-in-image}の低次正規化補正の消滅と
系\ref{cor:v16-evenization-twenty-second-class}の偶正規形との比較を合わせて初めて、
実際の二十二次接触類が零であることを結論する。

% END tate/v16_tate_notes_ja.tex


% BEGIN contact/twenty_fourth_normal_form_ja.tex
% 第16版の確定結論。22次像、全低次補正、fresh Tate値に依存する。

\section{二十二次の消去と二十四次開始の正規形}
\label{sec:v16-twenty-fourth-normal-form}

\begin{theorem}[二十四次開始の普遍代数的収束偶正規形]
\label{thm:v16-twenty-fourth-normal-form}
実際の境界の22次類は零であり、任意の滑らかな三次式 $F$ に対して
\begin{equation}\label{eq:v16-twenty-fourth-normal-form}
 \boxed{\displaystyle
 d^{\mathrm T}\sim
 P_F+\frac1{96}\det\operatorname{Hess}_F(q(a))
       +\sum_{\substack{n\ge24\\n\ \mathrm{even}}}R_n,
 \qquad R_n\in\mathcal C_n}
\end{equation}
という全次数の偶正規形が存在する。
接触変換は普遍同変な代数的冪級数であり、座標変更の線形項は恒等写像、
単元の定数項は1である。各形式係数は $F$ の係数に多項式的で、
各固定 $F$ では変換・逆変換・正規形の級数は原点近傍で収束する。
24次以降の実際の類の消滅は、この定理の主張に含めない。
\end{theorem}

\begin{proof}
第15版の定理\ref{thm:v15-twenty-second-normal-form}により
\[
 f_{\mathrm{ev}}=P_F+\frac1{96}H+R_{22}^{\mathrm{ev}}
                             +R_{24}^{\mathrm{ev}}+\cdots
\]
という普遍同変な代数的収束偶正規形がある。
系\ref{cor:v16-all-lower-corrections-in-image}は
\[
 R_{22}^{\mathrm{ev}}-d_{22}^{\mathrm T}\in J_{22}
\]
を示す。この比較には、17次の奇項の純粋消去と、
全偶化との対数作用による比較も含まれている。
命題\ref{prop:v16-exact-rho22}と
命題\ref{prop:v16-fresh-tate22-values}により
\[
 \rho_{22}(d_{22}^{\mathrm T})
   =\frac{16A_{22}-16A_{22}}{3459456000}=0,
 \qquad\ker\rho_{22}=J_{22}.
\]
従って $R_{22}^{\mathrm{ev}}\in J_{22}$ である。
ここで $I_{22}=J_{22}$ は構造定理でも確定しており、
異なる下位正規化による有効類と通常類の区別も解消している。

有限次元の普遍共変量空間で
\[
 \Xi_{17}\in\mathcal G_{17},\qquad
                       dP(\Xi_{17})=R_{22}^{\mathrm{ev}}
\]
と選ぶ。これは $F$ の判別式等で割る操作ではなく、
$F$ に16次、$a$ に17次の普遍多項式な場の持ち上げである。
奇座標変更 $a\mapsto a-\Xi_{17}(F,a)$ を行うと22次項が消える。
Hessian項への最初の寄与は26次、$P$ の二階Taylor項は38次からであるため、
既存の低次正規形を保ち、次に残り得る次数は24となる。
座標変更が奇関数なので全次数の偶性も保たれる。

今回の追加変更は有限多項式で、原点でのJacobianは恒等写像である。
その形式逆は代数的冪級数であり、固定した $F$ に対しては解析的にも局所収束する。
第15版の変換との有限合成は、代数性、局所収束性、普遍同変性、
各形式係数の $F$ 多項式性を保つ。
これで式\eqref{eq:v16-twenty-fourth-normal-form}を得る。
\end{proof}

\begin{remark}[確定した商と残る課題]
通常22次商は一次元で、非零な $[\Theta_8U_4^2J_6]$ によって生成される。
その商で実際のTate境界の係数は零であった。
また、下位jetが $P+cH$ である場合の $I_N=J_N$ は全次数の構造定理であるが、
これ自体は実際の高次係数の零性を与えない。
通常商の全次数の構造と、24次以降の実際の類は引き続き調べる必要がある。
特に全境界が $P_F+H/96$ と接触同値であることは、まだ結論していない。
\end{remark}

% END contact/twenty_fourth_normal_form_ja.tex

\clearpage
% BEGIN theory/v17_theta_annihilator_ja.tex
% 第17版：Theta8の全単元環におけるannihilator。
% 新しい外部入力は不要。第16版のPluecker素性・dq全射・横断単射を継承。

\section{行列束の恒等式と八次共変量の消滅イデアル}
\label{sec:v17-theta-annihilator}

第16版の記号
\[
 \Delta=\det V,\quad E=V^*\otimes\Delta^{-1},\quad
 L=\Delta^{-3},\quad
 \mathcal U=\bigoplus_{s\ge0}\mathcal U_s,\quad
 \mathcal N_s=\{u\in\mathcal U_s:u|_{\Sigma}=0\}
\]
を用いる。通常像を $J_n=dP(\mathcal G_{n-5})$ とする。
本節では、$H$ に対して既知だったPl\"uckerイデアルによる消去を
$\Theta_8$ に対しても示し、その消滅イデアルを正確に決める。
通常商の全ての元がこの二つの共変量から生成されるとは、まだ主張しない。

\subsection{対称行列と交代行列の束}

$T=\operatorname{Hess}_F(q(a))$、$g=dF_{q(a)}$ と置く。
交代テンソル $a$ も、$T$ と同じ型の双線形形式とみなせる。実際
\[
 \Lambda^2E^*\otimes L=\Lambda^2V\otimes\Delta^{-1}=W,
 \qquad
 \operatorname{Sym}^2E^*\otimes L=\operatorname{Sym}^2V\otimes\Delta^{-1}.
\]
従って独立なスカラー変数 $z$ に対して $T+za:E\to E^*\otimes L$
およびその行列式・余因子行列を自然に定義できる。
$\det E^*=\Delta^6$ なので
\[
 \det(T+za)\in L[z],\qquad
 \operatorname{adj}(T+za)\in(E\otimes E)[z].
\]
以後の行列表示は体積形式を選んだものであり、
このdeterminant twistを省略した別のテンソルを導入しているわけではない。

\begin{lemma}[行列束の三つの係数]
\label{lem:v17-matrix-pencil-coefficients}
普遍多項式恒等式
\begin{equation}\label{eq:v17-determinant-pencil}
 \det(T+za)=H+216z^2\Theta_8+6z^4P
\end{equation}
が成立する。さらに
\[
 B_2=[z^2]\operatorname{adj}(T+za)
 \in\operatorname{Sym}^2E
\]
は二重次数 $(\deg_F,\deg_a)=(2,6)$ の普遍共変量である。
\end{lemma}

\begin{proof}
まず $T$ を任意の対称五行五列行列、$a$ を任意の交代行列としてよい。
転置によって $z$ は $-z$ に変わるので、行列式は $z$ の偶多項式である。
定数項は $\det T$ であり、五次の係数は交代奇数行列の行列式なので零である。
対角な $T=\operatorname{diag}(t_0,\ldots,t_4)$ では、
二次の係数は
\[
 \sum_{i<j}a_{ij}^2\prod_{k\ne i,j}t_k
\]
であり、式\eqref{eq:boundary-theta-eight}の右辺の216倍と一致する。
四次の係数は $\sum_i t_iq_i^2=q^{\mathsf T}Tq$ である。
両辺は同じ合同変換則を持つ多項式であり、非退化対称行列は合同変換で対角化できる。
その集合は稠密なので、任意の $T$ に対して係数の等式が成立する。
三次式のEuler恒等式から $Tq=2g$、
$q^{\mathsf T}Tq=6F(q)=6P$ となり、表示した式を得る。

余因子行列についても
$\operatorname{adj}(T+za)^{\mathsf T}=\operatorname{adj}(T-za)$
なので偶数係数 $B_2$ は対称である。
四つの因子のうち二つが $T$、二つが $a$ であり、
$T$ の二重次数は $(1,2)$ なので、$B_2$ の二重次数は $(2,6)$ である。
\end{proof}

\begin{proposition}[$\Theta_8$ のPl\"ucker消去恒等式]
\label{prop:v17-theta-plucker-identity}
分母に $F$ や $a$ の多項式を用いず、
\begin{equation}\label{eq:v17-theta-plucker-identity}
 216\Theta_8\,q=2B_2g
\end{equation}
が普遍的に成立する。特に、普遍係数環 $R=\mathbb C[F,a]$ で
\[
 \Theta_8(q_0,\ldots,q_4)\subset(g_0,\ldots,g_4)
\]
である。
\end{proposition}

\begin{proof}
Pfaffianの恒等式 $aq=0$ と $Tq=2g$ を用いると
\[
 \det(T+za)\,q
 =\operatorname{adj}(T+za)(T+za)q
 =2\operatorname{adj}(T+za)g.
\]
$z^2$ の係数を比較すればよい。
\end{proof}

\subsection{全単元環に対する正確な消滅イデアル}

\begin{theorem}[$\Theta_8$ に作用する単元環]
\label{thm:v17-theta-unit-annihilator}
全ての $s\ge0$ について
\begin{equation}\label{eq:v17-theta-unit-annihilator}
 \Theta_8u\in J_{s+8}
 \quad\Longleftrightarrow\quad
 u\in\mathcal N_s
 \qquad(u\in\mathcal U_s).
\end{equation}
従って、通常接触商における $[\Theta_8]$ の $\mathcal U$-消滅イデアルは
正確に $\mathcal N=\bigoplus_s\mathcal N_s$ であり、
\[
 \mathcal U[\Theta_8]\simeq(\mathcal U/\mathcal N)[-8]
 \simeq\mathbb C[\overline U_4,\overline J_6][-8].
\]
ここで $[-8]$ は生成元を8次に置くという意味である。
\end{theorem}

\begin{proof}
$u\in\mathcal N_s$ とする。
補題\ref{lem:v16-plucker-prime-ideal}と有限次数での完全可約性によって、
$E^*$ 値の同変多項式 $b$ を用い
\[
 u=\langle b,q\rangle,\qquad
 (\deg_Fb,\deg_ab)=(s,s-2)
\]
と書ける。$s<2$ なら、この次数のイデアル部分は零である。
式\eqref{eq:v17-theta-plucker-identity}と $B_2$ の対称性から
\[
 \Theta_8u=\left\langle g,\frac1{108}B_2b\right\rangle.
\]
$v=B_2b/108$ は $E$ 値で、二重次数は $(s+2,s+4)$、
従って $v\in\mathcal D_{s+4}$ である。
命題\ref{prop:v12-uniform-pfaffian-surjectivity}の全射
$dq:\mathcal G_{s+3}\twoheadrightarrow\mathcal D_{s+4}$
を用いれば、$\Theta_8u\in J_{s+8}$ を得る。

逆に $\Theta_8u\in J_{s+8}$ とする。
第16版の横断切片に制限すると、補題\ref{lem:v16-contact-normal-cubic-jet}と
補題\ref{lem:v16-theta-normal-hessian}により
\[
 \overline u(f)\det(\operatorname{Hess}_f/6)
 \in\operatorname{Jac}(f)_3
\]
である。式\eqref{eq:v16-fermat-normal-minor}が与える非空開集合では
Hessian類は非零なので、$\overline u(f)=0$ となる。
多項式性によって全ての三元三次式で零であり、
階数二制限の記述から $u\in\mathcal N_s$ を得る。
環の最後の同型は定理\ref{thm:v16-unit-restriction-ring}による。
\end{proof}

\begin{remark}[この結論が決めていないこと]
$\mathcal U$ 全体による $\Theta_8$ の生成部分加群は正確に決まった。
ただし、通常商全体に別の生成元が存在しないことはこの定理から従わない。
第16版で得た $\mathbb C\cdot[H]$ は依然として別の非零方向であり、
さらに高次の未知の方向を排除するには別の議論が必要である。
\end{remark}

\subsection{独立な全係数検算}

\nolinkurl{theory/v17_structure_certificate.py} はPython標準ライブラリだけを用い、
対称 $T$ の15変数と交代 $a$ の10変数を全て独立な記号変数として計算する。
行列式の六つの係数、$aq=0$ の五成分、
$B_2Tq=216\Theta_8q$ の五成分を整数多項式の全係数で検証する。
さらに $[z^4]\operatorname{adj}(T+za)=qq^{\mathsf T}$ も25成分で照合する。
別に35個の三次係数と5個の独立変数でEuler恒等式を検証する。
これは有限標本による推測ではなく、表示した普遍恒等式そのものの係数検算である。

% END theory/v17_theta_annihilator_ja.tex


% BEGIN theory/v17_normal_projection_ja.tex
% 第17版：normal三次jetの全像・直和分解・普遍二階作用。
% 外部入力：Fisher, The Hessian of a genus one curve, Theorem 2.7(i), n=3。

\section{横断三次項の射影と通常商に残る未知の部分}
\label{sec:v17-normal-projection}

$\mathcal A=\mathcal U/\mathcal N
=\mathbb C[\overline U_4,\overline J_6]$ と置く。
第16版の単射を、本節では通常商からの射影と直和分解へ強める。
一方、横断三次項に現れない部分を明示して残すため、
この議論だけで全通常商の次元を決定したと取り違えない。

\subsection{全境界共変量の横断三次項}

$V=L_2\oplus Q_3$、$a_0=e_0\wedge e_1$ とし、
第16版の切片
\[
 a_\varepsilon=a_0+
 \varepsilon(p_4e_2\wedge e_3-p_3e_2\wedge e_4+p_2e_3\wedge e_4),
 \qquad F_f=x_0^3+x_1^3+f
\]
を使う。体積線の自明化は第16版と同じである。

\begin{lemma}[三次未満の横断項は現れない]
\label{lem:v17-covariant-normal-order}
$C\in\mathcal C_n$、$n\ge6$ とする。
$C(F,a_\varepsilon)$ の $\varepsilon^0,\varepsilon^1,\varepsilon^2$ 係数は、
全ての $F$ に対して零である。
三次係数は $F$ の $f\in\operatorname{Sym}^3Q_3$ 成分だけに依存し、
\[
 \gamma_n(C)(f,p)
 =[\varepsilon^3]C(F_f,a_\varepsilon)
\]
は $f$ の次数 $n-5$、$p$ の次数3の多項式である。
$f\mapsto\gamma_n(C)(f,-)$ は $\mathrm{SL}(Q_3)$ 共変写像
$\operatorname{Sym}^3Q_3\to\operatorname{Sym}^3Q_3$ となる。
\end{lemma}

\begin{proof}
$\lambda(t)=\operatorname{diag}(t,t,1,1,1)$ は $a_0$ を固定する。
横断変数 $b\in\Lambda^2Q_3\otimes\Delta^{-1}$ の重みは $-2$、
出力 $L=\Delta^{-3}$ の重みは $-6$ である。
従って $\varepsilon^k$ の係数に現れる $F$ の重みは $2k-6$ でなければならない。
$F$ の重みは非負なので $k<3$ は零であり、
$k=3$ では重み零の三元成分 $f$ だけが現れる。
この段階では $F_f$ に分離してから論じる必要はない。

$h\in\mathrm{SL}(Q_3)$ を $L_2$ 上の恒等写像で延長すると $a_0$ は固定され、
横断Pfaffian座標 $p$ は $h^{-\mathsf T}p$ に変わる。
従って
$\gamma_n(C)(hf,h^{-\mathsf T}p)=\gamma_n(C)(f,p)$
であり、出力三次式は $f$ と同じ表現で変換する。
\end{proof}

三元Hessianを
\[
 h_f=\frac1{216}\det\operatorname{Hess}_f
     =\det(\operatorname{Hess}_f/6)
\]
と正規化する。これは五変数の共変量 $H=H_{10}$ とは異なる。

\begin{theorem}[横断三次項が与える全次数の分裂射影]
\label{thm:v17-split-normal-projection}
各 $C\in\mathcal C_n$ は、唯一の
$A_C\in\mathcal A_{n-6}$、$B_C\in\mathcal A_{n-8}$ によって
\begin{equation}\label{eq:v17-normal-covariant-decomposition}
 \gamma_n(C)=A_C(\overline U_4,\overline J_6)\,f
            +B_C(\overline U_4,\overline J_6)\,h_f
\end{equation}
と書ける。負の次数の成分は零とする。
写像 $C\mapsto B_C$ は通常商に降り、次数付き $\mathcal U$ 加群の全射
\begin{equation}\label{eq:v17-normal-projection}
 \pi:\bigoplus_{n\ge6}\mathcal C_n/J_n
       \longrightarrow\mathcal A[-8]
\end{equation}
となる。右辺への $\mathcal U$ の作用は $\mathcal U\to\mathcal A$ を通じる。
$\overline u\mapsto[\Theta_8u]$ はこの全射の切断であり、
\begin{equation}\label{eq:v17-normal-splitting}
 \bigoplus_{n\ge6}\mathcal C_n/J_n
   =\mathcal U[\Theta_8]\oplus\mathcal K^{\mathrm{nor}},
 \qquad \mathcal K^{\mathrm{nor}}=\ker\pi
\end{equation}
という直和分解を得る。
\end{theorem}

\begin{proof}
Fisherの一次資料\cite[Definition~2.5, Lemma~2.6, Theorem~2.7(i)]{FisherHessian2012}
を $n=3$、標数零で用いる。
この定理は三元三次式の全係数空間上の多項式共変写像について、
恒等共変量 $f$ とHessian共変量が
$\mathbb C[c_4,c_6]$ 上の自由基底をなすというものである。
判別式を可逆にした環に限定した主張ではない。
Hessianを非零定数倍して $h_f$ とし、
定理\ref{thm:v16-unit-restriction-ring}の同じ不変式環
$\mathcal A=\mathbb C[\overline U_4,\overline J_6]$ を用いれば、
前補題から式\eqref{eq:v17-normal-covariant-decomposition}を得る。
$f$ と $h_f$ の係数次数は1と3なので、表示した次数も従う。

$C\in J_n$ なら、補題\ref{lem:v16-contact-normal-cubic-jet}により
$\gamma_n(C)\in\operatorname{Jac}(f)_3$ である。
$f=\frac13\sum p_i\partial_i f$ もこの空間に属する。
第16版の非空開集合では $[h_f]\ne0$ なので $B_C=0$ となり、
多項式性によって恒等的に零である。従って $\pi$ は通常商に降りる。

$u\in\mathcal U_s$ とすると、
前補題で $C$ の三次未満の横断項が消えることから
\[
 \gamma_{n+s}(uC)=\overline u\,\gamma_n(C)
\]
である。これが $\mathcal U$-線形性を示す。
補題\ref{lem:v16-theta-normal-hessian}により $\gamma_8(\Theta_8)=h_f$ なので、
$\pi([\Theta_8u])=\overline u$ となる。
定理\ref{thm:v17-theta-unit-annihilator}により切断は持上げ $u$ の選択に依存しない。
よって全射性と直和分解が従う。
\end{proof}

\begin{proposition}[未知の部分は横断四次以上で表せる]
\label{prop:v17-normal-invisible-representatives}
$\mathcal C_n^{\ge4}$ を、全 $F$ に対して上の切片で横断三次係数も零になる
$C\in\mathcal C_n$ の部分空間とする。この条件は非零階数二軌道の各点で
横断次数が少なくとも4であることと同値である。
すると
\begin{equation}\label{eq:v17-normal-invisible-representatives}
 \mathcal K_n^{\mathrm{nor}}
 =(\mathcal C_n^{\ge4}+J_n)/J_n.
\end{equation}
特に $[H]\in\mathcal K_{10}^{\mathrm{nor}}$ である。
\end{proposition}

\begin{proof}
$B_C=0$ なら $\gamma_n(C)=A_C f$ である。
$A_C$ を同じ多項式 $A_C(U_4,J_6)\in\mathcal U_{n-6}$ で持ち上げ、
$C-PA_C(U_4,J_6)$ と置き換える。
$PA_C(U_4,J_6)\in J_n$ であり、その三次係数は $A_Cf$ なので、
この代表は $\mathcal C_n^{\ge4}$ に入る。逆は定義から明らかである。
三次係数が $f$ だけに依存することと非零階数二軌道の推移性により、
全 $F$・軌道上の全点で同じ消失を得る。
$H$ は $q$ の五次式なので、錐に沿う次数は少なくとも5である。
\end{proof}

\begin{remark}[全商の生成問題の正確な残り]
第16版の単射は全次数の直和因子へ強まったが、
$\mathcal K^{\mathrm{nor}}=\mathbb C\cdot[H]$ はまだ証明していない。
これが全通常商の生成予想の残る主張である。
例えば24次では
\[
 \mathcal C_{24}/J_{24}
 =\mathbb C\cdot[\Theta_8U_4^4]\oplus
   \mathbb C\cdot[\Theta_8U_4J_6^2]\oplus\mathcal K_{24}^{\mathrm{nor}}.
\]
従って上界2は $\mathcal K_{24}^{\mathrm{nor}}=0$ と同値であり、
三元共変量の自由生成定理だけではこの上界は得られない。
本版では別の明示階数計算でこの消失を確定する
（命題\ref{prop:v17-ordinary-contact-rank24}）。
本節の $\mathcal K^{\mathrm{nor}}$ はKoszul複体の項
$\mathcal K_{k,n}$ やそのhomologyとは別の記号である。
\end{remark}

\subsection{普遍な場の横断一次作用はスカラーである}

\begin{lemma}[三元三次式からのEnd共変量]
\label{lem:v17-ternary-endomorphism-scalar}
三元三次式の係数に多項式的な
$\mathrm{SL}(Q_3)$ 同変写像
$M:\operatorname{Sym}^3Q_3\to\operatorname{End}(Q_3^*)$
は、全ての $f$ に対してスカラー行列である。
同変写像が $\mathrm{GL}(Q_3)$ のdeterminant指標でねじれていても同じ結論となる。
\end{lemma}

\begin{proof}
まず $M$ を $f$ の次数 $d$ の斉次成分 $M_d$ に分ける。
各成分も $\mathrm{SL}_3$ 同変であり、$h=\lambda h_0$、
$h_0\in\mathrm{SL}_3$ と書くと
\[
 M_d(hf)=(\det h)^d h^{-\mathsf T}M_d(f)h^{\mathsf T}.
\]
従って各斉次成分は自動的にdeterminant指標でねじれた
$\mathrm{GL}_3$ 共変写像へ延びる。以下は各 $M_d$ に適用してから和を取ればよい。
$f_t=x^3+y^3+z^3+6txyz$ とし、$\omega$ を原始三乗根とする。
$D=\operatorname{diag}(1,\omega,\omega^2)$ と三つの座標の巡回置換 $C$
は $\mathrm{SL}_3$ に属し、全ての $f_t$ を固定する。
$Q_3^*$ 上で $D$ は異なる三固有値を持つため、これと可換な行列は対角である。
さらに $C$ と可換なら三対角成分は等しい。
従って $M(f_t)$ はスカラーである。

写像
\[
 \mathrm{GL}_3\times\mathbb A^1\longrightarrow\operatorname{Sym}^3Q_3,
 \qquad(h,t)\longmapsto hf_t
\]
は優勢である。実際 $(1,0)$ での九つの行列方向の微分は
$3p_i^2p_j$（$i,j=1,2,3$）、最後の方向の微分は $6p_1p_2p_3$ である。
この十個は三次単項式の基底であり、対応する行列式は
$6\cdot3^9=118098\ne0$ である。
同変性とdeterminantによるスカラー倍は行列のスカラー性を保つ。
従って稠密な集合で $M$ はスカラーであり、多項式性により全 $f$ でそうなる。
\end{proof}

\begin{proposition}[Pfaffian方向の正確な横断一次項]
\label{prop:v17-scalar-normal-pfaffian-action}
$w\in\mathcal D_{r+1}$ とする。全 $F$ に対して
$w=Mq$ と同変に表示でき、
分離切片上の $Q_3^*$ 成分は
\begin{equation}\label{eq:v17-scalar-normal-pfaffian-action}
 w_Q(F_f,a_\varepsilon)
  =\varepsilon\alpha_w(f)p+O(\varepsilon^2)
\end{equation}
となる。$\alpha_w$ は $f$ に多項式的である。
この一次作用は $M$ の非一意性によらず、$w$ から定まる。
\end{proposition}

\begin{proof}
$w=Mq$ は補題\ref{lem:v16-plucker-vanishing-D}による。
$\lambda(t)$ に対する $Q_3^*\to Q_3^*$ ブロックの重みは零なので、
$M_{QQ}(F,a_0)$ は $F$ の三元成分 $f$ だけに依存する。
$\mathrm{SL}(Q_3)$ ではこれは共役作用に関して同変である。
$h\in\mathrm{GL}(Q_3)$ まで広げると、
$a_0\mapsto(\det h)^{-1}a_0$ と $a$ 次数 $r-1$ による
因子 $(\det h)^{r-1}$ が付くが、それは前補題で許したdeterminant指標である。
従って $M_{QQ}(f)=\alpha_w(f)\operatorname{id}$ となる。

切片では $q=(0,0,\varepsilon p)$ が正確に成立するので表示した一次項を得る。
二つの $M$ が同じ $w$ を表すなら、その $QQ$ ブロックの差は
全 $p$ を零に送るため零である。
これは $dq_{a_0}$ が三次元の横断空間から $Q_3^*$ への同型を与えることでも説明できる。
\end{proof}

\begin{theorem}[全二階接触補正は横断射影で消える]
\label{thm:v17-quadratic-actions-normal-zero}
$r,s\ge1$、$Y\in\mathcal G_r$、$Z\in\mathcal G_s$、$n=r+s+4$ とすると
\begin{equation}\label{eq:v17-quadratic-actions-normal-zero}
 \pi_n([P''(Y,Z)])=0.
\end{equation}
従って $P''(Y,Z)$ の通常商における類は
$\mathcal K_n^{\mathrm{nor}}$ に入る。
特に $\mathcal K_n^{\mathrm{nor}}=0$ と別に分かっている次数では、
全てのこの二階補正が $J_n$ に入る。
\end{theorem}

\begin{proof}
連鎖律により
\[
 P''(Y,Z)=\operatorname{Hess}_F(q)(dqY,dqZ)
              +dF_q\bigl(d^2q(Y,Z)\bigr).
\]
$d^2q(Y,Z)$ は $E$ 値で、二重次数は
$(r+s-2,r+s)$、従って $\mathcal D_{r+s}$ に属する。
全次数全射 $dq:\mathcal G_{r+s-1}\twoheadrightarrow\mathcal D_{r+s}$
から、第二項は全体として $J_{r+s+4}$ に入る。
単なる横断評価上の零性に置き換えているのではない。

第一項を分離切片に制限する。
$\operatorname{Hess}_{F_f}(q)$ の最初の二行・二列は零であり、
残りのブロックは $\varepsilon\operatorname{Hess}_f(p)$ である。
前命題からその横断三次係数は
\[
 \alpha_{dqY}(f)\alpha_{dqZ}(f)\,
       p^{\mathsf T}\operatorname{Hess}_f(p)p
 =6\alpha_{dqY}(f)\alpha_{dqZ}(f)f(p)
\]
となる。従ってHessian共変量の係数は零である。
\end{proof}

\begin{corollary}[二十四次の全二階補正の通常像所属]
\label{cor:v17-twentyfourth-corrections-conditional}
二十四次の四つの二階補正族について
\[
 P''(\mathcal G_7,\mathcal G_{13}),\quad
 P''(\mathcal G_8,\mathcal G_{12}),\quad
 P''(\mathcal G_9,\mathcal G_{11}),\quad
 P''(\mathcal G_{10},\mathcal G_{10})
 \ \subset J_{24}.
\]
これは四族の代表を個別の評価点で全て消すことを仮定しない。
\end{corollary}

\begin{proof}
命題\ref{prop:v17-ordinary-contact-rank24}により、
通常二十四次商は $[\Theta_8U_4^4],[\Theta_8U_4J_6^2]$ の二つで張られる。
従って $\mathcal K_{24}^{\mathrm{nor}}=0$ であり、
各組は $r+s+4=24$ を満たすので前定理を適用する。
この次数の核の消失は、後続節の明示像の独立性と
既存のnormal-slice単射を合わせた結果であり、
本節の定性的な分裂だけから推論したものではない。
\end{proof}

\begin{remark}[証明と有限監査の範囲]
Heisenberg二生成元の共通commutantの方程式を
$\mathbb Q(\omega)$ から有理数の方程式へ展開すると階数8であり、
核は単位行列の一次元空間となる。
\nolinkurl{theory/v17_structure_certificate.py} はこの有理行列と
優勢性の十行十列行列・行列式118098を独立に記録する。
三元共変量加群の全生成は有限標本から推測せず、明示したFisherの定理を入力とする。
全通常商の上界、$\mathcal K^{\mathrm{nor}}$ の消失、
実際のTate係数の消滅をこの監査だけから結論してはいない。
\end{remark}

% END theory/v17_normal_projection_ja.tex


% BEGIN theory/v17_discriminant_torsion_ja.tex
% 第17版：rank-four support、Schur補チャート、normal核の判別式torsion。
% 固定したsmooth Fの全Jacobian商や、原点での正則な単元変更とは区別する。

\section{横断射影の核は三元判別式上のtorsionである}
\label{sec:v17-discriminant-torsion}

本節では、前節で残した $\mathcal K^{\mathrm{nor}}$ の支持をさらに限定する。
普遍係数環を $R=\mathbb C[F,a]$、二つのイデアルを
\[
 I=(q_0,\ldots,q_4),\qquad G=(g_0,\ldots,g_4),\quad
 g_i=(\partial_iF)(q(a))
\]
と書く。通常商は有限同変成分を取った $R/G$ と一致する。
実際、$C\in\mathcal C_n$ が $G$ に入るなら、二重次数をそろえ、
完全可約性により同変な係数 $w\in\mathcal D_{n-4}$ を選べる。
全次数 $dq$ 全射から $C=dF(w)\in J_n$ を得る。
逆包含は定義から従う。

\subsection{通常共変量の類はPl\"ucker錐に支持される}

\begin{lemma}[階数四開集合での通常像所属]
\label{lem:v17-rank-four-support}
全ての $C\in\mathcal C_n$ と各 $i$ について
\[
 C\in G R_{q_i}.
\]
従って、$R/G$ における $C$ の類は $I$-torsionである。
すなわち、$C$ に依存する整数 $m$ が存在して
$I^mC\subset G$ となる。
\end{lemma}

\begin{proof}
階数四の点
$a_*=e_0\wedge e_1+e_2\wedge e_3$ を取る。
一パラメータ部分群
\[
 \mu(t)=\operatorname{diag}(t,t,t,t,t^{-2})
\]
では $\det\mu=t^2$ であり、$W=\Lambda^2V\otimes\Delta^{-1}$ の
$a_*$ は固定される。$q(a_*)=e_4^*$（体積線を自明化した表示）である。
$F$ のテンソル単項式の重みは、添字4の個数を $k$ として $3-3k$ となる。
$g(F,a_*)=0$ は $f_{i44}=0$（$i<4$）および $f_{444}=0$ を意味し、
負の重み $-3,-6$ の全成分を消す。
残る $F$ の重みは3と0であるのに対し、出力 $L$ の重みは $-6$ である。
従って多項式同変性により $C(F,a_*)=0$ となる。
階数四の点は一つの $\mathrm{GL}(V)$ 軌道なので、
$C$ は $q\ne0,\ g=0$ の全点で零である。

$q_i\ne0$ の開集合では、$F\mapsto g(F,a)$ は五成分全射である。
これらは $F$ の係数に線形な方程式なので、零点スキームは
$a$ の開集合上のベクトル束の全空間であり、被約である。
例えば対称テンソル係数 $f_{iii},f_{jii}$（$j\ne i$）を選べば、
対応する五行五列係数行列は、非零定数倍の $q_i^{10}$ を行列式に持つ。
従って零点上の消失から $C\in GR_{q_i}$ が従う。
各 $i$ についてある $m_i$ が $q_i^{m_i}C\in G$ を満たす。
$m=\sum_i(m_i-1)+1$ とすれば、$I^m$ の各単項式に
ある $q_i^{m_i}$ が因数として現れ、最後の主張を得る。
\end{proof}

\begin{remark}[支持を意味する範囲]
全ての $R/G$ が $I$-torsionだとは述べていない。
例えば $q\ne0$ にも $g=0$ の普遍incidenceは存在する。
上の結論は、指定した値域 $L$ の普遍同変成分を取ったときの消失である。
また $I$-torsion性だけから $IC\subset G$ や $C\in J_n$ は従わない。
\end{remark}

\subsection{Schur補チャートではPfaffianが正確に線形になる}

$x=a_{01}$ を可逆にし、$i,j\in\{2,3,4\}$ について
\[
 l_i=a_{0i},\qquad m_i=a_{1i},\qquad
 b_{ij}=a_{ij}-\frac{l_im_j-l_jm_i}{x}
\]
と置く。これらは $R_x$ の座標変換であり、
$I R_x=(b_{23},b_{24},b_{34})$ となる。
$p=(b_{34},-b_{24},b_{23})$ とすれば
\begin{equation}\label{eq:v17-schur-linear-pfaffian}
 q_0=\sum_{i=2}^4m_ip_i,\qquad
 q_1=-\sum_{i=2}^4l_ip_i,\qquad
 q_i=xp_i\quad(i=2,3,4).
\end{equation}
これは初期項ではなく正確な恒等式である。

行列 $h$ を
\[
 he_0=e_0-\sum_{i=2}^4\frac{m_i}{x}e_i,\qquad
 he_1=e_1+\sum_{i=2}^4\frac{l_i}{x}e_i,\qquad
 he_i=e_i\quad(i=2,3,4)
\]
で定義すると $\det h=1$ であり、
$a=h(xe_0\wedge e_1+\sum b_{ij}e_i\wedge e_j)$ である。
$F'=h^{-1}F$ と置き、その三元部分を $f=F'|_{Q_3}$ とすると
\begin{equation}\label{eq:v17-schur-exact-ternary}
 P(F,a)=x^3f(p),\qquad
 x^3\partial_i f(p)=m_ig_0-l_ig_1+xg_i.
\end{equation}
従って、三つの三元偏微分は $G R_x$ に入る。

\begin{lemma}[滑らかな三元三次式では横断四次以上を消せる]
\label{lem:v17-smooth-ternary-fourth-power}
上のチャートの係数環で、三元式 $f$ の判別式を可逆にすると
\begin{equation}\label{eq:v17-normal-fourth-in-gradient}
 (b_{23},b_{24},b_{34})^4\subset G
\end{equation}
が成立する。
\end{lemma}

\begin{proof}
まず固定した滑らかな三元三次式 $f$ を考える。
Euler恒等式から、その三つの偏微分の射影共通零点は曲線の特異点に等しく、
したがって存在しない。アフィン共通零点は原点だけなので、
$\mathbb C[p_2,p_3,p_4]$ で三つの二次式はheight3の正則列となる。
ここでも斉次極大での局所判定と斉次性を用いて大域へ戻せる。
三つの二次式のKoszul自由分解からHilbert級数は
\[
 \frac{(1-t^2)^3}{(1-t)^3}=(1+t)^3
\]
である。従って四次成分は零で、全ての四次単項式が三偏微分のイデアルに入る。

次に $F'$ と $x^{\pm1},l_i,m_i$ を変数とする係数環を $B$ とする。
$R_x=B[p_2,p_3,p_4]$ であり、上の三元式の係数も $B$ に属する。
判別式を可逆にした $B$ 上で、三偏微分と二次単項式の積による写像
\[
 B^{18}\longrightarrow B^{15}
\]
を考える。前段の議論から、この写像は各閉点の剰余体上で全射である。
有限生成の余核が非零ならその支持は閉点を持ち、
Nakayamaの補題に反する。従って係数環上でも全射である。
全ての高い単項式は四次単項式の倍数なので、
式\eqref{eq:v17-schur-exact-ternary}から主張を得る。
無限回の $I$-進近似や完備化による持上げは用いていない。
\end{proof}

\subsection{未知の通常商成分の判別式torsion}

\begin{lemma}[正規化済み単元からの三元判別式]
\label{lem:v17-unit-discriminant}
\[
 \mathscr D=J_6^2-2304U_4^3\in\mathcal U_{12}
\]
と置く。$\overline{\mathscr D}\in\mathcal A_{12}$ は三元三次式の通常の判別式の
非零定数倍である。特に $\overline{\mathscr D}(f)\ne0$ と
$f$ が滑らかであることは同値である。
\end{lemma}

\begin{proof}
Fisherの不変式定理\cite[Theorem~4.4]{Fisher2008}では、
三元判別式は次数12であり、その非消失は滑らかさと同値である。
既知の制限環 $\mathcal A$ の12次成分は
$\overline U_4^3,\overline J_6^2$ の二つで張られる。
第14版のHesse制限では
\[
 \overline U_4(f_t)=t^4-t,\qquad
 \overline J_6(f_t)=6-120t^3-48t^6.
\]
$t=0$ は滑らかであり、$t=-1/2$ の
$x^3+y^3+z^3-3xyz$ は三直線の積である。
後者での値は $U_4=9/16$、$J_6=81/4$ で、$J_6^2/U_4^3=2304$ となる。
従って判別式の二つの係数の比が決まり、
滑らかな $t=0$ での非消失により定数倍は非零である。
係数の追加検算として
\[
 \overline{\mathscr D}(f_t)=36(1+8t^3)^3
\]
も成り立つ。
\end{proof}

\begin{theorem}[横断射影は判別式可逆化後に同型となる]
\label{thm:v17-normal-kernel-discriminant-torsion}
$\mathcal K^{\mathrm{nor}}=\ker\pi$ の任意の元 $\xi$ に対して、
ある整数 $m=m(\xi)\ge0$ が存在して
\begin{equation}\label{eq:v17-normal-kernel-discriminant-torsion}
 \mathscr D^m\xi=0
\end{equation}
となる。従って
\begin{equation}\label{eq:v17-discriminant-localized-quotient}
 \left(\bigoplus_{n\ge6}\mathcal C_n/J_n\right)[\mathscr D^{-1}]
 \simeq
 \mathcal A[\overline{\mathscr D}^{-1}]\,[\Theta_8]
\end{equation}
である。この局所化は次数を保ち、$\mathscr D^{-1}$ の次数を $-12$ とする。
\end{theorem}

\begin{proof}
元は有限個の斉次成分の和なので、各成分について示せば十分である。
$\xi$ を横断四次以上の $C\in\mathcal C_n^{\ge4}$ で表す。
補題\ref{lem:v17-rank-four-support}から、その $R/G$ における類は $I$-torsionである。
従って、$I$ を含まない素イデアルで局所化すれば零になる。

$I$ を含み $\mathscr D$ を含まない素イデアル $\mathfrak p$ を考える。
$\mathscr D$ の $a$ 次数は12なので、$\mathfrak p$ は全ての $a_{ij}$ を含むことはできない。
ある $a_{ij}$ を可逆にしたSchur補チャートで論じられる。
添字を替えて $x=a_{01}$ が可逆であるとしてよい。
錐上の点での同変性と $a$ 次数により
\[
 \mathscr D(F,a)\equiv x^{12}\overline{\mathscr D}(f)\pmod I.
\]
従って $\overline{\mathscr D}(f)$ は $R_{\mathfrak p}$ で可逆である。

横断四次以上という条件は、このチャートで
$C\in(b_{23},b_{24},b_{34})^4$ を意味する。
実際、$\det h=1$ と $a$ に関する斉次性により
\[
 C(F,a)=C(F',xe_0\wedge e_1+b)
       =x^n C(F',e_0\wedge e_1+b/x)
\]
であり、定義から四次未満の $b$ 係数は全て零である。
補題\ref{lem:v17-smooth-ternary-fourth-power}を適用すると、
$C\in GR_{\mathfrak p}$ となる。
これで $\mathscr D$ を可逆にした環の全素点で $C$ の類が零となり、
$C\in GR_{\mathscr D}$ を得る。
局所化の定義からある $m$ に対して $\mathscr D^mC\in G$ である。
同変な有限次数成分を取り、全次数 $dq$ 全射を使えば
$\mathscr D^mC\in J_{n+12m}$ となる。
直和分解\eqref{eq:v17-normal-splitting}を局所化して最後の同型を得る。
\end{proof}

\begin{remark}[局所化で見えなくなる部分は依然として未決定]
この定理は $\mathcal K^{\mathrm{nor}}=0$ を意味しない。
既知の $[H]$ も $\mathscr D[H]=0$ を満たす非零元である。
ほかの判別式torsionが高次に現れないかは、引き続き全商の生成問題として残る。
また $\mathscr D$ は原点 $a=0$ で零なので、
この局所化を原点で許される定数項1の単元変更と解釈してはならない。
$\mathscr D^mR_n$ が通常像に入ることだけでは、実際の $R_n$ を同じ次数で消去できない。
本定理は全ての次数で使える構造上の制限であり、24次商の上界2の代わりではない。
\end{remark}

% END theory/v17_discriminant_torsion_ja.tex


% BEGIN theory/v17_vertex_torsion_ja.tex
% 第17版：Cech--Koszulによるnormal核のvertex支持と有限次元性。
% 有限補正の全次元・終端はまだ計算しない。

\section{通常商の未知部分は頂点に支持される有限次元加群である}
\label{sec:v17-vertex-finiteness}

前節では未知の部分 $\mathcal K^{\mathrm{nor}}$ が判別式torsionであることを示した。
さらに普遍勾配のKoszul複体と横断変数の\v{C}ech複体を同時に使うと、
この部分は $a=0$ の頂点に支持され、有限次元であることまで分かる。
したがって十分高い次数では通常商は $\Theta_8$ の自由部分加群と一致する。
ただし、その「十分高い」次数の明示的な上限や、有限補正の全体は本節では決定しない。

\subsection{横断三変数での局所コホモロジーの次数上界}

\begin{lemma}[\v{C}ech--Koszulの零次端]
\label{lem:v17-cech-koszul-normal-degree-bound}
$B$ をNoether環、$S=B[p_1,p_2,p_3]$、
$\mathfrak m=(p_1,p_2,p_3)$ とする。
$g_1,\ldots,g_5$ は $p$ に関して二次斉次とし、
$K_\bullet=K(g_1,\ldots,g_5;S)$ と置く。
正次数Koszulホモロジー $H_i(K_\bullet)$ が全て
$\mathfrak m$-torsionであると仮定する。
このとき
\begin{equation}\label{eq:v17-local-torsion-degree-bound}
 \bigl(H^0_{\mathfrak m}(S/(g_1,\ldots,g_5))\bigr)_d=0
 \qquad(d>3)
\end{equation}
である。下付き $d$ は $p$ の次数であり、係数環 $B$ の次数ではない。
\end{lemma}

\begin{proof}
拡大\v{C}ech複体を $\check C^\bullet_{\mathfrak m}$ と書き、
$K_i$ をコホモロジー次数 $-i$ に置いて
\[
 \operatorname{Tot}
  \bigl(\check C^\bullet_{\mathfrak m}\otimes_S K_\bullet\bigr)
\]
を考える。両方向に有限な双複体であり、二つのスペクトル系列は収束する
\cite[Tag 012X, Lemmas 12.25.1 and 12.25.3]{StacksCechKoszul}.
局所化は平坦なので、先にKoszulホモロジーを取った系列の第二項は
\[
 E_2^{u,v}=H^u_{\mathfrak m}(H_{-v}(K_\bullet)),
 \qquad u\ge0,\quad v\le0,
\]
総次数は $u+v$ である。
$i>0$ の $H_i(K_\bullet)$ は $\mathfrak m$-torsionなので、
その\v{C}ech複体は次数0以外で零となる。
従って総次数0の非零になり得る項は
\[
 E_2^{0,0}=H^0_{\mathfrak m}(S/(g))
\]
だけである。$r\ge2$ の微分の出先は $(r,1-r)$ であり、
正次数Koszulホモロジーの正次数\v{C}echコホモロジーだから零である。
入って来る項は $u=-r<0$ なので存在しない。
よって総複体の零次コホモロジーはこの $E_2^{0,0}$ と同型である。

逆に、先に\v{C}echコホモロジーを取る。
$p_1,p_2,p_3$ は多項式環の正則変数なので、自由 $S$ 加群の
\v{C}echコホモロジーは次数3以外で零である
\cite[Tag 0G6L]{StacksCechKoszul}。
ここではさらに
\[
 H^3_{\mathfrak m}(S)
  =\bigoplus_{\alpha_1,\alpha_2,\alpha_3\ge1}
       B\,p_1^{-\alpha_1}p_2^{-\alpha_2}p_3^{-\alpha_3}
\]
と明示でき、その最高 $p$ 次数は $-3$ である。
従って同じ総複体の零次コホモロジーは
\[
 H_3\bigl(H^3_{\mathfrak m}(K_\bullet)\bigr)
\]
である。$K_3=S(-6)^{10}$ なので、その第三項
$H^3_{\mathfrak m}(K_3)$ の最高 $p$ 次数は $6-3=3$ である。
ホモロジーはこの項の部分商だから、次数 $d>3$ では零となる。
両順序の計算を比較して結論を得る。
\end{proof}

\subsection{普遍勾配への適用}

\begin{theorem}[横断射影の核は頂点に支持される]
\label{thm:v17-normal-kernel-vertex-support}
$R=\mathbb C[F,a]$、$G=(g_0,\ldots,g_4)$、
$\mathfrak a=(a_{ij}:i<j)$ とする。
任意の $\xi\in\mathcal K^{\mathrm{nor}}$ は $R/G$ において
$\mathfrak a$-torsionである。
すなわち、$a=0$ の頂点を除く全ての点でその類は零である。
\end{theorem}

\begin{proof}
斉次な $\xi$ を、命題\ref{prop:v17-normal-invisible-representatives}による
$C\in\mathcal C_n^{\ge4}$ で表す。一般の元は有限個の斉次成分の和として扱えばよい。
補題\ref{lem:v17-rank-four-support}により、その類は $I$-torsionである。
各 $a_{ij}$ を可逆にしたSchur補チャートを考え、
添字を替えて $x=a_{01}$ の場合で示す。

前節の変数変更により $R_x=B[p_2,p_3,p_4]$ と書ける。
五つの勾配の各成分は $p$ の二次斉次式である。
$p\ne0$ では $q\ne0$ であり、$F$ から五つの勾配への写像は全射である。
従って五つの勾配はこの開集合上で、$F$ の五つの独立な線形方程式として
正則列をなす。よって正次数Koszulホモロジーは
$\mathfrak m=(p_2,p_3,p_4)$ に支持され、有限生成性から
$\mathfrak m$-torsionである。
ここでは五つの勾配が $p=0$ でも正則列であるとは仮定していない。

$C$ の局所化は $\mathfrak m$-torsionであり、
\[
 C(F,a)=x^nC(F',e_0\wedge e_1+b/x)
\]
により $p$ の次数が4以上の項だけからなる。
勾配イデアルはこの次数に関して斉次なので、
各斉次成分も $\mathfrak m$-torsionである。
前補題によりその全てが零であり、$C\in GR_x$ となる。
全ての $a_{ij}$ チャートについて同じことが成り立つから、
頂点以外の全開集合で $\xi$ は零である。
局所化の定義から、ある $\ell$ で $\mathfrak a^\ell\xi=0$ となる。
\end{proof}

\begin{theorem}[未知部分の有限次元性と十分高い次数の通常商]
\label{thm:v17-finite-normal-kernel}
$\mathcal K^{\mathrm{nor}}$ は有限次元の次数付き複素ベクトル空間である。
従って、ある整数 $n_0$ が存在し、全ての $n\ge n_0$ について
\begin{equation}\label{eq:v17-eventual-contact-quotient}
 \mathcal C_n/J_n
  =\bigoplus_{\substack{a,b\ge0\\4a+6b=n-8}}
      \mathbb C[\Theta_8U_4^aJ_6^b]
\end{equation}
となる。特にHilbert級数は
\begin{equation}\label{eq:v17-contact-hilbert-finite-correction}
 \sum_{n\ge6}\dim(\mathcal C_n/J_n)t^n
   =\frac{t^8}{(1-t^4)(1-t^6)}+E_{\mathrm{vert}}(t),
 \qquad E_{\mathrm{vert}}(t)\in\mathbb Z_{\ge0}[t]
\end{equation}
と書ける。$E_{\mathrm{vert}}(t)$ は有限な多項式で、
既知の十次類によりその $t^{10}$ 係数は1である。
$E_{\mathrm{vert}}(t)=t^{10}$ は定理\ref{thm:pf-full-contact-module}で決定する。
\end{theorem}

\begin{proof}
$T_0=H^0_{\mathfrak a}(R/G)$ は、Noether加群 $R/G$ の部分加群なので有限生成である。
従って有限個の二重斉次生成元 $v_1,\ldots,v_e$ と整数 $\ell$ を選び、
\[
 \mathfrak a^\ell T_0=0
\]
とできる。$v_j$ の $a$ 次数を $b_j$ とすると、
$T_0$ の $a$ 次数は $\max_j b_j+\ell-1$ 以下に限られる。
実際、任意の元を生成元の $F,a$ 多項式倍で表し、
$a$ 次数が $\ell$ 以上の係数単項式を除けばよい。

前定理から $\mathcal K_n^{\mathrm{nor}}$ は $T_0$ の指定された同変成分に入り、
その二重次数は $(n-5,n)$ である。
従って有限個の $n$ についてしか非零にならない。
各 $n$ では $F$ 次数も固定されているので、その空間は有限次元である。
ここで $T_0$ 全体が複素ベクトル空間として有限次元であるとは述べていない。
必要なのは $a$ 次数の有界性と、この対角的な二重次数条件の組合せである。

分裂\eqref{eq:v17-normal-splitting}から十分高い次数の等式とHilbert級数を得る。
十次では $\mathcal A_2=0$ であり、
定理\ref{thm:boundary-eleven-jet}の通常商が既知の $[H]$ 一方向なので、
$E_{\mathrm{vert}}$ の該当係数は1である。
\end{proof}

\begin{corollary}[未知部分は正次数単元についての有限torsion]
\label{cor:v17-positive-unit-torsion}
$\mathcal U_+=\bigoplus_{s>0}\mathcal U_s$ とすると、
通常商の $\mathcal U_+$-torsion部分は正確に $\mathcal K^{\mathrm{nor}}$ である。
ここでtorsionとは、ある $k$ に対して $(\mathcal U_+)^k\xi=0$ となるという意味である。
単一の正次数単元による消失とは区別する。
その全体を消す $\mathcal U_+$ の冪が存在するが、この冪の最小値は未計算である。
\end{corollary}

\begin{proof}
有限個の次数しか持たない $\mathcal K^{\mathrm{nor}}$ は、
十分大きい正次数単元の積で消える。
一方、$\Theta_8\mathcal A$ では $\overline U_4$ のどの冪も非零な元を消さない。
$\mathcal A$ が整域であることと分裂から結論を得る。
\end{proof}

\begin{remark}[有効な有限計算の範囲はまだ別に必要である]
本節は、十分高い次数での通常商を定性的に決める。
しかし、閾値 $n_0$、$E_{\mathrm{vert}}$ の最高次数と係数、
および $\mathcal K^{\mathrm{nor}}=\mathbb C\cdot[H]$ かどうかは決めていない。
従って、個別の24次上界を本節だけから得たとは述べない。
また実際のTate係数が自由部分へ零に写ることも、別の計算または幾何学的論証を要する。
\end{remark}

% END theory/v17_vertex_torsion_ja.tex


% BEGIN characters/v17_character_notes_ja.tex
\section{二十四次の普遍共変量とKoszul各項}
\label{sec:v17-characters}

二十四次では、normal sliceから得られる二つの商の元を出発点とし、
普遍空間の次元と接触写像の階数を別々に求める。
本節の計算は整数の重み指標、Weyl交代和、Pieri則だけに依存する。

\subsection{次数とdeterminant twist}

$\Delta=\det V$、$A=\Lambda^2V$、$W=A\otimes\Delta^{-1}$ とし、
$F\in\operatorname{Sym}^3V$、$a\in W$、
$q(a)=\tfrac12a\wedge a\in V^*\otimes\Delta^{-1}$、
$P=F(q(a))\in\Delta^{-3}$ という規約を保つ。
中心元の重みから、$a$ 次数 $n$ の境界項は $F$ 次数 $n-5$、
$a$ 次数 $r$ の座標補正は $F$ 次数 $r-1$、
$a$ 次数 $s$ の単元補正は $F$ 次数 $s$ を持つ。
従って二十四次のKoszul第 $k$ 項は
\[
 K_{k,24}=\operatorname{Hom}_{\mathrm{GL}(V)}
 \left(\operatorname{Sym}^{19-k}(\operatorname{Sym}^3V)
 \otimes\operatorname{Sym}^{24-4k}A,
 \Lambda^kV^*\otimes\Delta^{21-2k}\right),\qquad 0\leq k\leq5.
\]
$K_{0,24}=C_{24}$、$K_{1,24}=D_{20}$、$K_{2,24}=B_{24}$ と書く。
第5項にも $a$ 次数4が残るので、終端の多重度は前の次数から引き継がない。

\begin{proposition}[二十四次の厳密次元]
\label{prop:v17-character-dimensions}
各普遍空間の次数・標的分割・次元は次の通りである。
\[
\begin{array}{c|r|r|c|r}
\text{空間}&\deg_F&\deg_a&\text{標的分割}&\dim\\\hline
C_{24}&19&24&(21,21,21,21,21)&2458\\
G_{19}&18&19&(19,19,18,18,18)&11676\\
U_{18}&18&18&(18,18,18,18,18)&1365\\
D_{20}&18&20&(19,19,19,19,18)&6262\\
B_{24}&17&16&(17,17,17,16,16)&5615\\
K_{3,24}&16&12&(15,15,14,14,14)&2079\\
K_{4,24}&15&8&(13,12,12,12,12)&282\\
K_{5,24}&14&4&(10,10,10,10,10)&11
\end{array}
\]
補助的な十七次の空間についても
\[
 \dim C_{17}=31,\qquad \dim G_{12}=155,\qquad
 \dim U_{11}=9,\qquad \dim D_{13}=92
\]
を得る。これらの値は接触写像の階数を仮定しない。
\end{proposition}

\begin{proof}
$\operatorname{Sym}^m(\operatorname{Sym}^3V)$ の指標を35個の三次単項式重み
$\alpha$ に対するEuler積
\[
 \prod_{|\alpha|=3}(1-tx^\alpha)^{-1}
\]
の $t^m$ 係数として計算する。
$\rho=(4,3,2,1,0)$ とするとSchur多重度は
\[
 \sum_{\sigma\in S_5}\operatorname{sgn}(\sigma)
 [x^{\lambda+\rho-\sigma\rho}]
 \operatorname{ch}\operatorname{Sym}^m(\operatorname{Sym}^3V)
\]
である。外積二乗側も10個の重みから指標を作り、
テンソル積の交代和を直接評価した。
さらにLittlewood分解
\[
 \operatorname{Sym}^{d}(\Lambda^2V)
 =\bigoplus_{j=0}^{\lfloor d/2\rfloor}S_{(d-j,d-j,j,j,0)}V
\]
の多重度1とWeyl次元和を照合した。

$C_{24}$ で必要な $F$ 側の分割は
$(21,18-j,18-j,j,j)$、$0\leq j\leq9$ であり、多重度は
\[
 (4,5,42,108,369,521,765,425,213,6).
\]
和は2458である。
$U_{18}$ の分割は $(18,18-j,18-j,j,j)$、$0\leq j\leq9$ で、多重度は
\[
 (2,1,16,44,179,275,447,261,136,4).
\]
和は1365である。

$G_{19}$ と $D_{20}$ はそれぞれ外積第3冪・第1冪を用いる垂直帯Pieri計算、
$K_{k,24}$ は外積第 $k$ 冪を用いるPieri計算でも計算し、
直接のテンソル積指標と照合した。
Koszul源の非零Pieri成分数は $k=2,3,4,5$ の順に29、24、13、3である。
第5項の $F$ 側分割と多重度は
\[
 (10,10,10,6,6):7,\qquad
 (10,9,9,7,7):2,\qquad
 (10,8,8,8,8):2
\]
となる。十七次の補助四空間も同じ二通りの計算で一致する。

全ての $F$ 次数 $1\leq m\leq19$ と $a$ 次数 $1\leq d\leq24$ について、
独立のNewton--Adams再帰
\[
 m h_m=\sum_{k=1}^{m}p_kh_{m-k}
\]
で全重み係数を再計算し、Euler積から得た辞書全体と一致することを確認した。
必要な $F$ 次数 $11,12,14,15,16,17,18,19$ の全Schur分解についても、
多重度とWeyl次元の積の和が $\binom{m+34}{m}$ に一致する。
この検証は浮動小数点階数を使わない有限の整数恒等式である。
\end{proof}

\subsection{Euler交代和と高次Koszul微分}

交代和は
\[
 2458-6262+5615-2079+282-11=3.
\]
前版の高次ホモロジー消失定理
\ref{thm:v16-koszul-high-homology-vanishing}を適用すると
\[
 \dim H_0-\dim H_1+\dim H_2=3,
 \qquad H_i=0\quad(i>2).
\]
従って上位三つの微分については
\[
 \operatorname{rank}\delta_5=11,\qquad
 \operatorname{rank}\delta_4=282-11=271,\qquad
 \operatorname{rank}\delta_3=2079-271=1808
\]
が確定する。一方、$H_1,H_2$ の消失はこの計算からは従わない。
通常商の次元は別の接触像の計算で求める必要がある。
Pfaffian微分の全次数全射性を用いれば
\[
 \dim\ker(\delta q:G_{19}\longrightarrow D_{20})=11676-6262=5414
\]
を得るが、$dF:D_{20}\to C_{24}$ の階数はこの差からは決まらない。

\begin{proposition}[Pl\"ucker錐への制限の標的次元]
\label{prop:v17-plucker-dimensions}
分解可能な $a$ の錐への制限で、$\operatorname{Sym}^dA$ をCartan成分
$S_{(d,d)}V$ に置き換えたHom空間の次元は
\[
\begin{array}{c|rrrrrrrr}
 &C_{24}&G_{19}&U_{18}&D_{20}&B_{24}&K_{3,24}&K_{4,24}&K_{5,24}\\\hline
 \dim&0&2&2&0&0&18&36&7
\end{array}
\]
である。また $U_{16}$ の制限標的は2次元である。
\end{proposition}

\begin{proof}
Littlewood分解の $j=0$ のみを残して同じPieri計算を行う。
$U_{18}$ では $F$ 側分割 $(18,18,18,0,0)$ の多重度2、
$U_{16}$ では $(16,16,16,0,0)$ の多重度2を得る。
$D_{20}$ は第一行20が比較する矩形の幅19を超えるため全成分が消える。
全成分を保存証明書の \texttt{Plucker} 欄に記録した。
\end{proof}

この表は制限先のHom空間の次元であり、特定の候補が通常商全体を尽くすことを意味しない。
二つの候補 $\Theta_8U_4^4$、$\Theta_8U_4J_6^2$ の通常商での独立性は
normal sliceの単射定理に依存し、それらが全商を尽くすかどうかは接触階数によって判定する。

Fermat固定群の成分別計算から得る評価階数の上界は
\[
\begin{array}{c|r|r|r}
 &\dim&\text{対角値の階数上界}&\text{対角値と一階法jetの階数上界}\\\hline
C_{24}&2458&621&1962\\
G_{19}&11676&2253&8319\\
U_{18}&1365&224&863\\
D_{20}&6262&1314&4594
\end{array}
\]
である。実際の評価階数との等号は主張しない。

\subsection{再現用証明書}

\texttt{characters/v17\_character\_certificate.py} が全重みEuler積・Newton再帰、
全Schur次元和、テンソル積の交代和とPieri照合を新規に実行する。
\texttt{v17\_koszul\_pieri\_check.py} は別実装の垂直帯列挙を使い、
保存Schur分解から全Koszul源を再計算する。
\texttt{v17\_character\_audit.py} は保存証明書のtwist、交代和、Weyl次元、
Pieri、補助十七次、制限上界を監査する。
保存監査と全重みの新規再計算は別の操作として記録する。

% END characters/v17_character_notes_ja.tex


% BEGIN rank/rank24_section_ja.tex
\section{二十四次の通常接触商の厳密決定}
\label{sec:v17-contact-rank}

二十四次の計算では、19頂点の新しい場だけを列挙する代わりに、
下位接触像の単元倍と既知のHessian作用の全次数定理を併用する。
全ての追加列が通常像に属することを先に保証し、
二階Taylor補正の未証明な所属を階数の下界へ混ぜない。

\begin{proposition}[二十四次の通常像と二次元商]
\label{prop:v17-ordinary-contact-rank24}
二十四次では
\[
 \dim J_{24}=2456,\qquad
 \dim(C_{24}/J_{24})=2,
\]
かつ
\[
 C_{24}/J_{24}
 =\mathbb C\cdot[\Theta_8U_4^4]
   \oplus\mathbb C\cdot[\Theta_8U_4J_6^2].
\]
従って
\[
 \dim\ker(dP:G_{19}\to C_{24})=9220,\qquad
 \dim\ker(dF:D_{20}\to C_{24})=3806.
\]
\end{proposition}

\begin{proof}
命題\ref{prop:v17-character-dimensions}より
$\dim C_{24}=2458$、$\dim U_{18}=1365$ である。
普遍係数環が整域で $P\neq0$ なので、$u\mapsto Pu$ は単射であり、
$\dim(PU_{18})=1365$ となる。
以下では $P=0$ 上の評価を用いて、この部分空間を除いた像の下界を求める。

まず、既存の低次の場 $\Gamma$ と単元グラフ $u$ に対する
$u\,dP(\Gamma)=dP(u\Gamma)$ を作る。
次数の組は $(\deg_a\Gamma,\deg_a u)=(15,4),(13,6),(11,8),$
$(9,10),(7,12),(10,9),(8,11)$ である。
629個の候補から502個の独立な列を選んだ。
この積は、単元グラフと場のグラフの互いに素な和を19頂点のグラフとして書いたものでもある。
四頂点の生の単元グラフは $24U_4$ であり、
これを $U_4$ と取り違えない。像の基底の構成では生のグラフを使う。

次に、定理\ref{thm:v16-positive-unit-hessian-annihilation}により
$HU_{14}\subset J_{24}$、系\ref{cor:v16-all-hessian-first-actions}により
$dH(G_{15})\subset J_{24}$ である。
これらから、それぞれ97列、241列を追加する。
最後に、19頂点の連結グラフで定義される $G_{19}$ の場から251列を追加した。
開いた交代頂点への二重辺も探索対象に含めた。
採用された列数の内訳は
\[
 \begin{array}{c|r}
 \text{通常像への所属が保証された列}&\text{採用数}\\\hline
 u\,dP(\Gamma)\text{ の下位積}&502\\
 H\,u,\ u\in U_{14}&97\\
 dH(\Gamma),\ \Gamma\in G_{15}&241\\
 dP(\Gamma),\ \Gamma\in G_{19}\text{ の連結グラフ}&251\\\hline
 \text{計}&1091
 \end{array}
\]
となる。各列は普遍的な多項式であり、Hessianの列も拡大した有効像ではなく通常像に属する。

評価点は1120個の有理点で、各点で $q(a)=(1,1,1,1,1)$、$P=0$ が有理数上で成立する。
選んだ1091行と1091列の行列式は
\[
 \det M_{24}\equiv579\pmod{1009}
\]
である。全分母は1009と互いに素であり、従ってこの有理行列式は非零である。
全ての $PU_{18}$ の元は評価点で消えるので、1091列は $PU_{18}$ を法として独立であり、
\[
 \dim J_{24}\ge1365+1091=2456
\]
を得る。

他方、normal sliceの単射定理\ref{thm:v16-hessian-socle-injection}より
$[\Theta_8U_4^4]$ と $[\Theta_8U_4J_6^2]$ は通常商で独立である。
従って $\dim J_{24}\le2458-2=2456$ であり、両側の評価が一致する。
$PU_{18}$ はEuler場で $dP(G_{19})$ に吸収でき、
Pfaffian微分 $\delta q:G_{19}\to D_{20}$ は全射なので、
二つの核の次元はそれぞれ $11676-2456=9220$、$6262-2456=3806$ となる。
\end{proof}

\subsection{独立監査と補元の有限証明書}

探索側は左の頂点で辺を番号付けし、サンプル軸を持つ整数行列積で縮約する。
独立監査側は右の頂点で番号付けし直し、元の貪欲な縮約経路と
\texttt{einsum} による別の整数縮約を用いる。
19頂点のグラフは57個の辺添字を持つため、各二項縮約で添字を局所的に振り直す。
各加算の前には
\[
 (1009-1)^2\,5^s<2^{63}
\]
を確認する。ここで $s$ はその二項縮約の総和添字の個数である。
有限体への剰余を各段階で取り、浮動小数点計算は使わない。

非連結グラフの成分は実際の二部グラフから分解し、
同じ成分の値は一度だけ独立縮約する。
905個の相異なる因子から、採用された小行列の
\[
 1091^2=1190281
\]
成分を全て復元して探索側と照合する。
Hessianの入力微分は探索側では余因子を使い、独立側では
$h(q+te_i)$ の次数5の補間から求める。
また有限体の行列式は、配列演算による消去とPython整数による別の消去で照合する。
最終独立監査では1190281成分が全て一致し、二つの消去はいずれも579を返した。
成分分解はグラフの定義から従う因数分解であり、行列の一部の検査を省略する操作ではない。

二つの補元の値は全1120点で有理数上に新規計算した。
$U_4$ は四頂点グラフを24で割ったもの、$J_6$ は既存の固定した六頂点グラフを使う。
補元を追加した1093次の小行列式は
\[
 \det(M_{24},W_{24,1},W_{24,2})\equiv56\pmod{1009}
\]
である。この有限行列は明示した部分空間からの独立性を再確認する。
全通常像に対する非所属は、有限体で列が増えたことだけからではなく、
前の証明で引用したnormal sliceの定理により保証される。

探索が停止したこと、有限体の値が零であったこと、あるいはEuler交代和は、
いずれも通常像の上界には使っていない。

\begin{corollary}[二十四次のnormal核とKoszulホモロジー]
\label{cor:v17-rank24-normal-koszul}
normal射影の直和分解\eqref{eq:v17-normal-splitting}に現れる核は
\[
 \mathcal K_{24}^{\mathrm{nor}}=0
\]
である。従って二十四次ではnormal三次jetの二つのスカラー係数が通常商の完全な検出子となる。
また $b=\operatorname{rank}\delta_2$ と書けば、
\[
 \dim H_1=3806-b,\qquad
 \dim H_2=3807-b,\qquad
 \dim H_2=\dim H_1+1\ge1.
\]
\end{corollary}

\begin{proof}
通常商の次元2と、直和因子
$\mathbb C\cdot[\Theta_8U_4^4]\oplus\mathbb C\cdot[\Theta_8U_4J_6^2]$
の次元2が一致するため、残るnormal核は零である。
Koszul第1微分の核の次元は3806であり、
第3微分の階数は前節の高次ホモロジー消失から1808である。
従って $\dim H_1=3806-b$、
$\dim H_2=(5615-b)-1808=3807-b$ となる。
\end{proof}

これは二十四次でもKoszul複体が正次数で完全でないことを意味する。
第2微分の階数 $b$ と $H_1,H_2$ の個別の次元は本計算では確定していない。

\subsection{再現の範囲}

\texttt{rank/contact24\_combined.json} は、全入力、各列の型とグラフ、
有限体値、採用行を収録する。
\texttt{v17\_independent\_modular\_audit.py} が採用小行列の全成分を新規に独立縮約し、
\texttt{contact24\_combined\_independent\_audit.json} に結果と実行脚本のhashを保存する。
\texttt{v17\_rank\_certificate.py} は保存された独立因子から全成分を再構成し、
二つの小行列式と入力・脚本のhashを監査する。
保存監査は大きなテンソル縮約を新規に実行する操作とは区別する。

% END rank/rank24_section_ja.tex


% BEGIN contact/contact_twenty_fourth_ja.tex
% 第17版：24次の厳密な有限台帳と四つの二階族。
% 本節単独では通常24次商の次元も、四族の像所属も仮定しない。

\section{二十四次の接触作用と低次補正の四族への縮約}
\label{sec:v17-contact-twenty-fourth}

$P=P_F=F(q(a))$、$H=\det\operatorname{Hess}_F(q(a))$、$c=1/96$、
$J_n=dP(\mathcal G_{n-5})$ とする。
第16版の定理\ref{thm:v16-twenty-fourth-normal-form}により、
実際の境界には $P+cH+\sum_{n\ge24,\ n\ {\rm even}}R_n$ という偶正規形がある。
本節では未正規化Tate係数 $f_n=d_n^{\mathrm T}$ からその24次類へ至る
補正を正確に追跡する。通常24次商が二次元であることも、
以下の二階補正が零であることも、この段階では用いない。

\subsection{二十四次の対数作用と十七次の補助基底}

\begin{proposition}[二十四次の保存条件と有効像]
\label{prop:v17-effective-twenty-fourth}
最小増分4の接触対数 $L=\sum_{k\ge4}Q_k$、
$Q_k=Y_{k+1}+b_k$ に対し、$P+cH$ の23次までを保存する条件は
\begin{equation}\label{eq:v17-twenty-fourth-preservation}
 Q_kP+cQ_{k-4}H=0\quad(4\le k\le17),\qquad Q_j=0\quad(j<4)
\end{equation}
である。24次作用は $Q_{18}P+cQ_{14}H$ であり、
\[
                            I_{24}=J_{24}.
\]
全普遍接触変換と、奇座標・偶単元に制限した変換は同じ24次像を与える。
\end{proposition}

\begin{proof}
補題\ref{lem:v14-contact-logarithm}により、
差が24次以上である条件は $L(P+cH)$ が24次以上である条件と同値で、
24次では指数作用と対数作用が一致する。
当該作用に関わる偶数増分の鎖は
\[
 Q_6P=0,\qquad Q_{10}P+cQ_6H=0,\qquad
 Q_{14}P+cQ_{10}H=0.
\]
$Q_4,Q_8,Q_{12},Q_{16}$ の鎖と奇数増分は独立なので零に選べる。
一方、構造定理\ref{thm:v16-positive-unit-hessian-annihilation}と
系\ref{cor:v16-all-hessian-first-actions}により
\[
 Q_{14}H\in dH(\mathcal G_{15})+H\mathcal U_{14}\subset J_{24}.
\]
$Q_{18}P$ の像はEuler場を含めて $J_{24}$ であり、
自由な19次場が逆包含を与える。
\end{proof}

\begin{proposition}[十七次通常像と二十二個の代表場]
\label{prop:v17-seventeenth-contact-basis}
\[
 \dim\mathcal C_{17}=31,\quad
 \dim\mathcal G_{12}=155,\quad
 \dim\mathcal U_{11}=9,\quad
 \dim\mathcal D_{13}=92,\qquad
                           J_{17}=\mathcal C_{17}.
\]
特に $J_{17}/P\mathcal U_{11}$ は22次元である。
公開する22個のグラフ場 $\Gamma_{12,k}$ の像はその基底となる。
$dP$ と $dF$ の核の次元は、それぞれ124と61である。
\end{proposition}

\begin{proof}
四つの源の次元は
\nolinkurl{v17_character_certificate.json} の
\texttt{auxiliary\_degree17} に記録した全重み指標と独立Pieri計算による。
$P=0$ の一般整数入力で22グラフの像を有理数上で再縮約すると、
公開した $22\times22$ 行列の行列式は非零で、1009を法とする値は529である。
この行列と有理数Gaussian消去を
\nolinkurl{v17_g12_graph_basis.py} とそのJSONで再現する。
従って22方向は $P\mathcal U_{11}$ を法として独立である。
$P\mathcal U_{11}$ は9次元で $J_{17}$ に含まれるので、
$\dim J_{17}\ge22+9=31=\dim\mathcal C_{17}$ となる。
全次数Pfaffian全射により $dF$ も同じ像を持つため、核の次元も従う。
\end{proof}

\subsection{merged変更の十三項と純粋な有限消去}

前版と同じ $X_7,X_8$ および $U_4,V_6$ を用い、
\[
               g=(1+3U_4-V_6)f(a+X_7+X_8)
\]
とする。単元を $(1+3U_4)(1-V_6)$ に置き換えない。
また $X_7,X_8$ は一つの座標変更の和であり、二つの変更の逐次合成ではない。
既存の恒等式は
\[
 f_6=P,\quad f_7=f_8=f_9=f_{11}=0,\qquad
 f_{10}=cH-3PU_4,\quad f_{12}=PV_6-dP(X_7),\quad f_{13}=-dP(X_8).
\]
従って $g_{12}=g_{13}=0$ である。
$A_{14}=g_{14}$、$A_{16}=g_{16}$ と置き、
\[
 dP(\Xi_9)=A_{14},\qquad dP(\Xi_{10})=f_{15},\qquad
 dP(\Xi_{11})=A_{16}
\]
という純粋な持ち上げで14、15、16次を順に消す。
続く18、20、22次の入力も既存の結果により通常像に入るので、
それぞれ $\Xi_{13},\Xi_{15},\Xi_{17}$ で消す。

\begin{proposition}[二十四次の正確な有限補正式]
\label{prop:v17-exact-twenty-fourth-correction}
初回変更の24次係数は
\begin{equation}\label{eq:v17-initial-twenty-fourth}
\begin{split}
 g_{24}={}&f_{24}+df_{18}(X_7)+df_{17}(X_8)
       +\tfrac12 f_{12}''(X_7,X_7)+\tfrac12 f_{10}''(X_8,X_8)\\
 &+\tfrac16P'''(X_7,X_7,X_7)\\
 &+3U_4\bigl(f_{20}+df_{14}(X_7)+df_{13}(X_8)
                         +\tfrac12P''(X_8,X_8)\bigr)\\
 &-V_6\bigl(f_{18}+df_{12}(X_7)+\tfrac12P''(X_7,X_7)\bigr).
\end{split}
\end{equation}
上の純粋な有限消去後は正確に
\begin{equation}\label{eq:v17-finite-twenty-fourth}
\begin{split}
 R_{24}^{\mathrm{fin}}={}&g_{24}-dA_{16}(\Xi_9)-df_{15}(\Xi_{10})
                  +\tfrac12P''(\Xi_{10},\Xi_{10})
                  -c\,dH(\Xi_{15})
\end{split}
\end{equation}
となる。純粋な17次・19次消去を挿入しても24次係数は変わらない。
\end{proposition}

\begin{proof}
初回場の次数増分は6と7である。24次へ到達する一次項は
$df_{18}(X_7)$、$df_{17}(X_8)$、二次項は
$f_{12}''(X_7,X_7)/2$、$f_{10}''(X_8,X_8)/2$ である。
混合項には零の $f_{11}$ が必要となる。
三階項は $6+3\cdot6=24$ の $P'''(X_7,X_7,X_7)/6$ が初めて現れ、
四階以上は30次以上である。
四次単元を掛ける前の20次係数と、六次単元を掛ける前の18次係数を展開すると、
\eqref{eq:v17-initial-twenty-fourth}の全十三項を得る。

14次消去の増分8は16次項へ作用する。
15次消去の増分9は15次項へ作用し、さらに $6+2\cdot9=24$ の
$P$ 二階項を生む。この段の15次係数は元の $f_{15}$ のままである。
16、18、22次消去の対応先は、既に零である14、12、8次係数である。
20次消去の増分14は $cH$ に作用する。
17次と19次の純粋消去の対応先は零の13次と11次であり、
各々の二階Taylor項も24次を超える。
これで\eqref{eq:v17-finite-twenty-fourth}を得る。
\end{proof}

\nolinkurl{v17_degree_ledger.py} はこの十三項と、
14次から23次までの各純粋消去・単元を残す消去の次数を独立に列挙する。
これは次数の検査であり、通常像への所属の証明とは区別する。

\begin{remark}[単元を別に残す場合]
14次消去を $(1-B_8)g(a-\Xi'_9)$ とするなら、
24次に $-B_8A_{16}$ が追加される。
15次の入力を $f_{15}=dP(\Xi'_{10})+PB_9$ として消すと、
自己作用から $-PB_9^2$ が追加される。
20次消去の単元は $-cB_{14}H$ を加える。
それぞれ $A_{16}\in J_{16}$、$P\mathcal U_{18}\subset J_{24}$、
$H\mathcal U_{14}\subset J_{24}$ によって通常商では消える。
別規約の場を用いた二階項は、Euler方向を明示して比較する。
\end{remark}

\subsection{三階Taylor項を含む四つの二階族}

$f_{18}\in J_{18}$ は第14版で確定しており、
命題\ref{prop:v17-seventeenth-contact-basis}により
$f_{17}\in\mathcal C_{17}=J_{17}$ である。
従って
\[
 dP(Z_{18})=f_{18},\quad Z_{18}\in\mathcal G_{13},\qquad
 dP(Z_{17})=f_{17},\quad Z_{17}\in\mathcal G_{12}
\]
と選べる。系\ref{cor:v16-all-hessian-first-actions}から
\[
 W_{12}:\mathcal G_8\longrightarrow\mathcal G_{12},\qquad
                       dP(W_{12}(X))=dH(X)
\]
という純粋な線形持ち上げも選べる。
第14版の $P''(\mathcal G_7,\mathcal G_7)\subset J_{18}$ を用いて
\[
 V_{13}\in\mathcal G_{13},\qquad dP(V_{13})=P''(X_7,X_7),
 \qquad A=DX_7\cdot X_7\in\mathcal G_{13}
\]
とする。これらは全て普遍共変量の有限次元空間で選び、固定した $F$ ごとの除算はしない。

\begin{theorem}[二十四次の実際の類の四族代表]
\label{thm:v17-four-family-twenty-fourth-representative}
通常像を法として
\begin{equation}\label{eq:v17-four-family-twenty-fourth}
\begin{split}
 [R_{24}^{\mathrm{fin}}]=\bigl[f_{24}
 &+P''\bigl(X_7,Z_{18}-\tfrac13(V_{13}+A)\bigr)\\
 &+P''\bigl(X_8,Z_{17}+\tfrac c2W_{12}(X_8)\bigr)\\
 &-P''(\Xi_9,\Xi_{11})-\tfrac12P''(\Xi_{10},\Xi_{10})\bigr]
 \quad\text{in }\mathcal C_{24}/J_{24}.
\end{split}
\end{equation}
従って必要な補正は
\begin{equation}\label{eq:v17-four-pairing-families}
 P''(\mathcal G_7,\mathcal G_{13}),\quad
 P''(\mathcal G_8,\mathcal G_{12}),\quad
 P''(\mathcal G_9,\mathcal G_{11}),\quad
 P''(\mathcal G_{10},\mathcal G_{10})
\end{equation}
の四族に含まれる。
\end{theorem}

\begin{proof}
第14版と式\eqref{eq:v15-quadratic-corrections-contact}の既存の所属により
\[
 P''(\mathcal G_7,\mathcal G_7)\subset J_{18},\qquad
 P''(\mathcal G_7,\mathcal G_9)+P''(\mathcal G_8,\mathcal G_8)\subset J_{20}.
\]
$f_{20}\in J_{20}$、$f_{18}\in J_{18}$、
$f_{14}\in J_{14}$、$f_{13}=-dP(X_8)$、
$f_{12}=PV_6-dP(X_7)$ を用いて微分すると、
\eqref{eq:v17-initial-twenty-fourth}の $3U_4$ 群と $-V_6$ 群は全て $J_{24}$ に入る。
ここでは $uJ_n\subset J_{n+s}$（$u\in\mathcal U_s$）を用いている。
同じ理由で $(PU_4)''(X_8,X_8)$ と $(PV_6)''(X_7,X_7)$ も $J_{24}$ に入る。

持ち上げの恒等式を微分して
\[
 df_{18}(X_7)\equiv P''(X_7,Z_{18}),\qquad
 df_{17}(X_8)\equiv P''(X_8,Z_{17})\pmod{J_{24}}
\]
を得る。$dH(\mathcal G_{15})\subset J_{24}$ を用いると
\[
 H''(X_8,X_8)+dH(DX_8\cdot X_8)
   =P''(X_8,W_{12}(X_8))+dP(D(W_{12}(X_8))\cdot X_8)
\]
から $H''(X_8,X_8)\equiv P''(X_8,W_{12}(X_8))$ となる。

三階項の符号と係数を明記する。
点ごとの二階微分に対して
\[
 (dP(X_7))''(X_7,X_7)
  =P'''(X_7,X_7,X_7)+2P''(A,X_7)+dP(D^2X_7(X_7,X_7)).
\]
従って
\[
 \tfrac12f_{12}''(X_7,X_7)+\tfrac16P'''(X_7,X_7,X_7)
 \equiv-\tfrac13P'''(X_7,X_7,X_7)-P''(A,X_7).
\]
一方、$dP(V_{13})=P''(X_7,X_7)$ を $X_7$ で微分すると
\[
 P'''(X_7,X_7,X_7)+2P''(A,X_7)
 =P''(X_7,V_{13})+dP(DV_{13}\cdot X_7).
\]
二式を合わせて
\[
 \tfrac12f_{12}''(X_7,X_7)+\tfrac16P'''(X_7,X_7,X_7)
 \equiv-\tfrac13P''(X_7,V_{13}+A)\pmod{J_{24}}
\]
を得る。捨てた $dP$ 項の場は全て $\mathcal G_{19}$ に属する。

最後に $A_{16}=dP(\Xi_{11})$、$f_{15}=dP(\Xi_{10})$ により
\[
 -dA_{16}(\Xi_9)\equiv-P''(\Xi_9,\Xi_{11}),\qquad
 -df_{15}(\Xi_{10})+\tfrac12P''(\Xi_{10},\Xi_{10})
       \equiv-\tfrac12P''(\Xi_{10},\Xi_{10}).
\]
$-c\,dH(\Xi_{15})\in J_{24}$ と合わせると主張を得る。
\end{proof}

\subsection{十九次奇項と全偶化との比較}

\begin{lemma}[十九次の純粋消去]
\label{lem:v17-nineteenth-odd-removal}
上の有限正規化に17次の純粋消去を加えた後の19次係数 $R_{19}$ は $J_{19}$ に入る。
純粋な14次場でこれを消しても24次係数は変わらない。
\end{lemma}

\begin{proof}
17次消去は補題\ref{lem:v16-seventeenth-odd-removal}で可能である。
第16版の全偶正規形との接触変換を比較すると、差は19次以上であり、
最小増分4の対数補題から
\[
                        -R_{19}=Q_{13}P+cQ_9H.
\]
$dH(\mathcal G_{10})+H\mathcal U_9\subset J_{19}$ により右辺は $J_{19}$ に入る。
$dP(\Xi_{14})=R_{19}$ として $a-\Xi_{14}$ で消すと、
24次への一次作用には零の11次係数が必要である。
$P$ の二階項は32次からなので、24次は変わらない。
\end{proof}

\begin{corollary}[全偶正規形の二十四次類]
\label{cor:v17-evenization-twenty-fourth-class}
第16版の全偶正規形の24次係数を $R_{24}^{\mathrm{ev}}$ とすると
\[
                   R_{24}^{\mathrm{ev}}-R_{24}^{\mathrm{fin}}\in J_{24}.
\]
従って式\eqref{eq:v17-four-family-twenty-fourth}は実際の24次類を与える。
\end{corollary}

\begin{proof}
17次と19次を純粋に消した有限形は
$P+cH+R_{21}+R_{23}+R_{24}^{\mathrm{fin}}+O(a^{25})$ である。
全偶形との差は21次以上なので、指数作用と対数作用は24次まで一致する。
$L(R_{21}+R_{23})$ は25次以上であり、24次の差は $Q_{18}P+cQ_{14}H$ となる。
この二項はともに $J_{24}$ に属する。
\end{proof}

\begin{remark}[十九次を残す直接比較]
19次を残したまま比べると、24次には
$\tfrac12Q_5R_{19}$ が加わり得る。
$Q_5P=0$ から $Q_5$ は通常像を保存するので、この項も $J_{24}$ に入る。
ただし前版の「19次消去を明示しなくてよい」という結論だけを引用して、
二階対数項そのものを無視することはできない。
\end{remark}

\subsection{三百十組による有限の所属条件}

\begin{proposition}[二十四次補正の有限検査]
\label{prop:v17-three-hundred-ten-pairings}
$J_{12}/P\mathcal U_6$、$J_{13}/P\mathcal U_7$、
$J_{14}/P\mathcal U_8$、$J_{15}/P\mathcal U_9$、
$J_{16}/P\mathcal U_{10}$、$J_{17}/P\mathcal U_{11}$、
$J_{18}/P\mathcal U_{12}$ の基底を持ち上げる場を選ぶ。
それらの個数は順に
\[
                         2,\ 1,\ 7,\ 6,\ 21,\ 22,\ 60
\]
である。次の310個の共変量が全て $J_{24}$ に入るなら、
$R_{24}^{\mathrm{ev}}-d_{24}^{\mathrm T}\in J_{24}$ となる。
\[
\begin{array}{ll}
 \operatorname{Hess}_F(q)(dq(\Gamma_{7,i}),dq(\Gamma_{13,j})),
       &1\le i\le2,\quad1\le j\le60,\\
 \operatorname{Hess}_F(q)(dq(\Gamma_8),dq(\Gamma_{12,k})),
       &1\le k\le22,\\
 \operatorname{Hess}_F(q)(dq(\Gamma_{9,i}),dq(\Gamma_{11,j})),
       &1\le i\le7,\quad1\le j\le21,\\
 \operatorname{Hess}_F(q)(dq(\Gamma_{10,i}),dq(\Gamma_{10,j})),
       &1\le i\le j\le6.
\end{array}
\]
\end{proposition}

\begin{proof}
$dP(K)=0$ なら
$P''(K,Y)=-dP(DK\cdot Y)$ は通常像に入る。
Euler単元場 $au/6$ については
$P''(au/6,Y)=(5/6)u\,dP(Y)$ も通常像に入る。
従って各二階族の通常類は、各 $dP$ 像を $P$ 倍単元で割った商にのみ依存する。
さらに
\[
 P''(Y,Z)=\operatorname{Hess}_F(q)(dq(Y),dq(Z))
                 +dF(d^2q(Y,Z))
\]
の第二項は、$d^2q(Y,Z)\in\mathcal D_{20}$ と全次数Pfaffian全射により $J_{24}$ に入る。
必要な組数は $2\cdot60+1\cdot22+7\cdot21+6\cdot7/2=310$ である。
各代表の所属から四族全体の所属を得て、前定理と偶化比較を適用する。
\end{proof}

この命題の仮定は普遍な $J_{24}$ 所属であり、有限点での零値と置き換えてはならない。
実際、今回のHesse入力では二階補正の一部が点ごとには非零となる。
後の検出子による消去、または通常商全体を尽くす構造定理が別に必要である。
本版では通常24次階数の確定とnormal射影を合わせ、
系\ref{cor:v17-actual-twenty-fourth-class-zero}で全四族の $J_{24}$ 所属を結論する。

% END contact/contact_twenty_fourth_ja.tex


% BEGIN tate/v17_tate_notes_ja.tex
\section{二十四次Tate係数の新規再帰}
\label{sec:v17-tate24}
この節では、通常接触商への射影とは独立に、四つの明示的入力で
raw Tate係数を有理数体上で計算する。
規約は従来通り $d^{\mathrm T}=\tfrac12\det K$ とする。

\begin{proposition}[二十四次に必要な再帰の範囲]\label{prop:v17-tate-truncation}
二十四次以下の行列式を決めるには $K_1,\ldots,K_{20}$ で十分である。
これらの定数ブロックは、$A_1,\ldots,A_{16}$ を全項で求め、
曲率 $R_{17},R_{18},R_{19},R_{20}$ をそれぞれ内部$x$次数
$4,3,2,1$ 以下へ切断してからEulerホモトピーを適用すれば厳密に求まる。
対応する $A_{17},A_{18},A_{19},A_{20}$ の保持範囲は $3,2,1,0$ である。
\end{proposition}
\begin{proof}
五つの行列成分はいずれも$a$次数一以上から始まるため、
二十四次以下の積に現れる一因子の次数は高々 $24-4=20$ である。
内部重み $2\deg_x-\#\theta+\#\eta=3$ から全 $A_i$ の$x$次数は高々四である。
Clifford積は$x$次数を下げず、Eulerホモトピーはそれを一つ下げる。
第$n$段階の定数ブロックへ最終的に影響し得る範囲を逆算すると、
$n\ge17$ では曲率の$x$次数 $21-n$ 以下、補正の$x$次数 $20-n$ 以下を
残せばよい。既に切断した第$i$段階 ($i<n$) の上限 $20-i$ は
次の曲率に必要な上限 $21-n$ 以上であり、捨てた項は戻らない。
第16版で切断した $A_{15},A_{16}$ はここで全項として流用せず、
本節の再帰では最初から全項を新たに生成する。
\end{proof}

\subsection{二つのHesse型入力}
$m\in\{-1,-2\}$ とし、
\[
 F_m=x_0^3+x_1^3+x_2^3+x_3^3+(-2-6m)x_4^3+6m x_2x_3x_4,
 \qquad
 a_b=b\,e_{01}+b^{-1}(e_2-e_4)\wedge(e_3-e_4)
\]
を $b=1,2$ で評価する。対称三階テンソルの規約では $F_{234}=m$ である。
$q(a_b)=(0,0,1,1,1)$ かつ $F_m(q)=0$ であり、
三元部分の滑らかさの判別条件に現れる数
$(-2-6m)+8m^3=2(m-1)(2m+1)^2$ はそれぞれ $-4,-54$ と非零である。
従ってどちらの五変数三次式も滑らかである。

\begin{proposition}[四入力のraw二十四次係数]\label{prop:v17-tate24-values}
次の有理数を置く。
\begin{align*}
 A_{24,1}&=-\frac{2504247107738347979241135831211}
                      {1458881718210000000000000},\\
 A_{24,2}&=-\frac{6596059248103221614470055785281}
                      {131766880000000000000}.
\end{align*}
四入力での係数は
\[
 \begin{array}{c|cc}
       &b=1&b=2\\\hline
  m=-1&A_{24,1}&64A_{24,1}\\
  m=-2&A_{24,2}&64A_{24,2}
 \end{array}
\]
である。この表はraw係数 $f_{24}=d^{\mathrm T}_{24}$ の値であり、
正規化後の類との比較は接触作用の節で別に行う。
\end{proposition}
\begin{proof}
\texttt{v17\_hesse\_tate\_twentyfourth.py} は各入力から $A_1$ を構成し、
既存の再帰データを入力として読み込まず、命題
\ref{prop:v17-tate-truncation} の全段階を有理数体上で計算する。
共通分母をもつ整数多項式として積とEulerホモトピーを実装し、
各段階の閉性とMaurer--Cartan方程式を保持範囲で厳密に検査する。
混合Hesse項については自然な外積$\eta$記号とPBW記号の変換を先に適用し、
PBW順序のままホモトピーを取らない。
独立な32次元外積作用素で、素数 $1000003$ を法として
全 $0,\ldots,16$ 次の方程式と全20定数ブロックを照合する。
行列式は有理数体上で置換展開と部分集合動的計画の二方法により照合した。
四入力の全テンソル、交代行列、20ブロック、25行列式係数、
各段階の記録は同梱のJSONとログに保存している。
\end{proof}

独立な保存値監査は、上の再帰とは別の標準ライブラリだけを用いて
20ブロックから全25係数を再構成する。
$m=-2,b=1$ では共通分母法とは別に各係数を\texttt{Fraction}で扱う再帰も行い、
全ブロックと全行列式係数を比較する。
$m=-1,b=1,2$ は第16版と同一入力なので、旧18ブロックと22次以下の全係数が
一致することを確認する。保存値を再読する監査と、最初からの再帰の再実行とは区別する。

\subsection{法方向の三次jetを直接生成する再帰}
\label{subsec:v17-tate-normal-jets}
上の四点値とは別に、$\varepsilon$ を独立変数として
\[
 a(\varepsilon)=e_{01}+\varepsilon(e_2-e_4)\wedge(e_3-e_4),
 \qquad q(a(\varepsilon))=\varepsilon(0,0,1,1,1)
\]
と置く。係数環を $\mathbb Q[\varepsilon]/(\varepsilon^4)$ に置き換えて
同じMaurer--Cartan再帰を最初から行う。
Clifford積、微分、Eulerホモトピーのいずれも$\varepsilon$次数を下げないため、
この切断は行列式の$\varepsilon^3$係数を厳密に保持する。
$x$次数に関する切断は命題\ref{prop:v17-tate-truncation}と独立に併用できる。

\begin{proposition}[二十四次の二つの法三次jet]
\label{prop:v17-tate24-normal-jets}
$m=-1,-2$ の両方について
\[
 [\varepsilon^3]\,d_{24}^{\mathrm T}(F_m,a(\varepsilon))=0.
\]
これは四つの有限点の値からの補間ではなく、
$\varepsilon$を形式変数とした再帰の厳密な計算結果である。
\end{proposition}
\begin{proof}
\nolinkurl{tate/v17_tate_normal_jet.py} は各 $A_n$ を
$\varepsilon^0,\ldots,\varepsilon^3$ の四つの係数多項式として保持する。
共通分母整数演算によって全 $A_1,\ldots,A_{16}$ と所定の$x$切断をもつ
$A_{17},\ldots,A_{20}$ を新たに生成し、各$\varepsilon$係数で閉性と
Maurer--Cartan方程式を有理数体上で確認する。
独立な32次元外積作用素では $a$次数 $0,\ldots,16$ と
$\varepsilon$次数 $0,\ldots,3$ の全組を素数 $1000003$ 上で確認する。
さらに20ブロックの四係数それぞれを独立な外積作用素で照合する。

$a$次数と$\varepsilon$次数の二重多項式として行列式を展開し、
置換展開と部分集合動的計画の二方法で全 $25\times4$ 係数を比較する。
両入力で二十四次・$\varepsilon^3$ の係数は零である。
入力、全80係数行列、100行列式係数、再帰段階の記録を
\nolinkurl{v17_tate_normal_m1.json} と
\nolinkurl{v17_tate_normal_m2.json} に保存した。
\end{proof}

この二点は三元三次式自身の値が零であり、そのHessianの値は非零である。
したがって横断射影定理と二つの不変式単項式の評価行列を組み合わせれば、
raw二十四次係数の $\Theta_8$ 成分を決められる。
命題\ref{prop:v17-ordinary-contact-rank24}によって
$\mathcal K_{24}^{\mathrm{nor}}=0$ も確定している。
この階数定理と低次正規化の補正を合わせる段階は、次の二節で説明する。

% END tate/v17_tate_notes_ja.tex


% BEGIN detector/v17_detector_notes_ja.tex
% 第17版：二点normal三次jetでTheta二方向を読む。
% 通常24次像の全階数とは独立に、normal射影の検出を確定する。

\section{横断三次項による二十四次の二つの係数検出}
\label{sec:v17-normal-third-jet-detectors}

通常商全体を評価グラフの大行列から復元する前に、
定理\ref{thm:v17-split-normal-projection}の射影
\[
 \pi_{24}:\mathcal C_{24}/J_{24}\longrightarrow
 \mathcal A_{16}
 =\mathbb C\,\overline U_4^4\oplus
   \mathbb C\,\overline U_4\overline J_6^2
\]
を二つの横断三次項で読む。
この方法による全接触像の消去は構造定理から従い、
有限個の接触グラフを検査して全てを尽くしたと推測するものではない。

\subsection{二つのHesse入力と正規化}

$m=-1,-2$ に対して
\[
\begin{split}
 F_m&=x_0^3+x_1^3+x_2^3+x_3^3+(-2-6m)x_4^3+6m x_2x_3x_4,\\
 a_\varepsilon&=e_0\wedge e_1+
             \varepsilon(e_2-e_4)\wedge(e_3-e_4)
\end{split}
\]
と置く。$q(a_\varepsilon)=(0,0,\varepsilon,\varepsilon,\varepsilon)$ であり、
三元部分 $f_m$ は $f_m(1,1,1)=0$ を満たす。
$h_f=\det(\operatorname{Hess}_f/6)$ とする。
錐上の単元制限と三元Hessianの値は
\[
\begin{array}{c|rrr}
 m&\overline U_4(f_m)&\overline J_6(f_m)&h_{f_m}(1,1,1)\\ \hline
 -1&5&528&-4\\
 -2&36&7128&-54
\end{array}
\]
である。\nolinkurl{v17_normal_detector_certificate.py} は
$a_0=e_0\wedge e_1$ において二つの単元グラフを有理数上で再縮約し、
この表を検証する。$U_4$ はraw四頂点グラフの $1/24$ 倍であり、
その正規化も固定している。

線形汎関数
\[
 j_m(C)=[\varepsilon^3]C(F_m,a_\varepsilon)
 \qquad(C\in\mathcal C_{24})
\]
を定める。$W_1=\Theta_8U_4^4$、$W_2=\Theta_8U_4J_6^2$ とすると
\begin{equation}\label{eq:v17-normal-jet-matrix}
 M=\begin{pmatrix}
 j_{-1}(W_1)&j_{-1}(W_2)\\
 j_{-2}(W_1)&j_{-2}(W_2)
 \end{pmatrix}
 =
 \begin{pmatrix}
 -2500&-5575680\\
 -90699264&-98771498496
 \end{pmatrix}.
\end{equation}
\[
 \det M=-258781326059520\ne0.
\]

\begin{proposition}[二点の横断三次係数検出子]
\label{prop:v17-normal-jet-detectors}
\begin{align}
 \rho_1^{\mathrm{nor}}
   &=\frac{177147\,j_{-1}-10\,j_{-2}}{464125140},
       \label{eq:v17-normal-rho-one}\\
 \rho_2^{\mathrm{nor}}
   &=\frac{-22674816\,j_{-1}+625\,j_{-2}}{64695331514880}
       \label{eq:v17-normal-rho-two}
\end{align}
とすると、$\rho_i^{\mathrm{nor}}(J_{24})=0$、
$\rho_i^{\mathrm{nor}}(W_j)=\delta_{ij}$ である。
任意の $C\in\mathcal C_{24}$ に対して
\[
 \pi_{24}([C])
  =\rho_1^{\mathrm{nor}}(C)\,\overline U_4^4
       +\rho_2^{\mathrm{nor}}(C)\,\overline U_4\overline J_6^2.
\]
汎関数は $\mathbb Q$ 上で定義され、$\mathbb C$ へ係数拡大して用いる。
この命題は通常商全体の次元を仮定しない。
\end{proposition}

\begin{proof}
構造定理によって
\[
 \gamma_{24}(C)=A_C f+B_C h_f,\qquad
 B_C=b_1\overline U_4^4+b_2\overline U_4\overline J_6^2
\]
と書ける。二入力では $f_m(1,1,1)=0$ なので
\[
 \begin{pmatrix}j_{-1}(C)\\j_{-2}(C)\end{pmatrix}
       =M\begin{pmatrix}b_1\\b_2\end{pmatrix}.
\]
式\eqref{eq:v17-normal-rho-one}--\eqref{eq:v17-normal-rho-two}は
この $2\times2$ 行列の逆行列を有理数で計算したものである。
$C\in J_{24}$ では構造定理により $B_C=0$ なので、
二つの $j_m$ 自体が $J_{24}$ を消す。
これで全ての主張が従う。
\end{proof}

\begin{remark}[通常商との区別]
常に
\[
 \frac{\ker(\rho_1^{\mathrm{nor}},\rho_2^{\mathrm{nor}})}{J_{24}}
                      =\mathcal K_{24}^{\mathrm{nor}}
\]
である。命題\ref{prop:v17-ordinary-contact-rank24}は
$\dim J_{24}=2456$ と、通常商が $[W_1],[W_2]$ で尽くされることを証明した。
従って24次では
$\ker(\rho_1^{\mathrm{nor}},\rho_2^{\mathrm{nor}})=J_{24}$ となる。
これは横断検出子の構成に、別に証明した通常商の次元を合わせた結論である。
\end{remark}

\subsection{実際の係数と低次補正の横断射影}

\begin{proposition}[二十四次の実際のnormal射影は零である]
\label{prop:v17-actual-normal-projection-zero}
未正規化Tate係数と、第16版の偶正規形の24次係数は
\[
 \rho_1^{\mathrm{nor}}(d_{24}^{\mathrm T})
 =\rho_2^{\mathrm{nor}}(d_{24}^{\mathrm T})=0,\qquad
 \pi_{24}([R_{24}^{\mathrm{ev}}])
 =\pi_{24}([d_{24}^{\mathrm T}])=0
\]
を満たす。従って両者の通常類は $\mathcal K_{24}^{\mathrm{nor}}$ に入る。
この主張は通常24次階数の決定を必要としない。
\end{proposition}

\begin{proof}
命題\ref{prop:v17-tate24-normal-jets}のfreshな横断再帰により
\[
 [\varepsilon^3]d_{24}^{\mathrm T}(F_{-1},a_\varepsilon)
 =
 [\varepsilon^3]d_{24}^{\mathrm T}(F_{-2},a_\varepsilon)=0
\]
である。前命題の逆行列を適用すれば二つの検出値が零になる。

定理\ref{thm:v17-four-family-twenty-fourth-representative}と
系\ref{cor:v17-evenization-twenty-fourth-class}により、
$R_{24}^{\mathrm{ev}}-d_{24}^{\mathrm T}$ は $J_{24}$ と四つの二階族の和に入る。
定理\ref{thm:v17-quadratic-actions-normal-zero}は
全てのこの二階族の $\pi_{24}$ 射影を消す。
従って両者のnormal射影は一致し、零である。
\end{proof}

\begin{corollary}[実際の二十四次通常類の消滅]
\label{cor:v17-actual-twenty-fourth-class-zero}
\[
 d_{24}^{\mathrm T}\in J_{24},\qquad
 R_{24}^{\mathrm{ev}}\in J_{24},\qquad
 R_{24}^{\mathrm{ev}}-d_{24}^{\mathrm T}\in J_{24}.
\]
また24次を持つ全二階族 $P''(\mathcal G_r,\mathcal G_s)$ は $J_{24}$ に入る。
\end{corollary}

\begin{proof}
命題\ref{prop:v17-ordinary-contact-rank24}から
$\mathcal K_{24}^{\mathrm{nor}}=0$ であり、$\pi_{24}$ は同型である。
前命題のnormal射影零は、従って通常像への所属を意味する。
全二階族についても定理\ref{thm:v17-quadratic-actions-normal-zero}により
同じ同型の核に入り、零となる。
\end{proof}

\paragraph{保存ブロックの独立監査.}
\nolinkurl{v17_normal_saved_audit.py} は研究計算のPythonモジュールをimportせず、
標準ライブラリの有理数だけで保存した20次数・各4係数のブロックを読む。
全125個の三次テンソル成分と $a$ の定数・$\varepsilon$ 係数を検出側入力と照合し、
$K_1=-a$、$K_2=\operatorname{Hess}_F(q)/3$ も確認する。
120置換による疎多項式展開と、列部分集合による密係数の動的計画法を別に実装し、
各入力でTate次数 $0,\ldots,24$、$\varepsilon$ 次数 $0,\ldots,3$ の全100係数を
再計算して保存値と比較する。二入力の全200係数が一致し、
$d_{24}^{\mathrm T}$ のnormal三次値と二検出値は零であった。
入力JSONと生成ソースのSHA256も照合する。
これは保存ブロックからの監査であり、新規Maurer--Cartan再帰とは別である。

\subsection{有限評価による別の補助検査}

横断三次係数を使わない比較用データとして
\[
 a_b=b\,e_0\wedge e_1+b^{-1}(e_2-e_4)\wedge(e_3-e_4),
 \qquad m=-1,-2,\quad b=1,2
\]
の四点でも評価した。
$q(a_b)=(0,0,1,1,1)$、$P=0$ であり、五変数Hessianとその $q$ 勾配は点ごとに零である。
二候補 $W_1,W_2$ は $b^{18}$ に比例する。
選択した接触グラフには $b^6$ 型の値が現れ、
$64\operatorname{ev}_{(F_m,a_1)}-\operatorname{ev}_{(F_m,a_2)}$
はその検査済みの列を消す。
この有限グラフ検査だけを、全通常像を消す証明には用いない。

命題\ref{prop:v17-three-hundred-ten-pairings}の全310組を四点で有理数再縮約した。
非零なのは $G_7$ の第1代表と $G_{13}$ の第17、第46代表の二組だけである。
この二組をそれぞれ $B_{17},B_{46}$ と書くと、$b=1$ の値は
\[
\begin{array}{c|rr}
 m&B_{17}&B_{46}\\ \hline
 -1&-67584&-135168\\
 -2&-2244806784&-4489613568
\end{array}
\]
となった。$b=2$ ではそれぞれ64倍であり、残る308組は四点全てで零である。
従ってこの二階補正は点ごとに全て零なのではなく、表示した二点差分で消える。
グラフの順序と値は \nolinkurl{v17_lower_pairings_exact.json} に記録する。
全四族のnormal射影零は前述の構造定理によって証明し、
この有限評価はそれとは別の検査として保持する。

% END detector/v17_detector_notes_ja.tex


% BEGIN contact/twenty_sixth_normal_form_ja.tex
% 第17版：通常24次階数2456の独立監査を含む確定最終定理。

\section{二十四次の消去から二十六次開始の正規形へ}
\label{sec:v17-twenty-sixth-normal-form}

\begin{theorem}[二十六次開始の普遍代数的収束偶正規形]
\label{thm:v17-twenty-sixth-normal-form}
実際の境界の24次類は零であり、任意の滑らかな三次式 $F$ に対して
\begin{equation}\label{eq:v17-twenty-sixth-normal-form}
 \boxed{\displaystyle
 d^{\mathrm T}\sim
 P_F+\frac1{96}\det\operatorname{Hess}_F(q(a))
      +\sum_{\substack{n\ge26\\n\ \mathrm{even}}}R_n,
 \qquad R_n\in\mathcal C_n}
\end{equation}
という全次数の偶正規形が存在する。
接触変換は普遍同変な代数的冪級数であり、各形式係数は $F$ に多項式的である。
座標変更の線形項は恒等写像、単元の定数項は1であり、
各固定 $F$ では変換・逆変換・正規形は原点近傍で収束する。
26次以降の実際の類の消滅は主張しない。
\end{theorem}

\begin{proof}
第16版の定理\ref{thm:v16-twenty-fourth-normal-form}から
\[
 f_{\mathrm{ev}}=P_F+\frac1{96}H+R_{24}^{\mathrm{ev}}
                                      +R_{26}^{\mathrm{ev}}+\cdots
\]
という普遍代数的収束偶正規形を得る。
命題\ref{prop:v17-actual-normal-projection-zero}により
$\pi_{24}([R_{24}^{\mathrm{ev}}])=0$ である。
ここには十三項の初回台帳、三階Taylor項の二階族への還元、
十九次奇項の消去と全偶化との比較が含まれる。

命題\ref{prop:v17-ordinary-contact-rank24}により通常24次商は二次元であり、
$[\Theta_8U_4^4]$、$[\Theta_8U_4J_6^2]$ が基底である。
従って $\mathcal K_{24}^{\mathrm{nor}}=0$ であるため、
通常商上の $\pi_{24}$ は同型である。
系\ref{cor:v17-actual-twenty-fourth-class-zero}のとおり
$R_{24}^{\mathrm{ev}}\in J_{24}$ となる。
これはraw係数だけを正規化後の係数と同一視する議論ではなく、
四つの二階補正族のnormal射影零を経由した結論である。
また下位jetは $P+cH$ なので、構造定理の $I_{24}=J_{24}$ が適用できる。

普遍共変量の有限次元空間で
\[
 \Xi_{19}\in\mathcal G_{19},\qquad dP(\Xi_{19})=R_{24}^{\mathrm{ev}}
\]
と選び、奇多項式変更 $a\mapsto a-\Xi_{19}(F,a)$ を行う。
この場は $F$ 次数18、$a$ 次数19であり、判別式等による除算を必要としない。
$P$ への一次作用が24次を消し、
$H$ への最初の作用は28次、$P$ の二階項は42次からである。
従って既存の低次正規形を保ち、次に残り得る次数は26となる。
奇座標変更なので全次数の偶性も保たれる。

追加変更は有限多項式で原点でのJacobianが恒等写像であり、
形式逆は代数的、固定した $F$ では解析的な局所逆も存在する。
第16版の変換との有限合成によって、
代数性、局所収束性、普遍同変性と各係数の $F$ 多項式性が保たれる。
\end{proof}

\begin{remark}[次に残る次数と全Hessian形]
次に検討すべき実際の偶数次数は26である。
本定理は26次以降の類の消滅や、境界全体が $P+H/96$ に接触同値であることを
まだ主張しない。
また定理\ref{thm:v17-finite-normal-kernel}は未知部分の有限次元性を与えるが、
その全次数、安定域の開始次数、および
$\mathcal K^{\mathrm{nor}}=\mathbb C\cdot[H]$ かどうかは未決定である。
\end{remark}

% END contact/twenty_sixth_normal_form_ja.tex

\clearpage
% BEGIN contact/v18_normal_contact_functoriality_ja.tex
% 第18版：normal三次射影の全形式接触変換則。
% 第17版の二階P作用だけを扱う定理を、全高階作用へ強める。

\section{横断三次射影の全次数接触変換則}
\label{sec:v18-normal-contact-functoriality}

第17版では全二階 $P$ 補正のnormal射影が零であることを示した。
本節では任意の有限・形式接触変換に対する変換則を求める。
これにより、最初の非零normal係数は接触不変であることが分かり、
高次のraw Tate係数と正規化後の係数をnormal射影上で直接比較できる。
通常商の残る頂点部分の消失を、この変換則だけから結論するものではない。

$V=L_2\oplus Q_3$、$a_0=e_0\wedge e_1$ とし、
\[
 B(p)=p_4e_2\wedge e_3-p_3e_2\wedge e_4+p_2e_3\wedge e_4,
 \qquad a_\varepsilon=a_0+\varepsilon B(p),
 \qquad F_f=x_0^3+x_1^3+f
\]
と置く。第17版の記号
$\mathcal A=\mathbb C[\overline U_4,\overline J_6]$、
$\gamma_n(C)=A_Cf+B_Ch_f$、$h_f=\det(\operatorname{Hess}_f)/216$
を引き継ぐ。

\subsection{場の接線成分と法一次成分}

\begin{lemma}[rank二切片上の二つのスカラー]
\label{lem:v18-two-normal-field-scalars}
$Y_r\in\mathcal G_r$ に対して、一意な
$\eta_{Y_r},\beta_{Y_r}\in\mathcal A_{r-1}$ が存在し、
\begin{equation}\label{eq:v18-two-normal-field-scalars}
 Y_r(F_f,a_\varepsilon)
   =\eta_{Y_r}(f)a_0
     +\varepsilon\beta_{Y_r}(f)B(p)+O(\varepsilon^2)
\end{equation}
となる。右辺の誤差は $W$ の全成分についての誤差である。
第17版の $\alpha_{dqY_r}$ とは
\[
 \alpha_{dqY_r}=\eta_{Y_r}+\beta_{Y_r}
\]
という関係にある。特に $r$ が偶数なら両スカラーは零である。
\end{lemma}

\begin{proof}
$\lambda(t)=\operatorname{diag}(t,t,1,1,1)$ は $a_0$ を固定する。
$W=\Lambda^2V\otimes\Delta^{-1}$ の
$\Lambda^2L_2$、$L_2\otimes Q_3$、$\Lambda^2Q_3$ 成分の重みは順に
$0,-1,-2$、横断変数 $B$ の重みは $-2$ である。
従って $\varepsilon^k$ 係数に必要な $F$ の重みは、これらに $2k$ を加えた数である。
$k=0$ で非零になり得るのは $\Lambda^2L_2$ 成分だけで、
それは三元成分 $f$ にのみ依存する。
よって $Y_r(F_f,a_0)=\eta_{Y_r}(f)a_0$ と書ける。

$k=1$ の $\Lambda^2Q_3$ 成分は、$f$ のみに依存する $B$ の線形写像である。
$\mathrm{SL}(Q_3)$ に関してこれは $Q_3^*$ のEnd共変量だから、
補題\ref{lem:v17-ternary-endomorphism-scalar}によってスカラーである。
これを $\beta_{Y_r}$ と書く。
残る二成分の一次係数に必要な $F$ の重みは2と1である。
一方 $F_f$ の非零係数の重みは0と3だけなので、その両係数は零である。
両スカラーは $f$ の次数 $r-1$ の $\mathrm{SL}_3$ 不変式であり、
$\mathcal A_{r-1}$ に属する。
\[
 dq_{a_\varepsilon}(Y_r)
   =a_\varepsilon\wedge Y_r
   =\varepsilon(\eta_{Y_r}+\beta_{Y_r})a_0\wedge B+O(\varepsilon^2)
\]
が $\alpha$ との関係を与える。
$\mathcal A$ は次数4と6で生成されるので、奇数次成分は零である。
\end{proof}

\begin{proposition}[一つの斉次共変量の変換公式]
\label{prop:v18-homogeneous-normal-pullback}
$C_n\in\mathcal C_n$、$Y_i\in\mathcal G_{r_i}$ とし、$t_i$ を独立な形式変数とする。
\[
 \Psi_t(a)=a+\sum_i t_iY_i(F,a),\qquad
 s(t)=1+\sum_i t_i\eta_{Y_i},\qquad
 b(t)=1+\sum_i t_i\beta_{Y_i}
\]
と置けば
\begin{equation}\label{eq:v18-homogeneous-normal-pullback}
 [\varepsilon^3]C_n(F_f,\Psi_t(a_\varepsilon))
    =s(t)^{n-3}b(t)^3\gamma_n(C_n)(f,p)
\end{equation}
が成り立つ。
\end{proposition}

\begin{proof}
前補題から $\Psi_t(a_\varepsilon)=s(t)a_0+
\varepsilon b(t)B+O(\varepsilon^2)$ である。
$s(t)$ は形式単元なので、$a_{01}$ を可逆にしたSchur補チャートを用いる。
交差成分を除くdeterminant一の基底変更は $\varepsilon=0$ で恒等であり、
変換後の $F$ の三元成分は $f+O(\varepsilon)$ である。
Schur補の横断成分の一次項は正確に $\varepsilon b(t)B$ であり、
基点の $a_{01}$ 成分は $s(t)+O(\varepsilon)$ である。

補題\ref{lem:v17-covariant-normal-order}により $C_n$ の横断次数は三以上である。
従って三次項の計算では、基点の $O(\varepsilon)$ 変動、
$F$ の $O(\varepsilon)$ 変動、横断成分の $O(\varepsilon^2)$ 変動は寄与しない。
$a$ についての斉次性と横断三次性から
\[
 s(t)^n\left(\frac{b(t)}{s(t)}\right)^3\gamma_n(C_n)
\]
を得る。これが主張である。指数 $n-3$ は非負なので、
結果そのものは $s(t)$ の逆数を含まない多項式恒等式になっている。
\end{proof}

\begin{corollary}[全高階Taylor作用のnormal射影]
\label{cor:v18-all-higher-actions-normal-zero}
任意の $k\ge1$ と $Y_i\in\mathcal G_{r_i}$ に対し、
$n=6+\sum_i(r_i-1)$ と置くと
\[
 \pi_n\bigl([P^{(k)}(Y_1,\ldots,Y_k)]\bigr)=0.
\]
同様に、$H$ の全高階Taylor作用のnormal三次項は零である。
より一般に、$\pi([C_m])=0$ なら、$C_m$ の任意の点ごとの高階Taylor作用も
対応する次数でnormal射影が零となる。
\end{corollary}

\begin{proof}
前命題で $t_1\cdots t_k$ の係数を比較する。
右辺は $\gamma_m(C_m)$ の不変式倍なので、Hessian成分が零ならその性質を保つ。
$\gamma_6(P)=f$、$\gamma_{10}(H)=0$ を代入すれば二つの特例を得る。
ここで微分は点ごとの多重微分であり、場を逐次合成した微分とは区別する。
逐次合成した場合も連鎖律で現れる場の微分は普遍共変場であるため、同じ結論が従う。
\end{proof}

\subsection{全形式境界に対する変換則}

形式境界 $\mathfrak f=\sum_{n\ge6}C_n$ に対し、補助変数 $z$ を使って
\[
 \mathcal B_{\mathfrak f}(z)=\sum_{n\ge6}z^nB_{C_n}\in\mathcal A[[z]]
\]
と定める。これはnormal三次項のHessian係数を全次数でまとめたものである。

\begin{theorem}[normal係数の接触変換則]
\label{thm:v18-normal-contact-law}
普遍同変な接触変換を
\[
 \widetilde{\mathfrak f}(F,a)
   =u(F,a)\mathfrak f(F,\Phi(F,a)),\qquad
 u=1+\sum_{s>0}u_s,\quad
 \Phi=a+\sum_{r>1}Y_r
\]
と書く。各 $u_s\in\mathcal U_s$、$Y_r\in\mathcal G_r$ とし、
\[
 \overline u(z)=1+\sum_{s>0}z^s\overline u_s,
 \qquad
 \sigma(z)=1+\sum_{r>1}z^{r-1}\eta_{Y_r},
 \qquad
 \tau(z)=1+\sum_{r>1}z^{r-1}\beta_{Y_r}
\]
と置く。このとき
\begin{equation}\label{eq:v18-normal-contact-law}
 \boxed{\displaystyle
 \mathcal B_{\widetilde{\mathfrak f}}(z)
 =\overline u(z)\tau(z)^3\sigma(z)^{-3}
      \mathcal B_{\mathfrak f}\bigl(z\sigma(z)\bigr).}
\end{equation}
\end{theorem}

\begin{proof}
$\Phi(F,z a_\varepsilon)/z$
$=a_\varepsilon+\sum_{r>1}z^{r-1}Y_r(F,a_\varepsilon)$ と展開する。
前命題を各 $C_n$ に適用すると、そのnormal三次項は
$z^n\sigma(z)^{n-3}\tau(z)^3\gamma_n(C_n)$ である。
各 $C_n$ の横断次数は三以上なので、単元を掛けるときは
$u(F,z a_0)=\overline u(z)$ のみが寄与する。
Hessian側の係数を取り、$n$ について加えると公式を得る。
各 $z$ 次数で有限個の項しか寄与しないので、全ての操作は完備環上で定義される。
\end{proof}

\begin{corollary}[最初の非零normal係数は接触不変である]
\label{cor:v18-first-normal-contact-invariant}
上の接触変換について、$\mathcal B_{\mathfrak f}=0$ と
$\mathcal B_{\widetilde{\mathfrak f}}=0$ は同値である。
さらに一方の $N$ 次未満の係数が全て零なら他方もそうであり、
\[
 [z^N]\mathcal B_{\widetilde{\mathfrak f}}
 =[z^N]\mathcal B_{\mathfrak f}
\]
となる。
\end{corollary}

\begin{proof}
$\overline u,\sigma,\tau$ はいずれも定数項1の形式単元であり、
$z\sigma(z)=z+O(z^2)$ である。
公式の各係数は、同じ次数の元の係数と、それより低い次数の係数の和である。
次数に関する帰納法により結論を得る。
\end{proof}

\begin{corollary}[二十六次の実際のnormal係数]
\label{cor:v18-actual26-normal-equals-raw}
第17版の偶正規形を
\[
 \mathfrak f_{\mathrm{ev}}=P+\frac1{96}H+
          R_{26}^{\mathrm{ev}}+R_{28}^{\mathrm{ev}}+\cdots
\]
と書く。このとき未正規化Tate係数 $f_n=d_n^{\mathrm T}$ について
\[
 \pi_n([f_n])=0\quad(n<26),\qquad
 \pi_{26}([R_{26}^{\mathrm{ev}}])=\pi_{26}([f_{26}])
\]
が成り立つ。これは通常26次商の次元を仮定しない。
\end{corollary}

\begin{proof}
定理\ref{thm:v17-twenty-sixth-normal-form}の普遍形式接触変換へ前系を適用する。
正規形の26次未満には $P,H$ しかなく、両者のHessian側normal係数は零である。
よってraw側の同じ次数の係数も零であり、最初に残り得る26次係数は一致する。
\end{proof}

\begin{remark}[通常像への所属と区別する]
本節は $\mathcal K^{\mathrm{nor}}$ の消失も、通常26次商の二次元性も使わない。
従って $\pi_{26}([f_{26}])=0$ と分かれば実際の26次類は
$\mathcal K_{26}^{\mathrm{nor}}$ に入るが、その類自体の消滅には別の証明が必要である。
全次数でrawのnormal係数が零と証明できれば、この性質は任意の形式正規化後も保たれる。
\end{remark}

% END contact/v18_normal_contact_functoriality_ja.tex


% BEGIN theory/v18_actual_normal_vanishing_ja.tex
% 第18版：実際の境界の全次数normal消滅。有限標本からの推測ではない。
\section{実際の境界の横断Hessian係数は全次数で消える}
\label{sec:v18-actual-normal-vanishing}

第17版までの有限次数の消滅計算を、ここでは幾何学的な全次数の命題へ強める。
入力は、固定した平面三次曲線上のPicard族が階数二の直線を与えることと、
その非零点での境界の横断三次曲線が元の抽象曲線に一致することである。
Hesse pencilの $j$ 不変量が定数でないことを組み合わせると、
通常商の自由な直和因子に現れる実際の境界の係数が全て零になる。

\subsection{固定した平面のPicard族と横断三次式}

$V=L_2\oplus Q_3$、$a_0=e_0\wedge e_1$ とし、前節と同じ $B(p)$ を用いる。
$f\in\operatorname{Sym}^3Q_3$ を滑らかな平面三次曲線を定める式とし、
\[
 F_f=x_0^3+x_1^3+f,\qquad C_f=V(f)\subset\mathbb P(Q_3^*)
\]
と置く。この $F_f$ も滑らかである。
実際、五つの偏微分が同時に零なら $x_0=x_1=0$ となり、
残る三つの式は $f$ の滑らかさに反する。

先頭項が $P$ の普遍境界方程式を $d=\sum_{n\ge6}d_n$ とする。
定理\ref{thm:v17-split-normal-projection}により
\[
 \gamma_n(d_n)=A_n(f)f+B_n(f)h_f,
 \qquad h_f=\det(\operatorname{Hess}_f)/216
\]
と一意に書ける。以下では二つの形式級数
\begin{equation}\label{eq:v18-normal-pencil-series}
 \mathsf A_f(t)=\sum_{n\ge6}t^{n-6}A_n(f),\qquad
 \mathsf B_f(t)=\sum_{n\ge6}t^{n-6}B_n(f)
\end{equation}
を使う。$A_6=1$、$B_6=0$ なので
$\mathsf A_f(0)=1$、$\mathsf B_f(0)=0$ である。
$a$ に関する斉次性から、直線 $t a_0$ に沿う横断三次項は
\begin{equation}\label{eq:v18-geometric-normal-cubic}
 [\varepsilon^3]d(F_f,t a_0+\varepsilon B(p))
 =t^3\bigl(\mathsf A_f(t)f(p)+\mathsf B_f(t)h_f(p)\bigr).
\end{equation}

\begin{lemma}[横断三次曲線の抽象同型類]
\label{lem:v18-fixed-plane-normal-j}
式\eqref{eq:v18-geometric-normal-cubic}の括弧内が定める
$\mathbb C((t))$ 上の平面三次曲線の $j$ 不変量は $j(C_f)$ に等しい。
\end{lemma}

\begin{proof}
補題\ref{lem:universal-chart-plucker}の証明を、固定分解
$V=L_2\oplus Q_3$ に適用する。
固定平面の曲線 $C_f$ 上のPicard形式族は、可逆な形式パラメータ変更を除いて
直線 $t a_0$ 全体を与える。その非零形式一般点は $C_f$ 上の直線束の評価核に対応する。
適合部分空間は同じ普遍座標で標準Pl\"ucker錐であり、
$t\ne0$ では $\Lambda^2Q_3$ の三方向がその法空間を与える。
従って式\eqref{eq:v18-geometric-normal-cubic}は、その点のnormal coneの三次方程式である。

定理\ref{thm:local-bn-cubic}により、射影化したnormal coneは
$C_f$ を次数三の直線束によって埋め込んだものである。
直線束はPicardパラメータとともに変化してよいが、抽象曲線は同じ $C_f$ である。
従って $j$ は一定である。
形式一般点への適用は、定理\ref{thm:boundary-eighth-vanishing}で用いたものと同じであり、
ここではFermat三次式を一般の滑らかな $f$ に置き換えただけである。
方程式に掛かる単元と $t^3$ は、射影化した三次曲線を変えない。
\end{proof}

\subsection{Hesse pencilからの全次数消滅}

\begin{lemma}[Hesse pencilでの形式的な一意性]
\label{lem:v18-hesse-formal-rigidity}
$\lambda$ を超越元とし、$k=\mathbb C(\lambda)$ とする。
\[
 D=p_2^3+p_3^3+p_4^3,\qquad M=p_2p_3p_4,\qquad
 f_\lambda=D+6\lambda M
\]
と置く。$A(t),B(t)\in k[[t]]$、$A(0)=1$、$B(0)=0$ とし、
$A(t)f_\lambda+B(t)h_{f_\lambda}$ の $j$ 不変量が
$f_\lambda$ の $j$ 不変量に等しいなら、$B(t)=0$ である。
\end{lemma}

\begin{proof}
三行三列の行列式を展開すると
\[
 h_{f_\lambda}=-\lambda^2D+(1+2\lambda^3)M
\]
である。従って $A f_\lambda+B h_{f_\lambda}$ は、単元 $A-\lambda^2B$ を除いて
$D+6\mu(t)M$ と書ける。ただし
\begin{equation}\label{eq:v18-hesse-parameter-change}
 \mu(t)=\lambda+
 \frac{(1+8\lambda^3)B(t)}{6(A(t)-\lambda^2B(t))}.
\end{equation}
第\ref{sec:boundary-higher-jets}節で用いたHesse式を同じ規約に直すと
\begin{equation}\label{eq:v18-hesse-j-function}
 J(\lambda)=
 \frac{110592\lambda^3(1-\lambda^3)^3}{(1+8\lambda^3)^3}.
\end{equation}
これは定数でない有理関数なので、標数零の体 $k$ で $J'(\lambda)\ne0$ である。
$\mu(0)=\lambda$ と $J(\mu(t))=J(\lambda)$ から $\mu(t)=\lambda$ を得る。
例えば $\mu-\lambda$ の最初の非零項が $c t^r$ なら、
$J(\mu)-J(\lambda)$ の同じ次数の係数は $J'(\lambda)c\ne0$ となり矛盾する。
$1+8\lambda^3\ne0$ と式\eqref{eq:v18-hesse-parameter-change}から $B(t)=0$ が従う。
\end{proof}

\begin{theorem}[実際の境界の全次数normal消滅]
\label{thm:v18-all-order-normal-vanishing}
実際のBrill--Noether境界の、先頭項を $P$ に規格化した普遍方程式
$d=\sum_{n\ge6}d_n$ に対して
\begin{equation}\label{eq:v18-all-order-normal-vanishing}
 \boxed{\displaystyle
  \pi_n([d_n])=0\quad\text{for every }n\ge6,\qquad
  \mathcal B_d(z)=0.}
\end{equation}
特に未正規化Tate方程式 $d^{\mathrm T}=\frac12\det K$ にもこの結論が成立し、
任意の普遍形式接触変換の後も成立する。
\end{theorem}

\begin{proof}
まず幾何学的な普遍座標で考える。
前二補題を一般のHesse式に適用すると
$\mathsf B_{f_\lambda}(t)=0$ となる。
ここで超越的な $\lambda$ への移行は、各滑らかな複素パラメータで得た
$j$ の等式を各 $t$ 係数ごとの有理関数恒等式として読むことで正当化できる。
従って各 $n$ について $B_n(f_\lambda)=0$ が $\mathbb C(\lambda)$ で成立する。
各 $B_n$ は三元三次式の全係数空間上の斉次多項式不変量である。
補題\ref{lem:v17-ternary-endomorphism-scalar}の証明で確認した
\[
 \mathrm{GL}_3\times\mathbb A^1\longrightarrow\operatorname{Sym}^3Q_3,
 \qquad(h,\lambda)\longmapsto h f_\lambda
\]
の優勢性と、$\mathrm{GL}_3$ に対する相対不変性から、$B_n$ は恒等的に零となる。
この最後の多項式恒等式は、特異な三元三次式も含めて成立する。
滑らかな開集合上で $j$ を用いた段階と、この多項式的延長の段階を区別している。

定理\ref{thm:tate-boundary-determinant}および
定理\ref{thm:tate-tenth-coefficient}の証明が示すように、Tate座標も
全非零三次式上で正則かつ同変で、接写像を同じedge同型に固定した普遍座標である。
従ってその方程式と上の方程式は、定数項1の単元と一次項が恒等の
普遍形式座標変換で移り合う。
定理\ref{thm:v18-normal-contact-law}により $\mathcal B=0$ はこの変換で保存される。
これでTate方程式と、さらに任意の普遍形式接触変換後について結論を得る。
\end{proof}

\begin{corollary}[全ての実際の剰余類が入る部分]
\label{cor:v18-all-actual-classes-vertex}
実際の境界を有限次数まで普遍形式接触正規化したとき、
残る各斉次係数の通常商における類は $\mathcal K_n^{\mathrm{nor}}$ に入る。
特に、第17版の偶正規形について
\[
 [R_{26}^{\mathrm{ev}}]\in\mathcal K_{26}^{\mathrm{nor}}.
\]
これは $\mathcal K_{26}^{\mathrm{nor}}=0$ も $R_{26}^{\mathrm{ev}}\in J_{26}$ も主張しない。
\end{corollary}

\begin{proof}
全次数消滅と、定理\ref{thm:v17-split-normal-projection}で定義した核の等式による。
\end{proof}

\begin{remark}[全次数の証明と有限計算の役割]
本定理は、有限個のTate係数の零性から外挿したものではない。
幾何学的なnormal coneの同定、Hesse pencilの非定数な $j$、
および多項式共変量の分裂射影を用いた証明である。
後続節の26次の二入力再計算は、この全次数命題と実装規約との独立な整合性検査になる。
また、全てのnormal係数の消滅は、実際の境界が $P+H/96$ と形式接触同値であることを
まだ意味しない。その最後の問題は、有限次元の $\mathcal K^{\mathrm{nor}}$ に残る。
\end{remark}

% END theory/v18_actual_normal_vanishing_ja.tex


% BEGIN theory/v18_effective_vertex_bound_ja.tex
% 第18版：有限頂点部分の明示次数上界。
% 一次資料：Dubé, SIAM J. Comput. 19 (1990), Theorem 8.2, p.773。

\section{頂点部分の明示的な次数上界}
\label{sec:v18-effective-vertex-bound}

第17版で得た分裂を
\[
 \mathcal Q=\bigoplus_{n\ge6}\mathcal C_n/J_n
 =\mathcal A[\Theta_8]\oplus\mathcal K^{\mathrm{nor}},
 \qquad \mathcal A=\CC[\overline U_4,\overline J_6]
\]
と書く。また
\[
 R=\CC[F_{ijk},a_{ij}],\qquad
 G=(g_0,\ldots,g_4),\qquad g_i=(\partial_iF)(q(a)),
 \qquad\mathfrak m=(a_{ij})
\]
とし、$T_0=H^0_{\mathfrak m}(R/G)$ と置く。
$F$ の独立係数は35個、$a$ の座標は10個である。
各 $g_i$ の二重次数は $(\deg_F,\deg_a)=(1,4)$ なので、
45個の変数すべてを次数1とした通常次数では5となる。

第17版では $\mathcal K^{\mathrm{nor}}\subset T_0$ と
$T_0$ の $a$ 次数の有界性を証明したが、その上界は存在的であった。
本節は、巨大ながら具体的な整数上界を与える。
係数体を $\CC(F)$ に変えると $F$ 方向のtorsionを失う可能性があるため、
以下では終始45変数の多項式環を使う。

\subsection{初項の指数から相対torsionの次数を抑える}

\begin{lemma}[初項の箱に入る相対torsion]
\label{lem:v18-initial-box-torsion}
$S=k[y_1,\ldots,y_s,x_1,\ldots,x_r]$ とし、
$J\subset S$ は $y$ 次数と $x$ 次数に関して二重斉次なイデアルとする。
単項式順序を固定し、$L=\operatorname{in}(J)$ と置く。
$L$ の最小単項式生成元の $x_i$ 指数がすべて $D_i$ 以下であるとする。
各 $D_i\ge1$ なら
\begin{equation}\label{eq:v18-initial-box-bound}
 \bigl(H^0_{(x_1,\ldots,x_r)}(S/J)\bigr)_d=0
 \qquad\left(d>\sum_{i=1}^r(D_i-1)\right),
\end{equation}
ここで添字 $d$ は $x$ 次数であり、$y$ 次数は任意である。
ある $D_i=0$ なら、この相対torsion加群全体が零である。
\end{lemma}

\begin{proof}
非零の $x$ 斉次torsion類を取り、その剰余の標準形を $f$ とする。
二重斉次イデアルの被約Gr\"obner基底は二重斉次なので、
標準形も同じ $x$ 次数 $d$ を持つ。
$m=\operatorname{lm}(f)$ と置くと $m\notin L$ である。
各 $i$ に対し、torsion性からある $e_i$ について $x_i^{e_i}f\in J$、
従って $x_i^{e_i}m\in L$ となる。

もし $m$ の $x_i$ 指数が $D_i$ 以上なら、
$x_i^{e_i}m$ を割る $L$ の最小生成元はすでに $m$ を割る。
実際、他の変数の指数は変わらず、その生成元の $x_i$ 指数も $D_i$ 以下だからである。
これは $m\notin L$ に反する。従ってすべての $i$ について
$\deg_{x_i}m\le D_i-1$ であり、式\eqref{eq:v18-initial-box-bound}を得る。
ある $D_i=0$ の場合には同じ議論がどの非零標準単項式にも矛盾を与える。
\end{proof}

\begin{remark}[実際の初項計算で使える、より精密な形]
\label{rem:v18-initial-torsion-certificate}
証明は、torsion類の先頭標準単項式 $m$ が
\[
 m\notin L,\qquad
 \text{各 }i\text{ について、ある }e_i\text{ が存在して }x_i^{e_i}m\in L
\]
を満たすことも示す。
従って完全な初項イデアルが得られれば、単なる最大指数の和より強い上界を
$H^0_{(x)}(S/L)$ の標準単項式から得られる。
ただし、計算途中のGr\"obner基底や一部の初項だけでは、
上記の $D_i$ が全初項を抑えることにはならない。
有限体で得た候補についても、特性零でのGr\"obner基底証明を別に要する。
\end{remark}

\subsection{45変数の斉次Gr\"obner次数上界の適用}

\begin{theorem}[有限頂点部分の明示上界]
\label{thm:v18-effective-vertex-cutoff}
次の整数を定める。
\begin{equation}\label{eq:v18-effective-cutoff-definition}
 D_*=2\cdot18^{\,2^{43}},\qquad
 N_*=10(D_*-1)=20\cdot18^{\,2^{43}}-10.
\end{equation}
すると
\[
 (T_0)_d=0\quad(d>N_*),\qquad
 \mathcal K_n^{\mathrm{nor}}=0\quad(n>N_*).
\]
従って第17版の有限多項式 $E_{\mathrm{vert}}$ は
\begin{equation}\label{eq:v18-Evert-explicit-degree}
 \deg E_{\mathrm{vert}}\le N_*
\end{equation}
を満たす。
\end{theorem}

\begin{proof}
Dub\'eの斉次版\cite[Theorem~8.2, p.~773]{Dube1990}
によれば、$v$ 変数の多項式環において次数 $d$ 以下の斉次多項式で生成された
イデアルは、任意の単項式順序について総次数
\[
 2\left(\frac{d^2}{2}+d\right)^{2^{v-2}}
\]
以下のGr\"obner基底を持つ。
本件は $v=45,d=5$ なので、その上界は
\[
 2(35/2)^{2^{43}}<2\cdot18^{2^{43}}=D_*.
\]
特に各初項の、個々の $a_{ij}$ の指数も $D_*$ 以下である。
補題\ref{lem:v18-initial-box-torsion}を10個の $a$ 変数に適用すると
$(T_0)_d=0$ for $d>N_*$ を得る。
第17版の包含 $\mathcal K^{\mathrm{nor}}\subset T_0$ と、
$\mathcal K_n^{\mathrm{nor}}$ の $a$ 次数が $n$ であることから残りが従う。
\end{proof}

\begin{remark}[指数と実用性の監査]
Dub\'eの同じ論文の非斉次版Corollary~8.3では指数が $2^{v-1}$ となる。
これは斉次化のために変数を一つ追加した結果であり、
本件の斉次45変数の適用では $2^{43}$ が正しい。
一方、$F$ を係数体へ移して10変数として適用することはしていない。

式\eqref{eq:v18-effective-small-conjectures}のような小さい次数の結論は、
本定理からは出ない。
\begin{equation}\label{eq:v18-effective-small-conjectures}
 E_{\mathrm{vert}}=t^{10},\qquad
 \mathcal K_{26}^{\mathrm{nor}}=0.
\end{equation}
上界 $N_*$ は実際の計算範囲としては極めて大きい。
本定理の役割は、存在的な有限性を明示的な有限判定範囲に変える点にある。
通常商の全次数計算や、26次の階数証明を実行したことを意味しない。
\end{remark}

\begin{corollary}[全通常商の残余の有限判定範囲]
\label{cor:v18-Evert-finite-decision}
$E_{\mathrm{vert}}$ の全係数は $6\le n\le N_*$ の通常商の厳密な階数で決まる。
具体的には、その各次数で
\[
 [t^n]E_{\mathrm{vert}}
 =\dim(\mathcal C_n/J_n)
 -\#\{(r,s)\in\mathbb Z_{\ge0}^2:4r+6s=n-8\}
\]
を使えばよい。
これは実行可能性に関する見積りを与えるものではなく、
有限個の厳密な線形代数問題への理論的還元である。
\end{corollary}


% END theory/v18_effective_vertex_bound_ja.tex


% BEGIN theory/v18_vertex_hessian_chain_ja.tex
% 第18版：Hessianの頂点socleと、そのq微分が作る頂点torsion。

\section{Hessianの頂点socleとその一段下の類}
\label{sec:v18-vertex-hessian-chain}

この節では、通常商の選択された同変成分
$\mathcal K^{\mathrm{nor}}$ と、全加群
$T_0=H^0_{(a)}(R/G)$ を区別する必要性を、具体的な恒等式で示す。
特に、$\mathcal K^{\mathrm{nor}}=\CC\cdot[H]$ という予想を
$T_0=\CC[F]H$ へ強めることはできない。

独立な変数 $q_0,\ldots,q_4$ を一度使い、
\[
 g_i=\partial_iF(q),\qquad M=(\partial_i\partial_jF(q))_{i,j},
 \qquad C=\operatorname{adj}(M),\qquad H=\det M,
 \qquad H_j=\partial_{q_j}H
\]
と置く。以下の $H$ は従来の $H_{10}$ と同じ非零定数倍の正規化でよい。
恒等式を記した後で $q=q(a)$ を代入する。

\begin{lemma}[Hessianとその微分の乗法恒等式]
\label{lem:v18-hessian-derivative-identity}
$B=\CC[F,q]$、$G_q=(g_0,\ldots,g_4)\subset B$ とする。
各 $j,k$ について
\begin{align}
 q_kH&=2\sum_l C_{kl}g_l,
 \label{eq:v18-qH-socle}\\
 q_kH_j-\delta_{kj}H
 &=2\sum_l (\partial_{q_j}C_{kl})g_l.
 \label{eq:v18-qHj-duality}
\end{align}
従って $B/G_q$ で $qH=0$、$q_kH_j=\delta_{kj}H$ となり、
$(q)^2H_j=0$ である。
\end{lemma}

\begin{proof}
$g_i$ が二次式なのでEulerの恒等式は $Mq=2g$ となる。
$CM=H\operatorname{id}$ を用いると式\eqref{eq:v18-qH-socle}を得る。
これを $q_j$ で微分すれば
\[
 \delta_{kj}H+q_kH_j
 =2\sum_l(\partial_{q_j}C_{kl})g_l
   +2\sum_l C_{kl}M_{lj}
 =2\sum_l(\partial_{q_j}C_{kl})g_l+2\delta_{kj}H,
\]
従って式\eqref{eq:v18-qHj-duality}が成立する。
\end{proof}

\begin{theorem}[全座標に殺されるHessian類]
\label{thm:v18-H-vertex-socle}
Pfaffian代入 $q=q(a)$ の後では、全ての $i,j$ に対して
\begin{equation}\label{eq:v18-aH-explicit}
 a_{ij}H
 =-2\sum_{k,l}a_{ik}(\partial_{q_j}C_{kl})(q(a))g_l(q(a))
 \in G.
\end{equation}
従って $R/G$ における $H$ の類は
\[
 \mathfrak m[H]=0,
 \qquad\mathfrak m=(a_{ij})
\]
を満たす。第17版以前の通常10次非零性により、この類は非零である。
\end{theorem}

\begin{proof}
符号付きPfaffianの恒等式 $aq=0$ を使い、
式\eqref{eq:v18-qHj-duality}に $a_{ik}$ を掛けて $k$ について足す。
左辺の $\sum_k a_{ik}q_kH_j$ が零となるので、
式\eqref{eq:v18-aH-explicit}を得る。
この式は $F$ の局所化や特殊化を一切使わない普遍多項式恒等式である。
\end{proof}

\begin{proposition}[全頂点加群にはHessian以外も存在する]
\label{prop:v18-Hj-vertex-torsion}
各 $j=0,\ldots,4$ について、$H_j(q(a))$ の類は
$T_0=H^0_{\mathfrak m}(R/G)$ に属し、非零である。
その二重次数は $(\deg_F,\deg_a)=(5,8)$ であり、
\begin{equation}\label{eq:v18-Hj-vertex-chain}
 q_k[H_j]=\delta_{kj}[H],\qquad
 \mathfrak m[H]=0.
\end{equation}
特に
\[
 T_0\ne\CC[F]\,[H].
\]
\end{proposition}

\begin{proof}
補題\ref{lem:v18-hessian-derivative-identity}より、
$I=(q_0,\ldots,q_4)$ に対して $I^2[H_j]=0$ である。
第17版の各Schur chart $a_{uv}\ne0$ では、
$R_{a_{uv}}=B'[p_1,p_2,p_3]$、$IR_{a_{uv}}=(p_1,p_2,p_3)$ と書ける。
このとき $q$ は $p$ の一次式で、$F$ の線形座標変換は $p$ に依存しない。
従って $H_j(q(a))$ は $p$ 次数4の斉次式であり、
その類は $(p)$-torsionである。

第17版の局所Čech--Koszul補題では、
$H^0_{(p)}(R_{a_{uv}}/G)$ の $p$ 次数4以上は零である。
従って $[H_j]$ はすべての $a_{uv}\ne0$ で消える。
有限個の座標についてのこの消失から、ある $\mathfrak m$ の冪が $[H_j]$ を殺す。
よって $[H_j]\in T_0$ である。

式\eqref{eq:v18-Hj-vertex-chain}は前補題と前定理による。
特に $q_j[H_j]=[H]\ne0$ なので、各 $[H_j]$ も非零である。
$\CC[F][H]$ は $a$ 次数10にしか成分を持たず、
$[H_j]$ は $a$ 次数8なので、これに属さない。
\end{proof}

\begin{remark}[対角の同変成分との違い]
$\mathcal K_n^{\mathrm{nor}}$ は二重次数 $(n-5,n)$ の
指定されたdeterminant指標の成分である。
上の $[H_j]$ は $(5,8)$ なので、そのまま $\mathcal K_8^{\mathrm{nor}}$ に
入るものではない。一方、$q_jH_j=H$ は $(5,10)$ に入り、
既知の $\mathcal K_{10}^{\mathrm{nor}}$ を与える。
従って全頂点加群のこの非自明な構造は、
$E_{\mathrm{vert}}=t^{10}$ の予想と両立する。
本節は他の対角成分が零であることも、その非零性も主張しない。
\end{remark}

\subsection{独立な$q$変数の完全交叉と、その適用限界}

\begin{proposition}[代入前の普遍勾配は完全交叉である]
\label{prop:v18-universal-q-CI}
$B=\CC[F_{ijk},q_0,\ldots,q_4]$ では $g_0,\ldots,g_4$ は正則列であり、
\begin{equation}\label{eq:v18-q-torsion-bound}
 \bigl(H^0_{(q)}(B/G_q)\bigr)_d=0\qquad(d>5)
\end{equation}
となる。添字は $q$ 次数を表す。
さらにこのtorsionは $F$ 次数5未満には現れない。
\end{proposition}

\begin{proof}
$q\ne0$ では $F\mapsto dF(q)$ は五次元への全射な線形写像である。
例えば $q=e_0$ とすれば、$x_0^3,x_0^2x_1,\ldots,x_0^2x_4$ の係数が
五つの独立な値を与える。
従って $q\ne0$ 上の零点集合は次元 $5+(35-5)=35$ である。
$q=0$ 上も次元35なので、40変数環 $B$ における $G_q$ の高さは5である。
$B$ はCohen--Macaulayであり、五つの生成元は正則列となる。

そのKoszul複体を $L_i=B(-i,-2i)^{\binom5i}$ と書く。
$q$ に関するČech複体との二重複体を考えると、
Koszul方向を先に取れば $H^\bullet_{(q)}(B/G_q)$ を計算する。
Čech方向を先に取ると、自由加群の局所コホモロジーは次数5だけにある。
全次数0の候補は $H^5_{(q)}(L_5)$ であり、
$H^5_{(q)}B$ の最大 $q$ 次数は $-5$ なので、
$L_5$ の $q$ シフト10を加えると最大次数は5となる。
また $L_5$ の $F$ シフトは5であり、$F$ 次数5未満は存在しない。
\end{proof}

\begin{remark}[非平坦な代入を越える際の未解決点]
$\CC[q]\to\CC[a]$、$q\mapsto q(a)$ は平坦ではない。
$q\ne0$ のファイバーは次元5である一方、零ファイバーは
七次元のPl\"ucker錐である。
従って式\eqref{eq:v18-q-torsion-bound}の $q$ 次数を単純に二倍して、
$T_0$ の $a$ 次数が10以下であると結論することはできない。
実際、代入後の五つの $g_i$ の高さは3となり、
Koszulの正次数homology $H_1,H_2$ を無視できない。
この非平坦性を、指定された同変対角成分で制御することが、
巨大な一般上界を小さい構造上界へ置き換える具体的な残課題である。
\end{remark}


% END theory/v18_vertex_hessian_chain_ja.tex


% BEGIN theory/v18_finite_determinacy_ja.tex
% 第18版：normal消滅を満たす普遍形式境界の有限決定性。
\section{normal消滅を満たす境界の明示的有限決定性}
\label{sec:v18-finite-determinacy}

全次数normal消滅と頂点核の有効上界を組み合わせると、
実際の境界の問題を、明示された有限jetの問題へ還元できる。
上界は極めて大きく、実用的な計算範囲を与えるものではない。
また本節の無限合成については、形式的な同値を証明し、収束や代数化は主張しない。

\begin{theorem}[normal消滅部分における形式的有限決定性]
\label{thm:v18-relative-finite-determinacy}
$N\ge10$ が
\[
 \mathcal K_n^{\mathrm{nor}}=0\qquad(n>N)
\]
を満たすとする。先頭項が $P$ の普遍同変形式級数
$d=P+\sum_{n>6}C_n$ が $\mathcal B_d=0$ を満たせば、
一次項が恒等の普遍同変形式座標変換 $\Phi$ によって
\begin{equation}\label{eq:v18-formal-finite-truncation}
 d(F,\Phi(F,a))=d_{\le N}(F,a)
\end{equation}
となる。右辺は変換前の $d$ の $N$ 次以下の切断である。
各形式係数は $F$ の多項式であり、単元変更を使わずに構成できる。
特に二つのそのような境界の $N$-jet が等しければ、普遍形式右同値である。
\end{theorem}

\begin{proof}
$n>N$ とし、既に $n-1$ 次まで $d_{\le N}$ に一致したとする。
それまでの座標変換後の方程式も
定理\ref{thm:v18-normal-contact-law}により $\mathcal B=0$ を満たす。
$d_{\le N}$ も同じ性質を持つので、差の $n$ 次項 $R_n$ は
\[
 \pi_n([R_n])=0
\]
を満たす。仮定 $\mathcal K_n^{\mathrm{nor}}=0$ により $R_n\in J_n$ である。
従って普遍共変場 $Y_{n-5}\in\mathcal G_{n-5}$ を
$dP(Y_{n-5})=-R_n$ となるように選べる。
$a\mapsto a+Y_{n-5}$ を代入すると、六次の先頭項 $P$ の一次微分が
ちょうど $n$ 次で $R_n$ を消し、それより低い次数は変えない。
他の係数への作用と $P$ の二階以上のTaylor項は全て $n$ 次より高い。

この操作を $n=N+1,N+2,\ldots$ に対して行う。
場の次数 $n-5$ は狭義に増加するため、合成は $a$-進位相で収束して
一次項が恒等の形式自己同型を与える。
各 $Y_{n-5}$ は有限次元の複素ベクトル空間における線形方程式を解いて選ぶため、
$F$ 係数の除算は必要ない。
有限jetに寄与する合成は有限個だけであり、各係数は $F$ の多項式のままである。
同じ $N$-jet を持つ二つの方程式をその共通の切断へ移し、
一方の形式自己同型を逆にすれば、最後の主張を得る。
\end{proof}

\begin{corollary}[実際の境界に対する一つの普遍上界]
\label{cor:v18-explicit-finite-boundary}
\begin{equation}\label{eq:v18-explicit-finite-degree}
 N_* =20\cdot18^{\,2^{43}}-10
\end{equation}
と置く。実際のTate境界は、その $N_*$ 次以下の切断に普遍形式右同値である。
また、第17版の偶正規形から出発すれば、普遍形式接触同値の下で
\begin{equation}\label{eq:v18-finite-even-normal-form}
 P+\frac1{96}H+
 \sum_{\substack{26\le n\le N_*\\n\text{ even}}}T_n,
 \qquad [T_n]\in\mathcal K_n^{\mathrm{nor}},
\end{equation}
という有限多項式にできる。各 $T_n$ は、あらかじめ選んだ
$\mathcal K_n^{\mathrm{nor}}$ の代表空間から取れる。
\end{corollary}

\begin{proof}
定理\ref{thm:v18-effective-vertex-cutoff}により
$n>N_*$ では $\mathcal K_n^{\mathrm{nor}}=0$ であり、
定理\ref{thm:v18-all-order-normal-vanishing}により実際の境界のnormal係数は全て零である。
従って前定理を適用できる。

偶正規形では、各偶数次数で通常像 $J_n$ の成分を同じ方法で除き、
核の代表だけを残す。$n$ が偶数なら場の次数 $n-5$ は奇数なので、
$a\mapsto a+Y_{n-5}$ は反転 $a\mapsto-a$ と可換であり、偶性を保つ。
核の代表を選ぶ操作は有限次元の線形代数でできる。
$n>N_*$ では代表空間は零なので、その後の項は全て消せる。
\end{proof}

\begin{remark}[有限決定性の範囲と残る問題]
この有限決定性は $\mathcal B=0$ を満たす普遍同変方程式の範囲内の主張である。
任意の形式摂動についての通常の有限決定性ではない。
自由な直和因子 $\mathcal A[\Theta_8]$ は任意の高次数にも現れ得るので、
その条件を除いた主張は本証明からは導けない。
また $N_*$ は計算済みの最小閾値ではなく、Gr\"obner基底の一般上界から得た安全な上界である。
\[
 E_{\mathrm{vert}}(t)=t^{10}
\]
が証明されれば、式\eqref{eq:v18-finite-even-normal-form}の和は全て消え、
実際の境界が $P+H/96$ と普遍形式接触同値であることが従う。
この二つは定理\ref{thm:pf-full-contact-module}の全次数証明で確定する。
\end{remark}

% END theory/v18_finite_determinacy_ja.tex


% BEGIN characters/v18_character_notes_ja.tex
\section{二十六次の普遍指標と通常商の未確定部分}
\label{sec:v18-characters}

二十六次の普遍空間を、接触写像の階数やKoszul複体の正次数完全性を
仮定せずに計算する。通常商の自由成分には二つの既知の元があるが、
頂点支持成分 $K^{\mathrm{nor}}_{26}$ の消失は本節では確定しない。

\subsection{次数、標的表現、独立指標計算}

$\Delta=\det V$、$A=\Lambda^2V$、$W=A\otimes\Delta^{-1}$ と置く。
$F\in\operatorname{Sym}^3V$、$a\in W$、
$q(a)=\tfrac12a\wedge a\in V^*\otimes\Delta^{-1}$、
$P=F(q(a))\in\Delta^{-3}$ の規約を維持する。
二十六次のKoszul第 $k$ 項は
\[
 K_{k,26}=\operatorname{Hom}_{\mathrm{GL}(V)}
 \left(\operatorname{Sym}^{21-k}(\operatorname{Sym}^3V)
 \otimes\operatorname{Sym}^{26-4k}A,
 \Lambda^kV^*\otimes\Delta^{23-2k}\right),\qquad 0\leq k\leq5.
\]
$K_{0,26}=C_{26}$、$K_{1,26}=D_{22}$、$K_{2,26}=B_{26}$ と記す。
終端 $K_{5,26}$ にも $a$ 次数6が残る。

\begin{proposition}[二十六次の厳密次元]
\label{prop:v18-character-dimensions}
普遍空間の次元は次の通りである。
\[
\begin{array}{c|r|r|c|r}
\text{空間}&\deg_F&\deg_a&\text{標的分割}&\dim\\\hline
C_{26}&21&26&(23,23,23,23,23)&6377\\
G_{21}&20&21&(21,21,20,20,20)&32805\\
U_{20}&20&20&(20,20,20,20,20)&3712\\
D_{22}&20&22&(21,21,21,21,20)&17393\\
B_{26}&19&18&(19,19,19,18,18)&17043\\
K_{3,26}&18&14&(17,17,16,16,16)&7121\\
K_{4,26}&17&10&(15,14,14,14,14)&1151\\
K_{5,26}&16&6&(12,12,12,12,12)&52
\end{array}
\]
補助空間では $\dim C_{17}=31$、$\dim G_{12}=155$、
$\dim U_{11}=9$、$\dim D_{13}=92$ を再確認する。
\end{proposition}

\begin{proof}
三次単項式の35個の重み $\alpha$ からEuler積
$\prod_{|\alpha|=3}(1-tx^\alpha)^{-1}$ を作る。
その $t^m$ 係数に、$\rho=(4,3,2,1,0)$ を用いたWeyl交代和
\[
 \sum_{\sigma\in S_5}\operatorname{sgn}(\sigma)
 [x^{\lambda+\rho-\sigma\rho}]
 \operatorname{ch}\operatorname{Sym}^m(\operatorname{Sym}^3V)
\]
を適用してSchur多重度を得る。
外積二乗側も10個の重みから独立に指標を生成し、
テンソル積の交代和を直接評価する。
Littlewood分解
\[
 \operatorname{Sym}^d(\Lambda^2V)
 =\bigoplus_{j=0}^{\lfloor d/2\rfloor}S_{(d-j,d-j,j,j,0)}V
\]
の各多重度1と、Weyl次元和も照合した。

$C_{26}$ に現れる $F$ 側の分割は
$(23,20-j,20-j,j,j)$、$0\leq j\leq10$ である。
その多重度列は
\[
 (4,6,50,137,532,942,1699,1502,1170,280,55)
\]
となり、和は6377である。
$U_{20}$ の分割 $(20,20-j,20-j,j,j)$ に対する多重度列は
\[
 (2,2,19,54,255,493,983,923,756,187,38)
\]
となり、和は3712である。
$G_{21}$ と $D_{22}$ はそれぞれ外積第3冪・第1冪の垂直帯Pieri計算で、
Koszul各項は外積第 $k$ 冪のPieri計算で、直接のテンソル積指標と一致する。
Koszul源の非零Pieri成分数は $k=2,3,4,5$ の順に33、28、16、4である。

全ての $1\leq m\leq21$ の $F$ 指標と $1\leq d\leq26$ の $a$ 指標を
Newton--Adams再帰 $m h_m=\sum_{k=1}^m p_kh_{m-k}$ により別途生成し、
Euler積の全重み係数辞書と一致することを確認した。
必要な $F$ 次数 $11,12,16,17,18,19,20,21$ の全Schur分解は、
各多重度とWeyl次元の積の和が $\binom{m+34}{m}$ に一致する。
このうち前版にもある次数は、保存JSONの表現形式を正規化してから
全Schur分解を照合した。今回の全重み再計算と、旧証明書との保存値照合は
異なる検証として記録している。
\end{proof}

\subsection{Euler交代和が決める範囲}

今回の交代和は
\[
 6377-17393+17043-7121+1151-52=5.
\]
既証明の $H_i=0$（$i>2$）を適用すると
\[
 \dim H_0-\dim H_1+\dim H_2=5,
\]
および
\[
 \operatorname{rank}\delta_5=52,\qquad
 \operatorname{rank}\delta_4=1151-52=1099,\qquad
 \operatorname{rank}\delta_3=7121-1099=6022
\]
を得る。Pfaffian微分の全射性からは
\[
 \dim\ker(\delta q:G_{21}\to D_{22})=32805-17393=15412
\]
が従う。これらはいずれも $\operatorname{rank}(dP)$ の計算ではない。

$\kappa_{26}=\dim K^{\mathrm{nor}}_{26}$、$b_{26}=\operatorname{rank}\delta_2$ と置く。
既存の分裂
$Q=A_0[\Theta_8]\oplus K^{\mathrm{nor}}$、
$A_0=\mathbf C[\bar U_4,\bar J_6]$ を用いると
\[
 \dim Q_{26}=2+\kappa_{26},\qquad
 \dim J_{26}=6375-\kappa_{26}.
\]
従って
\[
 \dim H_1=11018+\kappa_{26}-b_{26},\qquad
 \dim H_2=11021-b_{26},\qquad
 \dim H_2-\dim H_1=3-\kappa_{26}.
\]
特に、交代和5から $\kappa_{26}=0$ や $H_1=0$ を推論することはできない。
$\kappa_{26}=0$ を別に証明できた場合に限り、
$\dim H_2=\dim H_1+3\geq3$ が従う。

\subsection{単元像、自由成分、必要な通常階数}

多項式環 $\mathbf C[F,a]$ における非零多項式 $P$ の乗法は単射であるから、
$PU_{20}\subset J_{26}$ は3712次元である。
この所属はEuler恒等式 $dP(a)=6P$ からも分かる。
通常商の自由成分の二つの元は
\[
 [\Theta_8U_4^3J_6],\qquad [\Theta_8J_6^3].
\]
これらの独立性は既存のnormal射影により確定している。
全通常商をこの二元が張ることは $\kappa_{26}=0$ と同値であり、
本版の有限rank計算では証明していない。

$P=0$ 上の評価で通常像の2663本の独立列を構成できれば、
\[
 \dim J_{26}\geq3712+2663=6375
\]
と理論上界 $\dim J_{26}\leq6375$ が一致する。
この条件を満たす全minorの新規構成・独立監査は今回は実施していない。
96点でのpilotの非零minorは、その有限の下界だけを与える。
小標本の階数飽和や候補列の枯渇を通常商の上界には用いない。

補助的に、$U_s$ 内の $P$ 可除部分を矩形補分割で計算した。
$P$ の係数空間は
\[
 \operatorname{Hom}_{\mathrm{GL}(V)}
 \left(\operatorname{Sym}^{s-1}(\operatorname{Sym}^3V)
 \otimes\operatorname{Sym}^{s-6}A,\Delta^{s-3}\right)
\]
であり、$F$ 側の分割は
$(s-3,s-3-j,s-3-j,3+j,3+j)$ となる。
保存された完全Schur分解から得る値は
\[
\begin{array}{c|r|r|r}
 &\dim U_s&\dim(P\text{ 可除部分})&\dim(U_s/(P\text{ 可除部分}))\\\hline
U_{13}&47&5&42\\
U_{16}&486&141&345\\
U_{20}&3712&1406&2306
\end{array}
\]
である。この表は単元の源空間の計算であり、候補グラフの基底性を主張しない。

\subsection{制限の標的次元と証明書の範囲}

Pl\"ucker錐へ制限したHom空間の次元は
\[
\begin{array}{c|rrrrrrrr}
 &C_{26}&G_{21}&U_{20}&D_{22}&B_{26}&K_{3,26}&K_{4,26}&K_{5,26}\\\hline
\dim&0&2&2&0&0&22&66&21
\end{array}
\]
である。これはLittlewood分解のCartan成分のみを残したPieri計算であり、
任意の普遍制限写像の全射性を仮定していない。

Fermat固定群の成分別計算による評価階数の上界は
\[
\begin{array}{c|r|r|r}
 &\dim&\text{対角値の階数上界}&\text{対角値と一階法jetの階数上界}\\\hline
C_{26}&6377&1332&4478\\
G_{21}&32805&4854&19814\\
U_{20}&3712&487&1991\\
D_{22}&17393&2834&11092
\end{array}
\]
となる。実際の評価階数との等号は主張しない。

\nolinkurl{v18_character_certificate.py} は全重みの新規整数計算、
\nolinkurl{v18_character_audit.py} は保存データの監査を担う。

\nolinkurl{v18_koszul_pieri_check.py} は別実装による垂直帯列挙を担う。
\nolinkurl{v18_unit_divisibility.py} の補助次数13・15の表には旧版の出典を残す。
通常商のrank pilotと選択された小minorの独立再縮約は
\texttt{rank/} 以下に分けて記録する。

% END characters/v18_character_notes_ja.tex


% BEGIN detector/v18_normal_detector_ja.tex
% 第18版：26次 normal 検出子と独立保存監査。
\section{二十六次normal検出子と厳密な零判定}
\label{sec:v18-normal26-detectors}

normal射影の26次の標的は
\[
 \mathcal A_{18}[\Theta_8]
 =\mathbb C\cdot[\Theta_8U_4^3J_6]\oplus\mathbb C\cdot[\Theta_8J_6^3]
\]
である。以下の二つの検出子は、この二次元成分を完全に読む。
通常商全体の次元は、この標的の次元だけからは決まらない。

$m=-1,-2$ に対して
\[
 F_m=x_0^3+x_1^3+x_2^3+x_3^3+(-2-6m)x_4^3+6mx_2x_3x_4,
 \qquad a_\varepsilon=e_{01}+\varepsilon(e_2-e_4)\wedge(e_3-e_4)
\]
とし、$L_m(C)=[\varepsilon^3]C(F_m,a_\varepsilon)$ と定める。
どちらも三元部分 $f_m$ と $p=(1,1,1)$ について $f_m(p)=0$ を満たす。
三元式の滑らかさの判別条件 $(-2-6m)+8m^3$ はそれぞれ $-4,-54$ で非零である。
第17版のnormal射影定理により、$L_m$ は $\mathcal J_{26}$ を消す。

\begin{proposition}[二つの有理normal検出子]\label{prop:v18-normal-jet-detectors}
基底を $W_1=\Theta_8U_4^3J_6$、$W_2=\Theta_8J_6^3$ とすると
\[
 \begin{array}{c|rrr}
 m&\overline U_4&\overline J_6&h_{f_m}(p)\\\hline
 -1&5&528&-4\\
 -2&36&7128&-54
 \end{array}
\]
であり、評価行列とその行列式は
\[
 M_{26}=
 \begin{pmatrix}
 -264000&-588791808\\
 -17958454272&-19556756702208
 \end{pmatrix},\qquad
 \det M_{26}=-5410806990313291776.
\]
従って
\begin{align*}
 \rho_{26,1}(C)&=
 \frac{531441L_{-1}(C)-16L_{-2}(C)}{147034844352},\\
 \rho_{26,2}(C)&=
 \frac{-8503056L_{-1}(C)+125L_{-2}(C)}{2561935127989248}
\end{align*}
は $\rho_{26,i}(W_j)=\delta_{ij}$ を満たし、normal射影の二係数を与える。
\end{proposition}
\begin{proof}
\texttt{v18\_normal\_detector\_certificate.py} は旧版のtensor graphの定義から
rank二切片上の $U_4,J_6$ を再評価し、三元Hessianを有理数演算で直接求める。
$U_4$ のgraph値は規約通り24で割る。
各行の値は $h_f(p)(\overline U_4^3\overline J_6,\overline J_6^3)$ であり、
逆行列を厳密に求めると上の二つの有理式になる。
出力には全入力テンソル、rank二点、法方向、旧入力とgraph定義のSHA-256を保存した。
\end{proof}

\begin{proposition}[二十六次の新規normal Tate計算と保存監査]
\label{prop:v18-normal26-tate-zero}
二入力で新規計算したraw Tate境界のnormal三次係数は
\[
 L_{-1}(d_{26}^{\mathrm T})=L_{-2}(d_{26}^{\mathrm T})=0.
\]
従って
\[
 \pi_{26}([d_{26}^{\mathrm T}])=0,\qquad
 \pi_{26}([R_{26}^{\mathrm{ev}}])=0.
\]
\end{proposition}
\begin{proof}
\texttt{v18\_tate\_normal\_jet.py} は
$\mathbb Q[\varepsilon]/(\varepsilon^4)$ 上でMaurer--Cartan再帰を最初から行う。
全 $A_1,\ldots,A_{18}$ と、曲率 $R_{19},\ldots,R_{22}$ の$x$次数
$4,3,2,1$ 以下を新規計算し、22個の定数ブロックを得た。
行列式の規約は $d^{\mathrm T}=\tfrac12\det K$ である。
二入力それぞれで、$a$次数0から26、$\varepsilon$次数0から3の
$27\times4=108$ 個の行列式係数は全て零となった。

別の \texttt{v18\_normal\_saved\_audit.py} は研究計算モジュールを一切importせず、
標準ライブラリの \texttt{Fraction} によって、保存ブロックから
置換展開と稠密な部分集合動的計画の二方法で全216係数を再構成した。
各入力の125成分の対称テンソル、$a_0$、法方向、$K_1=-a$ と
$K_2=2F(q,-,-)$ の等式、次数切断のmetadata、計算コードのhashも照合した。
最後に前命題の逆行列を適用して $(\rho_{26,1},\rho_{26,2})=(0,0)$ を確認した。
これは保存監査であり、新規再帰そのものは前記Tateスクリプトによる。
実際の26次係数との等式は系\ref{cor:v18-actual26-normal-equals-raw}による。
\end{proof}

\begin{remark}[結論の範囲]
ここでの零判定は $[R_{26}^{\mathrm{ev}}]\in\mathcal K_{26}^{\mathrm{nor}}$ を意味する。
通常像への所属には $\mathcal K_{26}^{\mathrm{nor}}=0$、通常26次の必要なrank、
またはこの一つの類の直接の消去を別に示す必要がある。
新規normal再帰と独立保存監査のいずれも、その追加主張を行わない。
\end{remark}

% END detector/v18_normal_detector_ja.tex

\clearpage
\section{今後の課題}

第18版では、実際の境界の横断Hessian係数が全次数で零であることを証明した。
従って全次数Hessian正規形の問題は、有限次元の頂点部分
$\mathcal K^{\mathrm{nor}}$ に集中する。
明示的上界 $N_*=20\cdot18^{2^{43}}-10$ は得たが、実用的な閾値ではない。
以下では、この理論的な有限還元を、小さい構造定理または有限証明書へ置き換える。
\begin{enumerate}
 \item 指定された同変頂点成分 $\mathcal K^{\mathrm{nor}}$ を決定する。
 定理\ref{thm:pf-full-contact-module}で $E_{\mathrm{vert}}=t^{10}$ を証明した。
 全次数normal消滅と合わせた $P+H/96$ 正規形とその代数化も
 第\ref{sec:pf-equivariant-algebraization}節で与える。
 全頂点加群 $T_0$ には $H_j=\partial_{q_j}H$ も存在するので、
 $T_0=\mathbb C[F]H$ という強すぎる簡約は使えない。
 独立な $q$ 変数での完全交叉と、非平坦なPfaffian代入の差を
 二重次数とdeterminant指標を保ったまま制御する。

 \item 二十六次の残る頂点類を決める。
 通常26次商の二次元性は定理\ref{thm:pf-full-contact-module}から従う。
 $[R_{26}^{\mathrm{ev}}]\in\mathcal K_{26}^{\mathrm{nor}}$ の直接消去、
 小さい理論的上界、または十分な通常接触minorのいずれかが必要になる。
 96列の有限minorは部分的な下界証拠であり、必要な2663本の独立方向を与えていない。
 24次の核の消失を26次へ外挿せず、Euler交代和を完全性として扱わない。

 \item 全次数normal消滅の幾何学的証明をTate模型内でも説明する。
 split三次式に沿う定数ブロックの構造、Picard形式対数、Hesse pencilの一定な
 抽象曲線との対応を直接表す式を求める。
 今回の証明は全次数を扱うため、normal成分の零判定だけを目的とした
 次数ごとの大型計算は必要ない。
 ただし頂点成分を読むためには、法三次jetを超える計算が引き続き必要である。

 \item 普遍Koszulの $H_1,H_2$ を決める。
 高次 $H_i$ ($i>2$) の消失は既知であるが、二つの個別次元は未確定である。
 26次のEuler和は5であり、通常商の次元を単独で与えない。
 五つの勾配のheightは3なので、代入後も五元正則列であるとは扱わない。

 \item 有限形式正規形の収束・代数化と、表示の縮約を調べる。
 特定の $N_*$-jet への無限合成自体の収束は別問題だが、同じ十次正規形への代数的・収束的な同値の存在を第\ref{sec:pf-equivariant-algebraization}節で示す。
 既知の有限次数の代数的・収束的正規化と、そのまま同一視しない。
 また4929補助変数と共役の5890追加未知数を含む既存表示を小さくし、
 $16\times5$ 行列によるincidenceと相対Picardの幾何を直接比較する課題が残る。

 \item 被約特異部分は命題\ref{thm:global-bn-singular}で $C$ と同定したが、
 Jacobian idealが定める非被約な特異スキームの構造を決定する。
 一段階blowup内でincidenceを定める $p^*\OO(3)$ の切断の束方向6次以上の項を求め、
 二段階解消の記述と比較する。
 相対Picardによるblowupの形式記述を大域的な代数モジュライへ延長すること、
 および $G$ の相対対象モジュライの分離性・propernessも別の課題である。
 \item 明示的族の線形束 $L_R$ のPicard座標を、Sklyaninの平行移動点 $p(h)$ と
 直接結ぶ公式を求める。
 現在は固定形式曲線を介した対応と、その一次係数が非零であることが分かっている。
 Abel--Jacobi座標での公式は、gluing対象と非被約点の局所環との関係を精密にする。
 \item 構造対象や次数シフトを指定しない、全体dg圏の同値類を区別する。
 第\ref{sec:five-generator}節の分類は次数付き代数と偏極付きtailsについてのものである。
 抽象point schemeからは四つの次数付き同型類が残る。
 六次境界は元の $X$ を復元するが、その境界を残余圏だけから内在的に指定できるかは
 別に検討する必要がある。gluing対象から非被約点の局所環への普遍的な変形写像も未構成である。
 \item 一般の三次式に対し、二つの中心パラメータを持つ有限表示の
 代数的量子化を構成し、残余圏へ写る二次項を明示する。
 point moduleの計算をline module・理想層へ拡張し、Fano surfaceとの関係を調べる。
 評価核から曲線と線形束を $F_L^{**}/F_L=j_*L$ として回復する構成を、
 他の数値類や非自由支持型へ広げることも考えられる。
\end{enumerate}
